% Recovered from the previously completed PDF; no new mathematical content.
% Prose is editable; formula boxes retain original symbol positions.
% Build from repository root: xelatex -no-shell-escape -halt-on-error -output-directory=... print/fudan/NAME.tex
\documentclass[11pt,a4paper]{article}
\usepackage[margin=23mm,headheight=14pt,headsep=8.5mm,footskip=10mm]{geometry}
\usepackage{amsmath,amssymb}
\DeclareMathAlphabet{\mathsf}{OT1}{cmss}{m}{n}
\usepackage[no-math]{fontspec}
\usepackage{xeCJK}
\setmainfont{Latin Modern Roman}
\setCJKmainfont{FandolSong-Regular.otf}[Path=assets/fonts/fandol/,BoldFont=FandolSong-Bold.otf,ItalicFont=FandolSong-Regular.otf,Script=CustomDefault]
\xeCJKsetup{PunctStyle=quanjiao,CheckSingle=true}
\usepackage{setspace}\setstretch{1.12}
\setlength{\parindent}{2em}\setlength{\parskip}{3pt}
\emergencystretch=6em
\hfuzz=9pt
\usepackage{titlesec}
\titleformat{\section}{\centering\normalfont\Large\bfseries}{}{0pt}{}
\titlespacing*{\section}{0pt}{1.8\baselineskip}{1.0\baselineskip}
\usepackage{fancyhdr}\pagestyle{fancy}\fancyhf{}
\fancyhead[L]{\small 数学手记·高等代数}\fancyhead[R]{\small 21A10–21A15、22S01–22S01}
\fancyfoot[C]{\thepage}\renewcommand{\headrulewidth}{0.4pt}
\usepackage[hidelinks,unicode]{hyperref}
\hypersetup{pdftitle={21A10–21A15、22S01–22S01：题目、解答与小结},pdfauthor={Vincent Arqain}}
\newcommand{\RB}[4]{\raisebox{0pt}[#2bp][#3bp]{\hbox to #1bp{#4\hss}}}
\newcommand{\RG}[5]{\hbox to 0pt{\kern#1bp\raisebox{#2bp}[0pt][0pt]{\fontsize{#3bp}{#3bp}\csname font#4\endcsname\char#5}\hss}}
\newcommand{\RT}[4]{\hbox to 0pt{\kern#1bp\raisebox{#2bp}[0pt][0pt]{\fontsize{#3bp}{#3bp}\selectfont#4}\hss}}
\newcommand{\RS}[4]{\hbox to 0pt{\kern#1bp\raisebox{#2bp}[0pt][0pt]{\fontsize{#3bp}{#3bp}\selectfont\ensuremath{#4}}\hss}}
\newcommand{\RR}[4]{\hbox to 0pt{\kern#1bp\raisebox{#2bp}[0pt][0pt]{\rule{#3bp}{#4bp}}\hss}}
\newcommand{\fontmi}{\usefont{OML}{cmm}{m}{it}}
\newcommand{\fontsy}{\usefont{OMS}{cmsy}{m}{n}}
\newcommand{\fontex}{\usefont{OMX}{cmex}{m}{n}}
\newcommand{\fontrm}{\usefont{OT1}{cmr}{m}{n}}
\newcommand{\fontbf}{\usefont{OT1}{cmr}{bx}{n}}
\newcommand{\fontam}{\usefont{U}{msa}{m}{n}}
\newcommand{\fontbb}{\usefont{U}{msb}{m}{n}}
\widowpenalty=10000\clubpenalty=10000\raggedbottom
\begin{document}
\begin{center}{\LARGE\bfseries 21A10–21A15、22S01–22S01\par}\vspace{8pt}复旦大学高等代数习题集·题目与解答\end{center}
% Original page 1.
\section*{2021A10}
设\allowbreak{}\hskip.1em \RB{36.077}{8.686}{3.223}{%
\RG{0.000}{0.004}{10.9091}{rm}{40}\RG{4.242}{0.004}{10.9091}{mi}{86}\RG{11.202}{0.004}{10.9091}{mi}{59}\RG{16.060}{0.004}{10.9091}{rm}{43}\RG{24.545}{0.004}{10.9091}{mi}{59}\RG{29.394}{0.004}{10.9091}{sy}{1}\RG{32.424}{0.004}{10.9091}{rm}{41}%
}\hskip.1em\allowbreak{} 是数域\allowbreak{}\hskip.1em \RB{10.198}{7.955}{0.500}{%
\RG{0.000}{0.004}{10.9091}{mi}{75}%
}\hskip.1em\allowbreak{} 上的线性空间，\allowbreak{}\hskip.1em \RB{49.405}{8.195}{2.874}{%
\RG{0.000}{0.004}{10.9091}{mi}{39}\RG{10.166}{0.004}{10.9091}{rm}{58}\RG{16.227}{0.004}{10.9091}{mi}{83}\RG{26.575}{0.004}{10.9091}{sy}{33}\RG{40.515}{0.004}{10.9091}{mi}{86}%
}\hskip.1em\allowbreak{} 是双射。在集合\allowbreak{}\hskip.1em \RB{7.536}{8.195}{0.736}{%
\RG{0.000}{0.004}{10.9091}{mi}{83}%
}\hskip.1em\allowbreak{} 上定义

\[
\RB{243.573}{10.305}{4.306}{%
\RG{56.782}{8.838}{10.9092}{ex}{0}\RG{116.807}{8.838}{10.9092}{ex}{1}\RG{203.567}{8.838}{10.9092}{ex}{0}\RG{235.978}{8.838}{10.9092}{ex}{1}\RG{45.463}{4.505}{7.9701}{sy}{0}\RG{52.050}{4.505}{7.9701}{rm}{49}\RG{192.248}{4.505}{7.9701}{sy}{0}\RG{198.835}{4.505}{7.9701}{rm}{49}\RG{0.000}{0.001}{10.9091}{mi}{97}\RG{8.190}{0.001}{10.9091}{sy}{8}\RG{19.099}{0.001}{10.9091}{mi}{98}\RG{26.812}{0.001}{10.9091}{rm}{61}\RG{38.327}{0.001}{10.9091}{mi}{39}\RG{61.783}{0.001}{10.9091}{mi}{39}\RG{68.919}{0.001}{10.9091}{rm}{40}\RG{73.162}{0.001}{10.9091}{mi}{97}\RG{78.928}{0.001}{10.9091}{rm}{41}\RG{85.582}{0.001}{10.9091}{rm}{43}\RG{96.504}{0.001}{10.9091}{mi}{39}\RG{103.640}{0.001}{10.9091}{rm}{40}\RG{107.883}{0.001}{10.9091}{mi}{98}\RG{112.564}{0.001}{10.9091}{rm}{41}\RG{121.807}{0.001}{10.9091}{mi}{59}\RG{148.458}{0.001}{10.9091}{mi}{107}\RG{156.920}{0.001}{10.9091}{sy}{14}\RG{164.799}{0.001}{10.9091}{mi}{97}\RG{173.596}{0.001}{10.9091}{rm}{61}\RG{185.112}{0.001}{10.9091}{mi}{39}\RG{208.567}{0.001}{10.9091}{mi}{107}\RG{214.578}{0.001}{10.9091}{mi}{39}\RG{221.726}{0.001}{10.9091}{rm}{40}\RG{225.969}{0.001}{10.9091}{mi}{97}\RG{231.735}{0.001}{10.9091}{rm}{41}\RG{240.978}{0.001}{10.9091}{mi}{58}%
}
\]
(1)证明\allowbreak{}\hskip.1em \RB{38.244}{8.681}{3.228}{%
\RG{0.000}{-0.001}{10.9091}{rm}{40}\RG{4.243}{-0.001}{10.9091}{mi}{83}\RG{10.952}{-0.001}{10.9091}{mi}{59}\RG{15.803}{-0.001}{10.9091}{sy}{8}\RG{24.288}{-0.001}{10.9091}{mi}{59}\RG{29.137}{-0.001}{10.9091}{sy}{14}\RG{34.591}{-0.001}{10.9091}{rm}{41}%
}\hskip.1em\allowbreak{} 为\allowbreak{}\hskip.1em \RB{10.198}{7.950}{0.501}{%
\RG{0.000}{-0.001}{10.9091}{mi}{75}%
}\hskip.1em\allowbreak{} 上的线性空间，并且\allowbreak{}\hskip.1em \RB{7.253}{5.321}{2.879}{%
\RG{0.000}{-0.001}{10.9091}{mi}{39}%
}\hskip.1em\allowbreak{} 为线性同构。

(2)设\allowbreak{}\hskip.1em \RB{7.504}{8.196}{0.735}{%
\RG{0.000}{0.005}{10.9091}{sy}{83}%
}\hskip.1em\allowbreak{} 是\allowbreak{}\hskip.1em \RB{6.729}{5.327}{0.615}{%
\RG{0.000}{0.005}{10.9091}{mi}{110}%
}\hskip.1em\allowbreak{} 阶实反对称阵全体，\allowbreak{}\hskip.1em \RB{9.598}{8.327}{1.237}{%
\RG{0.000}{0.005}{10.9091}{sy}{84}%
}\hskip.1em\allowbreak{} 是满足\allowbreak{}\hskip.1em \RB{32.018}{7.956}{2.211}{%
\RG{0.000}{0.005}{10.9091}{mi}{73}\RG{4.795}{-1.631}{7.9701}{mi}{110}\RG{12.619}{0.005}{10.9091}{rm}{43}\RG{23.292}{0.005}{10.9091}{mi}{80}%
}\hskip.1em\allowbreak{} 可逆的\allowbreak{}\hskip.1em \RB{6.729}{5.327}{0.615}{%
\RG{0.000}{0.005}{10.9091}{mi}{110}%
}\hskip.1em\allowbreak{} 阶正交阵全体。已知Cayley变换

\[
\RB{185.535}{10.302}{3.229}{%
\RG{125.092}{4.502}{7.9701}{sy}{0}\RG{131.679}{4.502}{7.9701}{rm}{49}\RG{0.000}{-0.002}{10.9091}{rm}{8}\RG{7.876}{-0.002}{10.9091}{rm}{40}\RG{12.122}{-0.002}{10.9091}{mi}{88}\RG{22.016}{-0.002}{10.9091}{rm}{41}\RG{29.271}{-0.002}{10.9091}{rm}{61}\RG{40.780}{-0.002}{10.9091}{rm}{40}\RG{45.046}{-0.002}{10.9091}{mi}{73}\RG{57.902}{-0.002}{10.9091}{sy}{0}\RG{68.812}{-0.002}{10.9091}{mi}{88}\RG{78.706}{-0.002}{10.9091}{rm}{41}\RG{82.938}{-0.002}{10.9091}{rm}{40}\RG{87.191}{-0.002}{10.9091}{mi}{73}\RG{100.047}{-0.002}{10.9091}{rm}{43}\RG{110.956}{-0.002}{10.9091}{mi}{88}\RG{120.850}{-0.002}{10.9091}{rm}{41}\RG{136.411}{-0.002}{10.9091}{mi}{59}\RG{152.153}{-0.002}{10.9091}{mi}{88}\RG{165.093}{-0.002}{10.9091}{sy}{50}\RG{175.381}{-0.002}{10.9091}{sy}{83}\RG{182.820}{-0.002}{10.9091}{mi}{59}\RG{49.841}{-1.638}{7.9701}{mi}{110}\RG{91.986}{-1.638}{7.9701}{mi}{110}%
}
\]
是从\allowbreak{}\hskip.1em \RB{7.504}{8.190}{0.741}{%
\RG{0.000}{-0.001}{10.9091}{sy}{83}%
}\hskip.1em\allowbreak{} 到\allowbreak{}\hskip.1em \RB{9.598}{8.321}{1.243}{%
\RG{0.000}{-0.001}{10.9091}{sy}{84}%
}\hskip.1em\allowbreak{} 的双射，其逆映射为\allowbreak{}\hskip.1em \RB{131.580}{9.757}{3.228}{%
\RG{0.000}{-0.001}{10.9091}{rm}{9}\RG{8.476}{-0.001}{10.9091}{rm}{40}\RG{12.728}{-0.001}{10.9091}{mi}{80}\RG{21.247}{-0.001}{10.9091}{rm}{41}\RG{28.502}{-0.001}{10.9091}{rm}{61}\RG{40.011}{-0.001}{10.9091}{rm}{40}\RG{44.277}{-0.001}{10.9091}{mi}{73}\RG{49.072}{-1.638}{7.9701}{mi}{110}\RG{56.972}{-0.001}{10.9091}{sy}{0}\RG{67.722}{-0.001}{10.9091}{mi}{80}\RG{76.241}{-0.001}{10.9091}{rm}{41}\RG{80.473}{-0.001}{10.9091}{rm}{40}\RG{84.726}{-0.001}{10.9091}{mi}{73}\RG{89.521}{-1.638}{7.9701}{mi}{110}\RG{97.421}{-0.001}{10.9091}{rm}{43}\RG{108.169}{-0.001}{10.9091}{mi}{80}\RG{116.688}{-0.001}{10.9091}{rm}{41}\RG{120.931}{3.957}{7.9701}{sy}{0}\RG{127.517}{3.957}{7.9701}{rm}{49}%
}\hskip.1em\allowbreak{} 。给出\allowbreak{}\hskip.1em \RB{9.598}{8.321}{1.243}{%
\RG{0.000}{-0.001}{10.9091}{sy}{84}%
}\hskip.1em\allowbreak{} 上加法和数乘的具体表达式，使\allowbreak{}\hskip.1em \RB{8.355}{7.956}{0.500}{%
\RG{0.000}{0.005}{10.9091}{rm}{9}%
}\hskip.1em\allowbreak{} 成为实线性同构。解：(1)定义\allowbreak{}\hskip.1em \RB{65.210}{9.762}{3.224}{%
\RG{0.000}{0.003}{10.9091}{rm}{48}\RG{5.454}{-1.686}{7.9701}{mi}{83}\RG{14.578}{0.003}{10.9091}{rm}{61}\RG{26.094}{0.003}{10.9091}{mi}{39}\RG{33.230}{3.962}{7.9701}{sy}{0}\RG{39.816}{3.962}{7.9701}{rm}{49}\RG{44.549}{0.003}{10.9091}{rm}{40}\RG{48.782}{0.003}{10.9091}{rm}{48}\RG{54.246}{-1.686}{7.9701}{mi}{86}\RG{61.558}{0.003}{10.9091}{rm}{41}%
}\hskip.1em\allowbreak{} 及\allowbreak{}\hskip.1em \RB{85.011}{9.762}{3.224}{%
\RG{0.000}{0.003}{10.9091}{sy}{9}\RG{8.485}{0.003}{10.9091}{mi}{97}\RG{17.282}{0.003}{10.9091}{rm}{61}\RG{28.797}{0.003}{10.9091}{mi}{39}\RG{35.933}{3.962}{7.9701}{sy}{0}\RG{42.520}{3.962}{7.9701}{rm}{49}\RG{47.252}{0.003}{10.9091}{rm}{40}\RG{51.495}{0.003}{10.9091}{sy}{0}\RG{59.980}{0.003}{10.9091}{mi}{39}\RG{67.117}{0.003}{10.9091}{rm}{40}\RG{71.359}{0.003}{10.9091}{mi}{97}\RG{77.126}{0.003}{10.9091}{rm}{41}\RG{81.358}{0.003}{10.9091}{rm}{41}%
}\hskip.1em\allowbreak{} ，则

\[
\RB{148.152}{8.682}{3.227}{%
\RG{0.000}{-0.000}{10.9091}{mi}{97}\RG{8.190}{-0.000}{10.9091}{sy}{8}\RG{19.099}{-0.000}{10.9091}{rm}{48}\RG{33.678}{-0.000}{10.9091}{rm}{61}\RG{45.193}{-0.000}{10.9091}{mi}{97}\RG{50.953}{-0.000}{10.9091}{mi}{59}\RG{77.605}{-0.000}{10.9091}{mi}{97}\RG{85.818}{-0.000}{10.9091}{sy}{8}\RG{96.727}{-0.000}{10.9091}{rm}{40}\RG{100.969}{-0.000}{10.9091}{sy}{9}\RG{109.454}{-0.000}{10.9091}{mi}{97}\RG{115.221}{-0.000}{10.9091}{rm}{41}\RG{122.476}{-0.000}{10.9091}{rm}{61}\RG{133.985}{-0.000}{10.9091}{rm}{48}\RG{145.557}{-0.000}{10.9091}{mi}{58}\RG{24.554}{-1.688}{7.9701}{mi}{83}\RG{139.463}{-1.688}{7.9701}{mi}{83}%
}
\]
双射的单射性使得\allowbreak{}\hskip.1em \RB{8.889}{7.949}{0.742}{%
\RG{0.000}{-0.002}{10.9091}{mi}{86}%
}\hskip.1em\allowbreak{} 上的运算规律可逐条转移到\allowbreak{}\hskip.1em \RB{7.536}{8.189}{0.742}{%
\RG{0.000}{-0.002}{10.9091}{mi}{83}%
}\hskip.1em\allowbreak{} 。例如，

\[
\RB{265.422}{67.109}{4.309}{%
\RG{7.137}{66.173}{10.9092}{ex}{0}\RG{62.459}{66.173}{10.9092}{ex}{1}\RG{82.004}{66.173}{10.9092}{ex}{0}\RG{142.028}{66.173}{10.9092}{ex}{1}\RG{0.000}{57.335}{10.9091}{mi}{39}\RG{12.137}{57.335}{10.9091}{rm}{40}\RG{16.379}{57.335}{10.9091}{mi}{97}\RG{24.570}{57.335}{10.9091}{sy}{8}\RG{35.479}{57.335}{10.9091}{mi}{98}\RG{40.161}{57.335}{10.9091}{rm}{41}\RG{46.827}{57.335}{10.9091}{sy}{8}\RG{57.737}{57.335}{10.9091}{mi}{99}\RG{70.489}{57.335}{10.9091}{rm}{61}\RG{87.004}{57.335}{10.9091}{mi}{39}\RG{94.140}{57.335}{10.9091}{rm}{40}\RG{98.383}{57.335}{10.9091}{mi}{97}\RG{104.149}{57.335}{10.9091}{rm}{41}\RG{110.803}{57.335}{10.9091}{rm}{43}\RG{121.725}{57.335}{10.9091}{mi}{39}\RG{128.861}{57.335}{10.9091}{rm}{40}\RG{133.104}{57.335}{10.9091}{mi}{98}\RG{137.785}{57.335}{10.9091}{rm}{41}\RG{149.452}{57.335}{10.9091}{rm}{43}\RG{160.361}{57.335}{10.9091}{mi}{39}\RG{167.498}{57.335}{10.9091}{rm}{40}\RG{171.740}{57.335}{10.9091}{mi}{99}\RG{176.461}{57.335}{10.9091}{rm}{41}\RG{116.725}{47.059}{10.9092}{ex}{0}\RG{175.703}{47.059}{10.9092}{ex}{1}\RG{202.386}{47.059}{10.9092}{ex}{0}\RG{257.707}{47.059}{10.9092}{ex}{1}\RG{70.489}{38.223}{10.9091}{rm}{61}\RG{82.004}{38.223}{10.9091}{mi}{39}\RG{89.140}{38.223}{10.9091}{rm}{40}\RG{93.383}{38.223}{10.9091}{mi}{97}\RG{99.149}{38.223}{10.9091}{rm}{41}\RG{105.803}{38.223}{10.9091}{rm}{43}\RG{121.725}{38.223}{10.9091}{mi}{39}\RG{128.861}{38.223}{10.9091}{rm}{40}\RG{133.104}{38.223}{10.9091}{mi}{98}\RG{137.785}{38.223}{10.9091}{rm}{41}\RG{144.440}{38.223}{10.9091}{rm}{43}\RG{155.361}{38.223}{10.9091}{mi}{39}\RG{162.498}{38.223}{10.9091}{rm}{40}\RG{166.740}{38.223}{10.9091}{mi}{99}\RG{171.461}{38.223}{10.9091}{rm}{41}\RG{183.735}{38.223}{10.9091}{rm}{61}\RG{195.250}{38.223}{10.9091}{mi}{39}\RG{207.386}{38.223}{10.9091}{mi}{97}\RG{215.577}{38.223}{10.9091}{sy}{8}\RG{226.486}{38.223}{10.9091}{rm}{40}\RG{230.728}{38.223}{10.9091}{mi}{98}\RG{237.834}{38.223}{10.9091}{sy}{8}\RG{248.743}{38.223}{10.9091}{mi}{99}\RG{253.464}{38.223}{10.9091}{rm}{41}\RG{262.707}{38.223}{10.9091}{mi}{59}\RG{8.865}{27.947}{10.9092}{ex}{0}\RG{62.459}{27.947}{10.9092}{ex}{1}\RG{88.027}{27.947}{10.9092}{ex}{0}\RG{148.051}{27.947}{10.9092}{ex}{1}\RG{1.729}{19.111}{10.9091}{mi}{39}\RG{13.865}{19.111}{10.9091}{mi}{107}\RG{22.312}{19.111}{10.9091}{sy}{14}\RG{30.191}{19.111}{10.9091}{rm}{40}\RG{34.434}{19.111}{10.9091}{mi}{97}\RG{42.624}{19.111}{10.9091}{sy}{8}\RG{53.533}{19.111}{10.9091}{mi}{98}\RG{58.216}{19.111}{10.9091}{rm}{41}\RG{70.489}{19.111}{10.9091}{rm}{61}\RG{82.004}{19.111}{10.9091}{mi}{107}\RG{93.027}{19.111}{10.9091}{mi}{39}\RG{100.163}{19.111}{10.9091}{rm}{40}\RG{104.405}{19.111}{10.9091}{mi}{97}\RG{110.172}{19.111}{10.9091}{rm}{41}\RG{116.827}{19.111}{10.9091}{rm}{43}\RG{127.748}{19.111}{10.9091}{mi}{39}\RG{134.884}{19.111}{10.9091}{rm}{40}\RG{139.126}{19.111}{10.9091}{mi}{98}\RG{143.808}{19.111}{10.9091}{rm}{41}\RG{170.755}{8.835}{10.9092}{ex}{0}\RG{249.159}{8.835}{10.9092}{ex}{1}\RG{70.489}{-0.002}{10.9091}{rm}{61}\RG{82.004}{-0.002}{10.9091}{mi}{107}\RG{88.015}{-0.002}{10.9091}{mi}{39}\RG{95.163}{-0.002}{10.9091}{rm}{40}\RG{99.405}{-0.002}{10.9091}{mi}{97}\RG{105.172}{-0.002}{10.9091}{rm}{41}\RG{111.827}{-0.002}{10.9091}{rm}{43}\RG{122.748}{-0.002}{10.9091}{mi}{107}\RG{128.759}{-0.002}{10.9091}{mi}{39}\RG{135.907}{-0.002}{10.9091}{rm}{40}\RG{140.149}{-0.002}{10.9091}{mi}{98}\RG{144.831}{-0.002}{10.9091}{rm}{41}\RG{152.086}{-0.002}{10.9091}{rm}{61}\RG{163.619}{-0.002}{10.9091}{mi}{39}\RG{175.756}{-0.002}{10.9091}{rm}{40}\RG{179.999}{-0.002}{10.9091}{mi}{107}\RG{188.445}{-0.002}{10.9091}{sy}{14}\RG{196.324}{-0.002}{10.9091}{mi}{97}\RG{202.091}{-0.002}{10.9091}{rm}{41}\RG{208.757}{-0.002}{10.9091}{sy}{8}\RG{219.666}{-0.002}{10.9091}{rm}{40}\RG{223.909}{-0.002}{10.9091}{mi}{107}\RG{232.356}{-0.002}{10.9091}{sy}{14}\RG{240.235}{-0.002}{10.9091}{mi}{98}\RG{244.916}{-0.002}{10.9091}{rm}{41}\RG{254.159}{-0.002}{10.9091}{mi}{58}%
}
\]
所以加法满足结合律，数乘对向量加法满足分配律。同理，由\allowbreak{}\hskip.1em \RB{8.889}{7.947}{0.744}{%
\RG{0.000}{-0.004}{10.9091}{mi}{86}%
}\hskip.1em\allowbreak{} 上的公理可分别得到加法交换律、标量加法的分配律、数乘结合律和\allowbreak{}\hskip.1em \RB{42.002}{7.768}{0.618}{%
\RG{0.000}{0.002}{10.9091}{rm}{49}\RG{7.879}{0.002}{10.9091}{sy}{14}\RG{15.757}{0.002}{10.9091}{mi}{97}\RG{24.554}{0.002}{10.9091}{rm}{61}\RG{36.069}{0.002}{10.9091}{mi}{97}%
}\hskip.1em\allowbreak{} 。因此\allowbreak{}\hskip.1em \RB{7.536}{8.193}{0.738}{%
\RG{0.000}{0.002}{10.9091}{mi}{83}%
}\hskip.1em\allowbreak{} 是线性空间。由定义，

\[
\RB{218.024}{8.682}{3.227}{%
\RG{0.000}{-0.000}{10.9091}{mi}{39}\RG{7.136}{-0.000}{10.9091}{rm}{40}\RG{11.379}{-0.000}{10.9091}{mi}{97}\RG{19.569}{-0.000}{10.9091}{sy}{8}\RG{30.478}{-0.000}{10.9091}{mi}{98}\RG{35.160}{-0.000}{10.9091}{rm}{41}\RG{42.415}{-0.000}{10.9091}{rm}{61}\RG{53.948}{-0.000}{10.9091}{mi}{39}\RG{61.085}{-0.000}{10.9091}{rm}{40}\RG{65.328}{-0.000}{10.9091}{mi}{97}\RG{71.094}{-0.000}{10.9091}{rm}{41}\RG{77.748}{-0.000}{10.9091}{rm}{43}\RG{88.670}{-0.000}{10.9091}{mi}{39}\RG{95.806}{-0.000}{10.9091}{rm}{40}\RG{100.049}{-0.000}{10.9091}{mi}{98}\RG{104.730}{-0.000}{10.9091}{rm}{41}\RG{108.973}{-0.000}{10.9091}{mi}{59}\RG{135.624}{-0.000}{10.9091}{mi}{39}\RG{142.776}{-0.000}{10.9091}{rm}{40}\RG{147.018}{-0.000}{10.9091}{mi}{107}\RG{155.465}{-0.000}{10.9091}{sy}{14}\RG{163.344}{-0.000}{10.9091}{mi}{97}\RG{169.110}{-0.000}{10.9091}{rm}{41}\RG{176.365}{-0.000}{10.9091}{rm}{61}\RG{187.899}{-0.000}{10.9091}{mi}{107}\RG{193.909}{-0.000}{10.9091}{mi}{39}\RG{201.058}{-0.000}{10.9091}{rm}{40}\RG{205.301}{-0.000}{10.9091}{mi}{97}\RG{211.067}{-0.000}{10.9091}{rm}{41}\RG{215.310}{-0.000}{10.9091}{mi}{59}%
}
\]
且\allowbreak{}\hskip.1em \RB{7.253}{5.319}{2.881}{%
\RG{0.000}{-0.003}{10.9091}{mi}{39}%
}\hskip.1em\allowbreak{} 为双射，故是线性同构。(2)对\allowbreak{}\hskip.1em \RB{43.696}{8.326}{2.612}{%
\RG{0.000}{0.004}{10.9091}{mi}{80}\RG{7.298}{0.004}{10.9091}{mi}{59}\RG{12.142}{0.004}{10.9091}{mi}{81}\RG{23.810}{0.004}{10.9091}{sy}{50}\RG{34.098}{0.004}{10.9091}{sy}{84}%
}\hskip.1em\allowbreak{} 和\allowbreak{}\hskip.1em \RB{27.525}{8.075}{0.932}{%
\RG{0.000}{0.004}{10.9091}{mi}{107}\RG{9.053}{0.004}{10.9091}{sy}{50}\RG{19.356}{0.004}{10.9091}{bb}{82}%
}\hskip.1em\allowbreak{} ，记

\[
\RB{332.194}{10.305}{3.226}{%
\RG{146.012}{4.505}{7.9701}{sy}{0}\RG{152.599}{4.505}{7.9701}{rm}{49}\RG{318.281}{4.505}{7.9701}{sy}{0}\RG{324.868}{4.505}{7.9701}{rm}{49}\RG{0.000}{0.001}{10.9091}{mi}{88}\RG{12.924}{0.001}{10.9091}{rm}{61}\RG{24.422}{0.001}{10.9091}{rm}{9}\RG{32.899}{0.001}{10.9091}{rm}{40}\RG{37.167}{0.001}{10.9091}{mi}{80}\RG{45.686}{0.001}{10.9091}{rm}{41}\RG{52.941}{0.001}{10.9091}{rm}{61}\RG{64.450}{0.001}{10.9091}{rm}{40}\RG{68.717}{0.001}{10.9091}{mi}{73}\RG{81.573}{0.001}{10.9091}{sy}{0}\RG{92.482}{0.001}{10.9091}{mi}{80}\RG{101.001}{0.001}{10.9091}{rm}{41}\RG{105.234}{0.001}{10.9091}{rm}{40}\RG{109.486}{0.001}{10.9091}{mi}{73}\RG{122.342}{0.001}{10.9091}{rm}{43}\RG{133.251}{0.001}{10.9091}{mi}{80}\RG{141.770}{0.001}{10.9091}{rm}{41}\RG{157.331}{0.001}{10.9091}{mi}{59}\RG{173.073}{0.001}{10.9091}{mi}{89}\RG{184.878}{0.001}{10.9091}{rm}{61}\RG{196.376}{0.001}{10.9091}{rm}{9}\RG{204.852}{0.001}{10.9091}{rm}{40}\RG{209.120}{0.001}{10.9091}{mi}{81}\RG{217.744}{0.001}{10.9091}{rm}{41}\RG{224.998}{0.001}{10.9091}{rm}{61}\RG{236.508}{0.001}{10.9091}{rm}{40}\RG{240.774}{0.001}{10.9091}{mi}{73}\RG{253.630}{0.001}{10.9091}{sy}{0}\RG{264.540}{0.001}{10.9091}{mi}{81}\RG{273.164}{0.001}{10.9091}{rm}{41}\RG{277.397}{0.001}{10.9091}{rm}{40}\RG{281.649}{0.001}{10.9091}{mi}{73}\RG{294.505}{0.001}{10.9091}{rm}{43}\RG{305.414}{0.001}{10.9091}{mi}{81}\RG{314.039}{0.001}{10.9091}{rm}{41}\RG{329.600}{0.001}{10.9091}{mi}{58}\RG{73.512}{-1.635}{7.9701}{mi}{110}\RG{114.281}{-1.635}{7.9701}{mi}{110}\RG{245.570}{-1.635}{7.9701}{mi}{110}\RG{286.444}{-1.635}{7.9701}{mi}{110}%
}
\]
由\allowbreak{}\hskip.1em \RB{22.929}{7.950}{2.606}{%
\RG{0.000}{-0.001}{10.9091}{mi}{88}\RG{9.273}{-0.001}{10.9091}{mi}{59}\RG{14.116}{-0.001}{10.9091}{mi}{89}%
}\hskip.1em\allowbreak{} 为实反对称阵可知\allowbreak{}\hskip.1em \RB{55.805}{7.950}{2.218}{%
\RG{0.000}{-0.001}{10.9091}{mi}{73}\RG{4.795}{-1.638}{7.9701}{mi}{110}\RG{12.856}{-0.001}{10.9091}{rm}{43}\RG{23.765}{-0.001}{10.9091}{mi}{88}\RG{36.083}{-0.001}{10.9091}{rm}{43}\RG{46.992}{-0.001}{10.9091}{mi}{89}%
}\hskip.1em\allowbreak{} 、\allowbreak{}\hskip.1em \RB{39.560}{8.070}{2.218}{%
\RG{0.000}{-0.001}{10.9091}{mi}{73}\RG{4.795}{-1.638}{7.9701}{mi}{110}\RG{12.856}{-0.001}{10.9091}{rm}{43}\RG{23.765}{-0.001}{10.9091}{mi}{107}\RG{29.776}{-0.001}{10.9091}{mi}{88}%
}\hskip.1em\allowbreak{} 均可逆。因此定义

\[
\RB{250.421}{10.300}{22.343}{%
\RG{236.387}{4.500}{7.9701}{sy}{0}\RG{242.973}{4.500}{7.9701}{rm}{49}\RG{0.000}{-0.004}{10.9091}{mi}{80}\RG{10.943}{-0.004}{10.9091}{sy}{8}\RG{21.852}{-0.004}{10.9091}{mi}{81}\RG{33.506}{-0.004}{10.9091}{rm}{61}\RG{45.005}{-0.004}{10.9091}{rm}{8}\RG{52.881}{-0.004}{10.9091}{rm}{40}\RG{57.143}{-0.004}{10.9091}{mi}{88}\RG{69.462}{-0.004}{10.9091}{rm}{43}\RG{80.371}{-0.004}{10.9091}{mi}{89}\RG{89.129}{-0.004}{10.9091}{rm}{41}\RG{96.383}{-0.004}{10.9091}{rm}{61}\RG{107.893}{-0.004}{10.9091}{rm}{40}\RG{112.159}{-0.004}{10.9091}{mi}{73}\RG{125.015}{-0.004}{10.9091}{sy}{0}\RG{135.924}{-0.004}{10.9091}{mi}{88}\RG{148.242}{-0.004}{10.9091}{sy}{0}\RG{159.151}{-0.004}{10.9091}{mi}{89}\RG{167.909}{-0.004}{10.9091}{rm}{41}\RG{172.142}{-0.004}{10.9091}{rm}{40}\RG{176.394}{-0.004}{10.9091}{mi}{73}\RG{189.251}{-0.004}{10.9091}{rm}{43}\RG{200.160}{-0.004}{10.9091}{mi}{88}\RG{212.478}{-0.004}{10.9091}{rm}{43}\RG{223.387}{-0.004}{10.9091}{mi}{89}\RG{232.144}{-0.004}{10.9091}{rm}{41}\RG{247.706}{-0.004}{10.9091}{mi}{59}\RG{116.954}{-1.640}{7.9701}{mi}{110}\RG{181.189}{-1.640}{7.9701}{mi}{110}\RG{188.182}{-14.612}{7.9701}{sy}{0}\RG{194.769}{-14.612}{7.9701}{rm}{49}\RG{5.631}{-19.116}{10.9091}{mi}{107}\RG{14.078}{-19.116}{10.9091}{sy}{14}\RG{21.957}{-19.116}{10.9091}{mi}{80}\RG{33.506}{-19.116}{10.9091}{rm}{61}\RG{45.005}{-19.116}{10.9091}{rm}{8}\RG{52.881}{-19.116}{10.9091}{rm}{40}\RG{57.143}{-19.116}{10.9091}{mi}{107}\RG{63.154}{-19.116}{10.9091}{mi}{88}\RG{73.060}{-19.116}{10.9091}{rm}{41}\RG{80.314}{-19.116}{10.9091}{rm}{61}\RG{91.823}{-19.116}{10.9091}{rm}{40}\RG{96.091}{-19.116}{10.9091}{mi}{73}\RG{108.947}{-19.116}{10.9091}{sy}{0}\RG{119.856}{-19.116}{10.9091}{mi}{107}\RG{125.867}{-19.116}{10.9091}{mi}{88}\RG{135.772}{-19.116}{10.9091}{rm}{41}\RG{140.005}{-19.116}{10.9091}{rm}{40}\RG{144.257}{-19.116}{10.9091}{mi}{73}\RG{157.113}{-19.116}{10.9091}{rm}{43}\RG{168.022}{-19.116}{10.9091}{mi}{107}\RG{174.033}{-19.116}{10.9091}{mi}{88}\RG{183.940}{-19.116}{10.9091}{rm}{41}\RG{199.501}{-19.116}{10.9091}{mi}{58}\RG{100.886}{-20.753}{7.9701}{mi}{110}\RG{149.053}{-20.753}{7.9701}{mi}{110}%
}
\]
将\allowbreak{}\hskip.1em \RB{22.929}{7.952}{2.605}{%
\RG{0.000}{0.001}{10.9091}{mi}{88}\RG{9.273}{0.001}{10.9091}{mi}{59}\RG{14.116}{0.001}{10.9091}{mi}{89}%
}\hskip.1em\allowbreak{} 展开，得到只含\allowbreak{}\hskip.1em \RB{20.715}{8.192}{2.615}{%
\RG{0.000}{0.001}{10.9091}{mi}{80}\RG{7.298}{0.001}{10.9091}{mi}{59}\RG{12.142}{0.001}{10.9091}{mi}{81}%
}\hskip.1em\allowbreak{} 的显式加法公式

\[
\RB{357.595}{10.303}{3.228}{%
\RG{85.792}{4.503}{7.9701}{sy}{0}\RG{92.378}{4.503}{7.9701}{rm}{49}\RG{302.911}{4.503}{7.9701}{sy}{0}\RG{309.499}{4.503}{7.9701}{rm}{49}\RG{0.000}{-0.001}{10.9091}{mi}{80}\RG{10.943}{-0.001}{10.9091}{sy}{8}\RG{21.852}{-0.001}{10.9091}{mi}{81}\RG{33.507}{-0.001}{10.9091}{rm}{61}\RG{45.006}{-0.001}{10.9091}{rm}{40}\RG{49.264}{-0.001}{10.9091}{mi}{73}\RG{62.120}{-0.001}{10.9091}{rm}{43}\RG{73.031}{-0.001}{10.9091}{mi}{80}\RG{81.550}{-0.001}{10.9091}{rm}{41}\RG{97.111}{-0.001}{10.9091}{rm}{40}\RG{101.353}{-0.001}{10.9091}{sy}{0}\RG{109.838}{-0.001}{10.9091}{mi}{73}\RG{122.694}{-0.001}{10.9091}{rm}{43}\RG{133.603}{-0.001}{10.9091}{mi}{80}\RG{144.546}{-0.001}{10.9091}{rm}{43}\RG{155.455}{-0.001}{10.9091}{mi}{81}\RG{166.504}{-0.001}{10.9091}{rm}{43}\RG{177.403}{-0.001}{10.9091}{rm}{51}\RG{182.867}{-0.001}{10.9091}{mi}{80}\RG{191.377}{-0.001}{10.9091}{mi}{81}\RG{200.012}{-0.001}{10.9091}{rm}{41}\RG{204.244}{-0.001}{10.9091}{rm}{40}\RG{208.477}{-0.001}{10.9091}{rm}{51}\RG{213.951}{-0.001}{10.9091}{mi}{73}\RG{226.807}{-0.001}{10.9091}{rm}{43}\RG{237.716}{-0.001}{10.9091}{mi}{80}\RG{248.659}{-0.001}{10.9091}{rm}{43}\RG{259.568}{-0.001}{10.9091}{mi}{81}\RG{270.617}{-0.001}{10.9091}{sy}{0}\RG{281.526}{-0.001}{10.9091}{mi}{80}\RG{290.035}{-0.001}{10.9091}{mi}{81}\RG{298.669}{-0.001}{10.9091}{rm}{41}\RG{314.231}{-0.001}{10.9091}{rm}{40}\RG{318.474}{-0.001}{10.9091}{mi}{73}\RG{331.330}{-0.001}{10.9091}{rm}{43}\RG{342.239}{-0.001}{10.9091}{mi}{80}\RG{350.758}{-0.001}{10.9091}{rm}{41}\RG{355.000}{-0.001}{10.9091}{mi}{58}\RG{54.060}{-1.637}{7.9701}{mi}{110}\RG{114.634}{-1.637}{7.9701}{mi}{110}\RG{218.747}{-1.637}{7.9701}{mi}{110}\RG{323.269}{-1.637}{7.9701}{mi}{110}%
}
\]
实际上，令\allowbreak{}\hskip.1em \RB{123.097}{8.187}{2.620}{%
\RG{0.000}{-0.004}{10.9091}{mi}{78}\RG{12.985}{-0.004}{10.9091}{rm}{61}\RG{24.500}{-0.004}{10.9091}{sy}{0}\RG{32.985}{-0.004}{10.9091}{mi}{73}\RG{37.780}{-1.640}{7.9701}{mi}{110}\RG{45.841}{-0.004}{10.9091}{rm}{43}\RG{56.750}{-0.004}{10.9091}{mi}{80}\RG{67.694}{-0.004}{10.9091}{rm}{43}\RG{78.603}{-0.004}{10.9091}{mi}{81}\RG{89.651}{-0.004}{10.9091}{rm}{43}\RG{100.549}{-0.004}{10.9091}{rm}{51}\RG{106.015}{-0.004}{10.9091}{mi}{80}\RG{114.524}{-0.004}{10.9091}{mi}{81}%
}\hskip.1em\allowbreak{} 、\allowbreak{}\hskip.1em \RB{113.993}{8.187}{2.620}{%
\RG{0.000}{-0.004}{10.9091}{mi}{68}\RG{12.365}{-0.004}{10.9091}{rm}{61}\RG{23.863}{-0.004}{10.9091}{rm}{51}\RG{29.335}{-0.004}{10.9091}{mi}{73}\RG{34.131}{-1.640}{7.9701}{mi}{110}\RG{42.191}{-0.004}{10.9091}{rm}{43}\RG{53.100}{-0.004}{10.9091}{mi}{80}\RG{64.044}{-0.004}{10.9091}{rm}{43}\RG{74.953}{-0.004}{10.9091}{mi}{81}\RG{86.002}{-0.004}{10.9091}{sy}{0}\RG{96.911}{-0.004}{10.9091}{mi}{80}\RG{105.420}{-0.004}{10.9091}{mi}{81}%
}\hskip.1em\allowbreak{} ，直接展开得到

\[
\RB{173.094}{8.686}{3.223}{%
\RG{0.000}{0.004}{10.9091}{mi}{78}\RG{12.985}{0.004}{10.9091}{rm}{61}\RG{24.483}{0.004}{10.9091}{rm}{40}\RG{28.742}{0.004}{10.9091}{mi}{73}\RG{41.598}{0.004}{10.9091}{rm}{43}\RG{52.507}{0.004}{10.9091}{mi}{80}\RG{61.027}{0.004}{10.9091}{rm}{41}\RG{65.259}{0.004}{10.9091}{rm}{40}\RG{69.512}{0.004}{10.9091}{mi}{73}\RG{82.368}{0.004}{10.9091}{sy}{0}\RG{93.277}{0.004}{10.9091}{mi}{88}\RG{105.595}{0.004}{10.9091}{sy}{0}\RG{116.504}{0.004}{10.9091}{mi}{89}\RG{125.262}{0.004}{10.9091}{rm}{41}\RG{129.495}{0.004}{10.9091}{rm}{40}\RG{133.747}{0.004}{10.9091}{mi}{73}\RG{146.603}{0.004}{10.9091}{rm}{43}\RG{157.512}{0.004}{10.9091}{mi}{81}\RG{166.136}{0.004}{10.9091}{rm}{41}\RG{170.379}{0.004}{10.9091}{mi}{59}\RG{33.538}{-1.632}{7.9701}{mi}{110}\RG{74.308}{-1.632}{7.9701}{mi}{110}\RG{138.542}{-1.632}{7.9701}{mi}{110}%
}
\]
\[
\RB{172.474}{8.683}{3.226}{%
\RG{0.000}{0.001}{10.9091}{mi}{68}\RG{12.365}{0.001}{10.9091}{rm}{61}\RG{23.863}{0.001}{10.9091}{rm}{40}\RG{28.122}{0.001}{10.9091}{mi}{73}\RG{40.978}{0.001}{10.9091}{rm}{43}\RG{51.887}{0.001}{10.9091}{mi}{80}\RG{60.407}{0.001}{10.9091}{rm}{41}\RG{64.639}{0.001}{10.9091}{rm}{40}\RG{68.892}{0.001}{10.9091}{mi}{73}\RG{81.748}{0.001}{10.9091}{rm}{43}\RG{92.657}{0.001}{10.9091}{mi}{88}\RG{104.975}{0.001}{10.9091}{rm}{43}\RG{115.885}{0.001}{10.9091}{mi}{89}\RG{124.642}{0.001}{10.9091}{rm}{41}\RG{128.875}{0.001}{10.9091}{rm}{40}\RG{133.127}{0.001}{10.9091}{mi}{73}\RG{145.983}{0.001}{10.9091}{rm}{43}\RG{156.892}{0.001}{10.9091}{mi}{81}\RG{165.516}{0.001}{10.9091}{rm}{41}\RG{169.759}{0.001}{10.9091}{mi}{59}\RG{32.918}{-1.635}{7.9701}{mi}{110}\RG{73.688}{-1.635}{7.9701}{mi}{110}\RG{137.922}{-1.635}{7.9701}{mi}{110}%
}
\]
% Original page 2.
于是所给公式成立，且\allowbreak{}\hskip.1em \RB{9.271}{7.950}{0.501}{%
\RG{0.000}{-0.001}{10.9091}{mi}{68}%
}\hskip.1em\allowbreak{} 可逆。数乘可化为

\[
\RB{269.238}{11.915}{4.311}{%
\RG{39.390}{8.833}{10.9092}{ex}{2}\RG{142.812}{8.833}{10.9092}{ex}{3}\RG{147.350}{8.833}{10.9092}{ex}{2}\RG{250.779}{8.833}{10.9092}{ex}{3}\RG{255.324}{6.115}{7.9701}{sy}{0}\RG{261.911}{6.115}{7.9701}{rm}{49}\RG{0.000}{-0.003}{10.9091}{mi}{107}\RG{8.447}{-0.003}{10.9091}{sy}{14}\RG{16.326}{-0.003}{10.9091}{mi}{80}\RG{27.875}{-0.003}{10.9091}{rm}{61}\RG{43.936}{-0.003}{10.9091}{rm}{40}\RG{48.169}{-0.003}{10.9091}{rm}{49}\RG{56.057}{-0.003}{10.9091}{sy}{0}\RG{66.967}{-0.003}{10.9091}{mi}{107}\RG{72.990}{-0.003}{10.9091}{rm}{41}\RG{77.232}{-0.003}{10.9091}{mi}{73}\RG{90.088}{-0.003}{10.9091}{rm}{43}\RG{100.986}{-0.003}{10.9091}{rm}{40}\RG{105.219}{-0.003}{10.9091}{rm}{49}\RG{113.095}{-0.003}{10.9091}{rm}{43}\RG{124.028}{-0.003}{10.9091}{mi}{107}\RG{130.050}{-0.003}{10.9091}{rm}{41}\RG{134.293}{-0.003}{10.9091}{mi}{80}\RG{151.903}{-0.003}{10.9091}{rm}{40}\RG{156.136}{-0.003}{10.9091}{rm}{49}\RG{164.012}{-0.003}{10.9091}{rm}{43}\RG{174.933}{-0.003}{10.9091}{mi}{107}\RG{180.957}{-0.003}{10.9091}{rm}{41}\RG{185.199}{-0.003}{10.9091}{mi}{73}\RG{198.055}{-0.003}{10.9091}{rm}{43}\RG{208.953}{-0.003}{10.9091}{rm}{40}\RG{213.186}{-0.003}{10.9091}{rm}{49}\RG{221.085}{-0.003}{10.9091}{sy}{0}\RG{231.995}{-0.003}{10.9091}{mi}{107}\RG{238.017}{-0.003}{10.9091}{rm}{41}\RG{242.260}{-0.003}{10.9091}{mi}{80}\RG{266.643}{-0.003}{10.9091}{mi}{58}\RG{82.028}{-1.640}{7.9701}{mi}{110}\RG{189.995}{-1.640}{7.9701}{mi}{110}%
}
\]
由构造有

\[
\RB{243.542}{8.683}{3.226}{%
\RG{0.000}{0.001}{10.9091}{rm}{9}\RG{8.476}{0.001}{10.9091}{rm}{40}\RG{12.727}{0.001}{10.9091}{mi}{80}\RG{23.670}{0.001}{10.9091}{sy}{8}\RG{34.580}{0.001}{10.9091}{mi}{81}\RG{43.204}{0.001}{10.9091}{rm}{41}\RG{50.459}{0.001}{10.9091}{rm}{61}\RG{61.968}{0.001}{10.9091}{rm}{9}\RG{70.445}{0.001}{10.9091}{rm}{40}\RG{74.720}{0.001}{10.9091}{mi}{80}\RG{83.239}{0.001}{10.9091}{rm}{41}\RG{89.893}{0.001}{10.9091}{rm}{43}\RG{100.792}{0.001}{10.9091}{rm}{9}\RG{109.268}{0.001}{10.9091}{rm}{40}\RG{113.542}{0.001}{10.9091}{mi}{81}\RG{122.166}{0.001}{10.9091}{rm}{41}\RG{126.408}{0.001}{10.9091}{mi}{59}\RG{153.075}{0.001}{10.9091}{rm}{9}\RG{161.552}{0.001}{10.9091}{rm}{40}\RG{165.802}{0.001}{10.9091}{mi}{107}\RG{174.249}{0.001}{10.9091}{sy}{14}\RG{182.128}{0.001}{10.9091}{mi}{80}\RG{190.648}{0.001}{10.9091}{rm}{41}\RG{197.902}{0.001}{10.9091}{rm}{61}\RG{209.436}{0.001}{10.9091}{mi}{107}\RG{215.459}{0.001}{10.9091}{rm}{9}\RG{223.935}{0.001}{10.9091}{rm}{40}\RG{228.186}{0.001}{10.9091}{mi}{80}\RG{236.705}{0.001}{10.9091}{rm}{41}\RG{240.947}{0.001}{10.9091}{mi}{58}%
}
\]
因此\allowbreak{}\hskip.1em \RB{9.598}{8.323}{1.241}{%
\RG{0.000}{0.001}{10.9091}{sy}{84}%
}\hskip.1em\allowbreak{} 是实线性空间，而\allowbreak{}\hskip.1em \RB{50.749}{8.323}{1.241}{%
\RG{0.000}{0.001}{10.9091}{rm}{9}\RG{11.498}{0.001}{10.9091}{rm}{58}\RG{17.576}{0.001}{10.9091}{sy}{84}\RG{29.304}{0.001}{10.9091}{sy}{33}\RG{43.245}{0.001}{10.9091}{sy}{83}%
}\hskip.1em\allowbreak{} 为线性同构。其加法零元为\allowbreak{}\hskip.1em \RB{10.173}{7.952}{2.215}{%
\RG{0.000}{0.001}{10.9091}{mi}{73}\RG{4.795}{-1.635}{7.9701}{mi}{110}%
}\hskip.1em\allowbreak{} ，\allowbreak{}\hskip.1em \RB{8.726}{7.952}{0.500}{%
\RG{0.000}{0.001}{10.9091}{mi}{80}%
}\hskip.1em\allowbreak{} 的加法逆元为\allowbreak{}\hskip.1em \RB{43.110}{9.760}{0.500}{%
\RG{0.000}{0.001}{10.9091}{mi}{80}\RG{8.519}{3.960}{7.9701}{sy}{0}\RG{15.106}{3.960}{7.9701}{rm}{49}\RG{22.868}{0.001}{10.9091}{rm}{61}\RG{34.384}{0.001}{10.9091}{mi}{80}%
}\hskip.1em\allowbreak{} T。

\[
\RB{8.366}{7.993}{0.502}{%
\RG{0.000}{-0.002}{10.9091}{am}{3}%
}
\]
方法提炼：以双射搬运代数结构。新定义的加法、数乘不必直接逐项验证八条公理；将它们映入已知线性空间后，由双射的单射性恢复全部等式。矩阵公式则须保持因子的乘法顺序。

题目来源：谢启鸿《复旦大学高等代数习题集（2023版）》，第108–109页。

% Original page 3.
\section*{2021A11}
设\allowbreak{}\hskip.1em \RB{69.752}{8.686}{3.223}{%
\RG{0.000}{0.004}{10.9091}{mi}{65}\RG{11.212}{0.004}{10.9091}{sy}{50}\RG{21.515}{0.004}{10.9091}{mi}{77}\RG{32.098}{-1.633}{7.9701}{mi}{109}\RG{39.589}{-1.633}{7.9701}{sy}{2}\RG{46.175}{-1.633}{7.9701}{mi}{110}\RG{51.811}{0.004}{10.9091}{rm}{40}\RG{56.054}{0.004}{10.9091}{mi}{75}\RG{66.100}{0.004}{10.9091}{rm}{41}%
}\hskip.1em\allowbreak{} ，rank\allowbreak{}\hskip.1em \RB{36.469}{8.686}{3.223}{%
\RG{0.000}{0.004}{10.9091}{rm}{40}\RG{4.242}{0.004}{10.9091}{mi}{65}\RG{12.424}{0.004}{10.9091}{rm}{41}\RG{19.679}{0.004}{10.9091}{rm}{61}\RG{31.212}{0.004}{10.9091}{mi}{114}%
}\hskip.1em\allowbreak{} 。证明：\allowbreak{}\hskip.1em \RB{8.366}{8.315}{0.500}{%
\RG{0.000}{0.004}{10.9091}{mi}{65}%
}\hskip.1em\allowbreak{} 可写成\allowbreak{}\hskip.1em \RB{40.333}{8.315}{0.736}{%
\RG{0.000}{0.004}{10.9091}{mi}{65}\RG{11.212}{0.004}{10.9091}{rm}{61}\RG{22.727}{0.004}{10.9091}{mi}{66}\RG{31.542}{0.004}{10.9091}{mi}{67}%
}\hskip.1em\allowbreak{} ，其中\allowbreak{}\hskip.1em \RB{69.311}{8.686}{3.223}{%
\RG{0.000}{0.004}{10.9091}{mi}{66}\RG{11.852}{0.004}{10.9091}{sy}{50}\RG{22.155}{0.004}{10.9091}{mi}{77}\RG{32.738}{-1.633}{7.9701}{mi}{109}\RG{40.229}{-1.633}{7.9701}{sy}{2}\RG{46.816}{-1.633}{7.9701}{mi}{114}\RG{51.370}{0.004}{10.9091}{rm}{40}\RG{55.613}{0.004}{10.9091}{mi}{75}\RG{65.659}{0.004}{10.9091}{rm}{41}%
}\hskip.1em\allowbreak{} 、\allowbreak{}\hskip.1em \RB{66.715}{8.686}{3.223}{%
\RG{0.000}{0.004}{10.9091}{mi}{67}\RG{11.608}{0.004}{10.9091}{sy}{50}\RG{21.911}{0.004}{10.9091}{mi}{77}\RG{32.494}{-1.633}{7.9701}{mi}{114}\RG{36.551}{-1.633}{7.9701}{sy}{2}\RG{43.137}{-1.633}{7.9701}{mi}{110}\RG{48.774}{0.004}{10.9091}{rm}{40}\RG{53.016}{0.004}{10.9091}{mi}{75}\RG{63.063}{0.004}{10.9091}{rm}{41}%
}\hskip.1em\allowbreak{} ，且这种分解的成立等价于可选取\allowbreak{}\hskip.1em \RB{8.747}{7.951}{0.500}{%
\RG{0.000}{-0.000}{10.9091}{mi}{66}%
}\hskip.1em\allowbreak{} 的\allowbreak{}\hskip.1em \RB{5.256}{5.322}{0.620}{%
\RG{0.000}{-0.000}{10.9091}{mi}{114}%
}\hskip.1em\allowbreak{} 个列向量作为\allowbreak{}\hskip.1em \RB{8.366}{8.311}{0.500}{%
\RG{0.000}{-0.000}{10.9091}{mi}{65}%
}\hskip.1em\allowbreak{} 的列空间的一组基；也等价于可选取\allowbreak{}\hskip.1em \RB{8.791}{8.191}{0.740}{%
\RG{0.000}{-0.000}{10.9091}{mi}{67}%
}\hskip.1em\allowbreak{}的\allowbreak{}\hskip.1em \RB{5.256}{5.318}{0.624}{%
\RG{0.000}{-0.004}{10.9091}{mi}{114}%
}\hskip.1em\allowbreak{} 个行向量作为\allowbreak{}\hskip.1em \RB{8.366}{8.307}{0.504}{%
\RG{0.000}{-0.004}{10.9091}{mi}{65}%
}\hskip.1em\allowbreak{} 的行空间的一组基。此时称\allowbreak{}\hskip.1em \RB{40.333}{8.307}{0.744}{%
\RG{0.000}{-0.004}{10.9091}{mi}{65}\RG{11.212}{-0.004}{10.9091}{rm}{61}\RG{22.727}{-0.004}{10.9091}{mi}{66}\RG{31.542}{-0.004}{10.9091}{mi}{67}%
}\hskip.1em\allowbreak{} 为满秩分解：

解：先设\allowbreak{}\hskip.1em \RB{24.841}{7.770}{1.991}{%
\RG{0.000}{0.004}{10.9091}{mi}{114}\RG{8.255}{0.004}{10.9091}{sy}{21}\RG{19.770}{0.004}{10.9091}{rm}{49}%
}\hskip.1em\allowbreak{} 。若\allowbreak{}\hskip.1em \RB{40.333}{8.315}{0.736}{%
\RG{0.000}{0.004}{10.9091}{mi}{65}\RG{11.212}{0.004}{10.9091}{rm}{61}\RG{22.727}{0.004}{10.9091}{mi}{66}\RG{31.542}{0.004}{10.9091}{mi}{67}%
}\hskip.1em\allowbreak{} ，则由矩阵乘积的秩不等式，

\[
\RB{267.717}{9.295}{3.225}{%
\RG{0.000}{0.002}{10.9091}{mi}{114}\RG{8.255}{0.002}{10.9091}{rm}{61}\RT{19.770}{0.002}{10.9091}{r}\RT{24.047}{0.002}{10.9091}{a}\RT{29.501}{0.002}{10.9091}{n}\RT{35.567}{0.002}{10.9091}{k}\RG{41.327}{0.002}{10.9091}{rm}{40}\RG{45.569}{0.002}{10.9091}{mi}{65}\RG{53.751}{0.002}{10.9091}{rm}{41}\RG{61.006}{0.002}{10.9091}{rm}{61}\RT{72.540}{0.002}{10.9091}{r}\RT{76.816}{0.002}{10.9091}{a}\RT{82.271}{0.002}{10.9091}{n}\RT{88.336}{0.002}{10.9091}{k}\RG{94.096}{0.002}{10.9091}{rm}{40}\RG{98.339}{0.002}{10.9091}{mi}{66}\RG{107.153}{0.002}{10.9091}{mi}{67}\RG{115.738}{0.002}{10.9091}{rm}{41}\RG{123.010}{0.002}{10.9091}{sy}{20}\RT{134.526}{0.002}{10.9091}{m}\RT{143.613}{0.002}{10.9091}{i}\RT{146.646}{0.002}{10.9091}{n}\RG{152.711}{0.002}{10.9091}{sy}{102}\RT{158.166}{0.002}{10.9091}{r}\RT{162.443}{0.002}{10.9091}{a}\RT{167.897}{0.002}{10.9091}{n}\RT{173.963}{0.002}{10.9091}{k}\RG{179.722}{0.002}{10.9091}{rm}{40}\RG{183.965}{0.002}{10.9091}{mi}{66}\RG{192.787}{0.002}{10.9091}{rm}{41}\RG{197.030}{0.002}{10.9091}{mi}{59}\RT{201.878}{0.002}{10.9091}{r}\RT{206.154}{0.002}{10.9091}{a}\RT{211.609}{0.002}{10.9091}{n}\RT{217.674}{0.002}{10.9091}{k}\RG{223.435}{0.002}{10.9091}{rm}{40}\RG{227.677}{0.002}{10.9091}{mi}{67}\RG{236.254}{0.002}{10.9091}{rm}{41}\RG{240.497}{0.002}{10.9091}{sy}{103}\RG{248.973}{0.002}{10.9091}{sy}{20}\RG{260.497}{0.002}{10.9091}{mi}{114}\RG{265.123}{0.002}{10.9091}{mi}{58}%
}
\]
所以rank\allowbreak{}\hskip.1em \RB{28.685}{8.681}{3.228}{%
\RG{0.000}{-0.001}{10.9091}{rm}{40}\RG{4.243}{-0.001}{10.9091}{mi}{66}\RG{13.065}{-0.001}{10.9091}{rm}{41}\RG{20.320}{-0.001}{10.9091}{rm}{61}%
}\hskip.1em\allowbreak{} rank\allowbreak{}\hskip.1em \RB{36.864}{8.681}{3.228}{%
\RG{0.000}{-0.001}{10.9091}{rm}{40}\RG{4.242}{-0.001}{10.9091}{mi}{67}\RG{12.819}{-0.001}{10.9091}{rm}{41}\RG{20.074}{-0.001}{10.9091}{rm}{61}\RG{31.607}{-0.001}{10.9091}{mi}{114}%
}\hskip.1em\allowbreak{} 。另一方面，

\[
\RB{201.503}{9.290}{3.230}{%
\RT{0.000}{-0.003}{10.9091}{C}\RT{7.876}{-0.003}{10.9091}{o}\RT{13.331}{-0.003}{10.9091}{l}\RG{16.364}{-0.003}{10.9091}{rm}{40}\RG{20.606}{-0.003}{10.9091}{mi}{65}\RG{28.788}{-0.003}{10.9091}{rm}{41}\RG{36.061}{-0.003}{10.9091}{sy}{18}\RT{47.576}{-0.003}{10.9091}{C}\RT{55.453}{-0.003}{10.9091}{o}\RT{60.907}{-0.003}{10.9091}{l}\RG{63.941}{-0.003}{10.9091}{rm}{40}\RG{68.183}{-0.003}{10.9091}{mi}{66}\RG{77.005}{-0.003}{10.9091}{rm}{41}\RG{81.247}{-0.003}{10.9091}{mi}{59}\RT{107.914}{-0.003}{10.9091}{R}\RT{115.943}{-0.003}{10.9091}{o}\RT{121.398}{-0.003}{10.9091}{w}\RG{129.274}{-0.003}{10.9091}{rm}{40}\RG{133.516}{-0.003}{10.9091}{mi}{65}\RG{141.698}{-0.003}{10.9091}{rm}{41}\RG{148.971}{-0.003}{10.9091}{sy}{18}\RT{160.486}{-0.003}{10.9091}{R}\RT{168.516}{-0.003}{10.9091}{o}\RT{173.970}{-0.003}{10.9091}{w}\RG{181.847}{-0.003}{10.9091}{rm}{40}\RG{186.090}{-0.003}{10.9091}{mi}{67}\RG{194.667}{-0.003}{10.9091}{rm}{41}\RG{198.909}{-0.003}{10.9091}{mi}{58}%
}
\]
由于两组包含关系两边的空间均为\allowbreak{}\hskip.1em \RB{5.256}{5.326}{0.616}{%
\RG{0.000}{0.004}{10.9091}{mi}{114}%
}\hskip.1em\allowbreak{} 维，必有

\[
\RB{201.503}{9.295}{3.225}{%
\RT{0.000}{0.002}{10.9091}{C}\RT{7.876}{0.002}{10.9091}{o}\RT{13.331}{0.002}{10.9091}{l}\RG{16.364}{0.002}{10.9091}{rm}{40}\RG{20.606}{0.002}{10.9091}{mi}{65}\RG{28.788}{0.002}{10.9091}{rm}{41}\RG{36.043}{0.002}{10.9091}{rm}{61}\RT{47.576}{0.002}{10.9091}{C}\RT{55.453}{0.002}{10.9091}{o}\RT{60.907}{0.002}{10.9091}{l}\RG{63.941}{0.002}{10.9091}{rm}{40}\RG{68.183}{0.002}{10.9091}{mi}{66}\RG{77.005}{0.002}{10.9091}{rm}{41}\RG{81.247}{0.002}{10.9091}{mi}{59}\RT{107.914}{0.002}{10.9091}{R}\RT{115.943}{0.002}{10.9091}{o}\RT{121.398}{0.002}{10.9091}{w}\RG{129.274}{0.002}{10.9091}{rm}{40}\RG{133.516}{0.002}{10.9091}{mi}{65}\RG{141.698}{0.002}{10.9091}{rm}{41}\RG{148.953}{0.002}{10.9091}{rm}{61}\RT{160.486}{0.002}{10.9091}{R}\RT{168.515}{0.002}{10.9091}{o}\RT{173.970}{0.002}{10.9091}{w}\RG{181.847}{0.002}{10.9091}{rm}{40}\RG{186.090}{0.002}{10.9091}{mi}{67}\RG{194.667}{0.002}{10.9091}{rm}{41}\RG{198.909}{0.002}{10.9091}{mi}{58}%
}
\]
从而\allowbreak{}\hskip.1em \RB{8.747}{7.950}{0.501}{%
\RG{0.000}{-0.001}{10.9091}{mi}{66}%
}\hskip.1em\allowbreak{} 的\allowbreak{}\hskip.1em \RB{5.256}{5.321}{0.621}{%
\RG{0.000}{-0.001}{10.9091}{mi}{114}%
}\hskip.1em\allowbreak{} 个列向量线性无关，构成\allowbreak{}\hskip.1em \RB{8.366}{8.310}{0.501}{%
\RG{0.000}{-0.001}{10.9091}{mi}{65}%
}\hskip.1em\allowbreak{} 的列空间的一组基；\allowbreak{}\hskip.1em \RB{8.791}{8.190}{0.741}{%
\RG{0.000}{-0.001}{10.9091}{mi}{67}%
}\hskip.1em\allowbreak{} 的\allowbreak{}\hskip.1em \RB{5.256}{5.321}{0.621}{%
\RG{0.000}{-0.001}{10.9091}{mi}{114}%
}\hskip.1em\allowbreak{} 个行向量亦构成\allowbreak{}\hskip.1em \RB{8.366}{8.310}{0.501}{%
\RG{0.000}{-0.001}{10.9091}{mi}{65}%
}\hskip.1em\allowbreak{} 的行空间的一组基。

反之，若预先选取Col\allowbreak{}\hskip.1em \RB{16.077}{8.684}{3.225}{%
\RG{0.000}{0.002}{10.9091}{rm}{40}\RG{4.243}{0.002}{10.9091}{mi}{65}\RG{12.424}{0.002}{10.9091}{rm}{41}%
}\hskip.1em\allowbreak{} 的一组基\allowbreak{}\hskip.1em \RB{45.498}{8.204}{2.614}{%
\RG{0.000}{0.002}{10.9091}{mi}{12}\RG{6.170}{-1.634}{7.9701}{rm}{49}\RG{10.902}{0.002}{10.9091}{mi}{59}\RG{15.735}{0.002}{10.9091}{mi}{58}\RG{20.579}{0.002}{10.9091}{mi}{58}\RG{25.422}{0.002}{10.9091}{mi}{58}\RG{30.255}{0.002}{10.9091}{mi}{59}\RG{35.099}{0.002}{10.9091}{mi}{12}\RG{41.315}{-1.634}{7.9701}{mi}{114}%
}\hskip.1em\allowbreak{} ，令

\[
\RB{186.765}{8.682}{3.227}{%
\RG{0.000}{0.000}{10.9091}{mi}{65}\RG{11.212}{0.000}{10.9091}{rm}{61}\RG{22.710}{0.000}{10.9091}{rm}{40}\RG{26.970}{0.000}{10.9091}{mi}{11}\RG{38.681}{0.000}{10.9091}{mi}{59}\RG{43.514}{0.000}{10.9091}{mi}{58}\RG{48.358}{0.000}{10.9091}{mi}{58}\RG{53.191}{0.000}{10.9091}{mi}{58}\RG{58.034}{0.000}{10.9091}{mi}{59}\RG{62.878}{0.000}{10.9091}{mi}{11}\RG{75.539}{0.000}{10.9091}{rm}{41}\RG{79.781}{0.000}{10.9091}{mi}{59}\RG{106.432}{0.000}{10.9091}{mi}{66}\RG{118.300}{0.000}{10.9091}{rm}{61}\RG{129.798}{0.000}{10.9091}{rm}{40}\RG{134.058}{0.000}{10.9091}{mi}{12}\RG{144.960}{0.000}{10.9091}{mi}{59}\RG{149.793}{0.000}{10.9091}{mi}{58}\RG{154.637}{0.000}{10.9091}{mi}{58}\RG{159.480}{0.000}{10.9091}{mi}{58}\RG{164.313}{0.000}{10.9091}{mi}{59}\RG{169.157}{0.000}{10.9091}{mi}{12}\RG{179.928}{0.000}{10.9091}{rm}{41}\RG{184.171}{0.000}{10.9091}{mi}{58}\RG{33.948}{-1.637}{7.9701}{rm}{49}\RG{69.902}{-1.637}{7.9701}{mi}{110}\RG{140.228}{-1.637}{7.9701}{rm}{49}\RG{175.373}{-1.637}{7.9701}{mi}{114}%
}
\]
对每个\allowbreak{}\hskip.1em \RB{4.831}{7.708}{2.739}{%
\RG{0.000}{-0.003}{10.9091}{mi}{106}%
}\hskip.1em\allowbreak{} ，存在唯一的\allowbreak{}\hskip.1em \RB{74.272}{7.948}{3.765}{%
\RG{0.000}{-0.003}{10.9091}{mi}{99}\RG{4.721}{-1.639}{7.9701}{rm}{49}\RG{8.955}{-1.639}{7.9701}{mi}{106}\RG{13.337}{-0.003}{10.9091}{mi}{59}\RG{18.170}{-0.003}{10.9091}{mi}{58}\RG{23.014}{-0.003}{10.9091}{mi}{58}\RG{27.857}{-0.003}{10.9091}{mi}{58}\RG{32.690}{-0.003}{10.9091}{mi}{59}\RG{37.534}{-0.003}{10.9091}{mi}{99}\RG{42.301}{-1.639}{7.9701}{mi}{114}\RG{46.350}{-1.639}{7.9701}{mi}{106}\RG{53.770}{-0.003}{10.9091}{sy}{50}\RG{64.074}{-0.003}{10.9091}{mi}{75}%
}\hskip.1em\allowbreak{} ，使得

\[
\RB{67.615}{17.648}{13.670}{%
\RG{31.756}{13.633}{7.9701}{mi}{114}\RG{25.906}{10.361}{10.9092}{ex}{88}\RG{0.000}{-0.003}{10.9091}{mi}{11}\RG{14.391}{-0.003}{10.9091}{rm}{61}\RG{43.482}{-0.003}{10.9091}{mi}{99}\RG{55.468}{-0.003}{10.9091}{mi}{12}\RG{65.020}{-0.003}{10.9091}{mi}{58}\RG{6.978}{-1.639}{7.9701}{mi}{106}\RG{48.203}{-1.639}{7.9701}{mi}{105}\RG{51.081}{-1.639}{7.9701}{mi}{106}\RG{61.638}{-1.639}{7.9701}{mi}{105}\RG{26.933}{-13.090}{7.9701}{mi}{105}\RG{29.816}{-13.090}{7.9701}{rm}{61}\RG{36.400}{-13.090}{7.9701}{rm}{49}%
}
\]
令\allowbreak{}\hskip.1em \RB{60.398}{8.681}{3.763}{%
\RG{0.000}{-0.001}{10.9091}{mi}{67}\RG{11.999}{-0.001}{10.9091}{rm}{61}\RG{23.890}{-0.001}{10.9091}{rm}{40}\RG{28.149}{-0.001}{10.9091}{mi}{99}\RG{32.870}{-1.637}{7.9701}{mi}{105}\RG{35.748}{-1.637}{7.9701}{mi}{106}\RG{40.135}{-0.001}{10.9091}{rm}{41}\RG{44.377}{-1.637}{7.9701}{mi}{114}\RG{48.434}{-1.637}{7.9701}{sy}{2}\RG{55.020}{-1.637}{7.9701}{mi}{110}%
}\hskip.1em\allowbreak{} ，逐列比较便得到\allowbreak{}\hskip.1em \RB{41.117}{8.310}{0.741}{%
\RG{0.000}{-0.001}{10.9091}{mi}{65}\RG{11.604}{-0.001}{10.9091}{rm}{61}\RG{23.511}{-0.001}{10.9091}{mi}{66}\RG{32.326}{-0.001}{10.9091}{mi}{67}%
}\hskip.1em\allowbreak{} 。如果改从\allowbreak{}\hskip.1em \RB{8.366}{8.310}{0.501}{%
\RG{0.000}{-0.001}{10.9091}{mi}{65}%
}\hskip.1em\allowbreak{} 的行空间的一组基出发，将基向量按行排列成\allowbreak{}\hskip.1em \RB{8.791}{8.187}{0.744}{%
\RG{0.000}{-0.004}{10.9091}{mi}{67}%
}\hskip.1em\allowbreak{} ，再将\allowbreak{}\hskip.1em \RB{8.366}{8.307}{0.504}{%
\RG{0.000}{-0.004}{10.9091}{mi}{65}%
}\hskip.1em\allowbreak{} 的每一行在该基下的坐标组成\allowbreak{}\hskip.1em \RB{8.747}{7.947}{0.504}{%
\RG{0.000}{-0.004}{10.9091}{mi}{66}%
}\hskip.1em\allowbreak{} ，也得到\allowbreak{}\hskip.1em \RB{40.333}{8.307}{0.744}{%
\RG{0.000}{-0.004}{10.9091}{mi}{65}\RG{11.212}{-0.004}{10.9091}{rm}{61}\RG{22.727}{-0.004}{10.9091}{mi}{66}\RG{31.542}{-0.004}{10.9091}{mi}{67}%
}\hskip.1em\allowbreak{} 。

两种构造所得的\allowbreak{}\hskip.1em \RB{22.438}{8.193}{2.603}{%
\RG{0.000}{0.002}{10.9091}{mi}{66}\RG{8.815}{0.002}{10.9091}{mi}{59}\RG{13.647}{0.002}{10.9091}{mi}{67}%
}\hskip.1em\allowbreak{} 都满足rank\allowbreak{}\hskip.1em \RB{29.088}{8.684}{3.225}{%
\RG{0.000}{0.002}{10.9091}{rm}{40}\RG{4.242}{0.002}{10.9091}{mi}{66}\RG{13.064}{0.002}{10.9091}{rm}{41}\RG{20.722}{0.002}{10.9091}{rm}{61}%
}\hskip.1em\allowbreak{} rank\allowbreak{}\hskip.1em \RB{37.661}{8.684}{3.225}{%
\RG{0.000}{0.002}{10.9091}{rm}{40}\RG{4.242}{0.002}{10.9091}{mi}{67}\RG{12.819}{0.002}{10.9091}{rm}{41}\RG{20.477}{0.002}{10.9091}{rm}{61}\RG{32.404}{0.002}{10.9091}{mi}{114}%
}\hskip.1em\allowbreak{} 。当\allowbreak{}\hskip.1em \RB{26.073}{7.768}{0.738}{%
\RG{0.000}{0.002}{10.9091}{mi}{114}\RG{8.653}{0.002}{10.9091}{rm}{61}\RG{20.555}{0.002}{10.9091}{rm}{48}%
}\hskip.1em\allowbreak{} 时，\allowbreak{}\hskip.1em \RB{32.097}{8.313}{0.738}{%
\RG{0.000}{0.002}{10.9091}{mi}{65}\RG{11.611}{0.002}{10.9091}{rm}{61}\RG{23.524}{0.002}{10.9091}{mi}{79}%
}\hskip.1em\allowbreak{} ；取空列、空行组成的矩阵，结论依空基的通常约定仍成立。\allowbreak{}\hskip.1em \RB{8.366}{7.993}{0.502}{%
\RG{0.000}{-0.002}{10.9091}{am}{3}%
}\hskip.1em\allowbreak{}

方法提炼：从像空间到坐标矩阵。满秩分解不是额外的矩阵技巧，而是选择列空间的一组基并记录各列坐标。行空间给出完全对偶的构造。

题目来源：谢启鸿《复旦大学高等代数习题集（2023版）》，第109页。

% Original page 4.
\section*{2021A12}
设\allowbreak{}\hskip.1em \RB{55.676}{8.686}{3.223}{%
\RG{0.000}{0.004}{10.9091}{mi}{65}\RG{11.212}{0.004}{10.9091}{sy}{50}\RG{21.515}{0.004}{10.9091}{mi}{77}\RG{32.099}{-1.633}{7.9701}{mi}{110}\RG{37.735}{0.004}{10.9091}{rm}{40}\RG{41.977}{0.004}{10.9091}{mi}{75}\RG{52.023}{0.004}{10.9091}{rm}{41}%
}\hskip.1em\allowbreak{} 、rank\allowbreak{}\hskip.1em \RB{36.469}{8.686}{3.223}{%
\RG{0.000}{0.004}{10.9091}{rm}{40}\RG{4.243}{0.004}{10.9091}{mi}{65}\RG{12.424}{0.004}{10.9091}{rm}{41}\RG{19.679}{0.004}{10.9091}{rm}{61}\RG{31.212}{0.004}{10.9091}{mi}{114}%
}\hskip.1em\allowbreak{} ，\allowbreak{}\hskip.1em \RB{29.554}{8.075}{0.932}{%
\RG{0.000}{0.004}{10.9091}{mi}{107}\RG{9.053}{0.004}{10.9091}{sy}{50}\RG{19.356}{0.004}{10.9091}{mi}{75}%
}\hskip.1em\allowbreak{} 。证明：\allowbreak{}\hskip.1em \RB{41.837}{9.762}{0.616}{%
\RG{0.000}{0.004}{10.9091}{mi}{65}\RG{8.182}{3.962}{7.9701}{rm}{50}\RG{15.945}{0.004}{10.9091}{rm}{61}\RG{27.460}{0.004}{10.9091}{mi}{107}\RG{33.471}{0.004}{10.9091}{mi}{65}%
}\hskip.1em\allowbreak{} 的充分必要条件是：对\allowbreak{}\hskip.1em \RB{8.366}{8.315}{0.500}{%
\RG{0.000}{0.004}{10.9091}{mi}{65}%
}\hskip.1em\allowbreak{} 的任一满秩分解\allowbreak{}\hskip.1em \RB{40.333}{8.311}{0.740}{%
\RG{0.000}{-0.000}{10.9091}{mi}{65}\RG{11.212}{-0.000}{10.9091}{rm}{61}\RG{22.727}{-0.000}{10.9091}{mi}{66}\RG{31.542}{-0.000}{10.9091}{mi}{67}%
}\hskip.1em\allowbreak{} ，均有

\[
\RB{49.912}{8.187}{2.220}{%
\RG{0.000}{-0.004}{10.9091}{mi}{67}\RG{8.564}{-0.004}{10.9091}{mi}{66}\RG{20.429}{-0.004}{10.9091}{rm}{61}\RG{31.945}{-0.004}{10.9091}{mi}{107}\RG{37.956}{-0.004}{10.9091}{mi}{73}\RG{47.317}{-0.004}{10.9091}{mi}{58}\RG{42.763}{-1.640}{7.9701}{mi}{114}%
}
\]
解：先设\allowbreak{}\hskip.1em \RB{24.841}{7.764}{1.997}{%
\RG{0.000}{-0.002}{10.9091}{mi}{114}\RG{8.255}{-0.002}{10.9091}{sy}{21}\RG{19.770}{-0.002}{10.9091}{rm}{49}%
}\hskip.1em\allowbreak{} 。由\allowbreak{}\hskip.1em \RB{40.333}{8.309}{0.742}{%
\RG{0.000}{-0.002}{10.9091}{mi}{65}\RG{11.212}{-0.002}{10.9091}{rm}{61}\RG{22.727}{-0.002}{10.9091}{mi}{66}\RG{31.542}{-0.002}{10.9091}{mi}{67}%
}\hskip.1em\allowbreak{} 为满秩分解可知

\[
\RB{280.161}{9.289}{3.231}{%
\RG{0.000}{-0.004}{10.9091}{mi}{66}\RG{11.852}{-0.004}{10.9091}{sy}{50}\RG{22.155}{-0.004}{10.9091}{mi}{77}\RG{49.018}{-0.004}{10.9091}{rm}{40}\RG{53.260}{-0.004}{10.9091}{mi}{75}\RG{63.307}{-0.004}{10.9091}{rm}{41}\RG{67.549}{-0.004}{10.9091}{mi}{59}\RG{83.291}{-0.004}{10.9091}{mi}{67}\RG{94.914}{-0.004}{10.9091}{sy}{50}\RG{105.217}{-0.004}{10.9091}{mi}{77}\RG{132.080}{-0.004}{10.9091}{rm}{40}\RG{136.322}{-0.004}{10.9091}{mi}{75}\RG{146.368}{-0.004}{10.9091}{rm}{41}\RG{150.610}{-0.004}{10.9091}{mi}{59}\RT{166.368}{-0.004}{10.9091}{r}\RT{170.645}{-0.004}{10.9091}{a}\RT{176.099}{-0.004}{10.9091}{n}\RT{182.165}{-0.004}{10.9091}{k}\RG{187.925}{-0.004}{10.9091}{rm}{40}\RG{192.167}{-0.004}{10.9091}{mi}{66}\RG{200.989}{-0.004}{10.9091}{rm}{41}\RG{208.243}{-0.004}{10.9091}{rm}{61}\RT{219.777}{-0.004}{10.9091}{r}\RT{224.053}{-0.004}{10.9091}{a}\RT{229.508}{-0.004}{10.9091}{n}\RT{235.573}{-0.004}{10.9091}{k}\RG{241.334}{-0.004}{10.9091}{rm}{40}\RG{245.576}{-0.004}{10.9091}{mi}{67}\RG{254.153}{-0.004}{10.9091}{rm}{41}\RG{261.408}{-0.004}{10.9091}{rm}{61}\RG{272.941}{-0.004}{10.9091}{mi}{114}\RG{277.567}{-0.004}{10.9091}{mi}{58}\RG{32.738}{-1.640}{7.9701}{mi}{110}\RG{37.876}{-1.640}{7.9701}{sy}{2}\RG{44.463}{-1.640}{7.9701}{mi}{114}\RG{115.800}{-1.640}{7.9701}{mi}{114}\RG{119.857}{-1.640}{7.9701}{sy}{2}\RG{126.443}{-1.640}{7.9701}{mi}{110}%
}
\]
因此\allowbreak{}\hskip.1em \RB{8.747}{7.954}{0.500}{%
\RG{0.000}{0.003}{10.9091}{mi}{66}%
}\hskip.1em\allowbreak{} 有左逆\allowbreak{}\hskip.1em \RB{65.562}{8.685}{3.224}{%
\RG{0.000}{0.003}{10.9091}{mi}{76}\RG{10.455}{0.003}{10.9091}{sy}{50}\RG{20.758}{0.003}{10.9091}{mi}{77}\RG{31.341}{-1.633}{7.9701}{mi}{114}\RG{35.398}{-1.633}{7.9701}{sy}{2}\RG{41.984}{-1.633}{7.9701}{mi}{110}\RG{47.621}{0.003}{10.9091}{rm}{40}\RG{51.863}{0.003}{10.9091}{mi}{75}\RG{61.910}{0.003}{10.9091}{rm}{41}%
}\hskip.1em\allowbreak{} 、\allowbreak{}\hskip.1em \RB{8.791}{8.194}{0.737}{%
\RG{0.000}{0.003}{10.9091}{mi}{67}%
}\hskip.1em\allowbreak{} 有右逆\allowbreak{}\hskip.1em \RB{66.505}{8.685}{3.224}{%
\RG{0.000}{0.003}{10.9091}{mi}{82}\RG{11.398}{0.003}{10.9091}{sy}{50}\RG{21.701}{0.003}{10.9091}{mi}{77}\RG{32.284}{-1.633}{7.9701}{mi}{110}\RG{37.422}{-1.633}{7.9701}{sy}{2}\RG{44.009}{-1.633}{7.9701}{mi}{114}\RG{48.564}{0.003}{10.9091}{rm}{40}\RG{52.806}{0.003}{10.9091}{mi}{75}\RG{62.853}{0.003}{10.9091}{rm}{41}%
}\hskip.1em\allowbreak{} ，满足

\[
\RB{110.244}{8.192}{2.604}{%
\RG{0.000}{0.001}{10.9091}{mi}{76}\RG{7.418}{0.001}{10.9091}{mi}{66}\RG{19.277}{0.001}{10.9091}{rm}{61}\RG{30.792}{0.001}{10.9091}{mi}{73}\RG{40.142}{0.001}{10.9091}{mi}{59}\RG{66.793}{0.001}{10.9091}{mi}{67}\RG{75.368}{0.001}{10.9091}{mi}{82}\RG{86.785}{0.001}{10.9091}{rm}{61}\RG{98.300}{0.001}{10.9091}{mi}{73}\RG{107.650}{0.001}{10.9091}{mi}{58}\RG{35.587}{-1.635}{7.9701}{mi}{114}\RG{103.095}{-1.635}{7.9701}{mi}{114}%
}
\]
若\allowbreak{}\hskip.1em \RB{41.836}{9.757}{0.621}{%
\RG{0.000}{-0.001}{10.9091}{mi}{65}\RG{8.182}{3.957}{7.9701}{rm}{50}\RG{15.944}{-0.001}{10.9091}{rm}{61}\RG{27.459}{-0.001}{10.9091}{mi}{107}\RG{33.470}{-0.001}{10.9091}{mi}{65}%
}\hskip.1em\allowbreak{} ，代入\allowbreak{}\hskip.1em \RB{40.333}{8.310}{0.741}{%
\RG{0.000}{-0.001}{10.9091}{mi}{65}\RG{11.212}{-0.001}{10.9091}{rm}{61}\RG{22.727}{-0.001}{10.9091}{mi}{66}\RG{31.542}{-0.001}{10.9091}{mi}{67}%
}\hskip.1em\allowbreak{} 得

\[
\RB{205.050}{8.678}{3.231}{%
\RG{0.000}{-0.004}{10.9091}{mi}{66}\RG{8.822}{-0.004}{10.9091}{rm}{40}\RG{13.064}{-0.004}{10.9091}{mi}{67}\RG{21.628}{-0.004}{10.9091}{mi}{66}\RG{30.463}{-0.004}{10.9091}{rm}{41}\RG{34.706}{-0.004}{10.9091}{mi}{67}\RG{46.313}{-0.004}{10.9091}{rm}{61}\RG{57.828}{-0.004}{10.9091}{mi}{107}\RG{63.839}{-0.004}{10.9091}{mi}{66}\RG{72.654}{-0.004}{10.9091}{mi}{67}\RG{80.618}{-0.004}{10.9091}{mi}{59}\RG{107.280}{-0.004}{10.9091}{mi}{66}\RG{116.134}{-0.004}{10.9091}{rm}{40}\RG{120.376}{-0.004}{10.9091}{mi}{67}\RG{128.940}{-0.004}{10.9091}{mi}{66}\RG{140.200}{-0.004}{10.9091}{sy}{0}\RG{151.109}{-0.004}{10.9091}{mi}{107}\RG{157.120}{-0.004}{10.9091}{mi}{73}\RG{166.482}{-0.004}{10.9091}{rm}{41}\RG{170.724}{-0.004}{10.9091}{mi}{67}\RG{182.333}{-0.004}{10.9091}{rm}{61}\RG{193.848}{-0.004}{10.9091}{mi}{79}\RG{202.455}{-0.004}{10.9091}{mi}{58}\RG{161.927}{-1.640}{7.9701}{mi}{114}%
}
\]
等式左乘\allowbreak{}\hskip.1em \RB{7.515}{7.955}{0.500}{%
\RG{0.000}{0.004}{10.9091}{mi}{76}%
}\hskip.1em\allowbreak{} 、右乘\allowbreak{}\hskip.1em \RB{8.736}{7.955}{0.736}{%
\RG{0.000}{0.004}{10.9091}{mi}{82}%
}\hskip.1em\allowbreak{} ，得到

\[
\RB{174.300}{8.683}{3.226}{%
\RG{0.000}{0.001}{10.9091}{mi}{67}\RG{8.564}{0.001}{10.9091}{mi}{66}\RG{19.823}{0.001}{10.9091}{sy}{0}\RG{30.732}{0.001}{10.9091}{mi}{107}\RG{36.743}{0.001}{10.9091}{mi}{73}\RG{49.136}{0.001}{10.9091}{rm}{61}\RG{60.652}{0.001}{10.9091}{mi}{76}\RG{68.070}{0.001}{10.9091}{mi}{66}\RG{76.898}{0.001}{10.9091}{rm}{40}\RG{81.140}{0.001}{10.9091}{mi}{67}\RG{89.703}{0.001}{10.9091}{mi}{66}\RG{100.964}{0.001}{10.9091}{sy}{0}\RG{111.873}{0.001}{10.9091}{mi}{107}\RG{117.884}{0.001}{10.9091}{mi}{73}\RG{127.246}{0.001}{10.9091}{rm}{41}\RG{131.488}{0.001}{10.9091}{mi}{67}\RG{140.052}{0.001}{10.9091}{mi}{82}\RG{151.463}{0.001}{10.9091}{rm}{61}\RG{162.978}{0.001}{10.9091}{mi}{79}\RG{171.586}{0.001}{10.9091}{mi}{59}\RG{41.551}{-1.635}{7.9701}{mi}{114}\RG{122.691}{-1.635}{7.9701}{mi}{114}%
}
\]
从而\allowbreak{}\hskip.1em \RB{46.946}{8.190}{2.218}{%
\RG{0.000}{-0.001}{10.9091}{mi}{67}\RG{8.564}{-0.001}{10.9091}{mi}{66}\RG{20.429}{-0.001}{10.9091}{rm}{61}\RG{31.944}{-0.001}{10.9091}{mi}{107}\RG{37.955}{-0.001}{10.9091}{mi}{73}\RG{42.763}{-1.638}{7.9701}{mi}{114}%
}\hskip.1em\allowbreak{} 。以上推理没有对满秩分解施加特殊选择，所以结论对每一种满秩分解均成立。

反之，若取\allowbreak{}\hskip.1em \RB{8.366}{8.316}{0.500}{%
\RG{0.000}{0.005}{10.9091}{mi}{65}%
}\hskip.1em\allowbreak{} 的任一个满秩分解\allowbreak{}\hskip.1em \RB{40.333}{8.316}{0.735}{%
\RG{0.000}{0.005}{10.9091}{mi}{65}\RG{11.212}{0.005}{10.9091}{rm}{61}\RG{22.727}{0.005}{10.9091}{mi}{66}\RG{31.542}{0.005}{10.9091}{mi}{67}%
}\hskip.1em\allowbreak{} ，并有\allowbreak{}\hskip.1em \RB{46.945}{8.196}{2.211}{%
\RG{0.000}{0.005}{10.9091}{mi}{67}\RG{8.564}{0.005}{10.9091}{mi}{66}\RG{20.429}{0.005}{10.9091}{rm}{61}\RG{31.944}{0.005}{10.9091}{mi}{107}\RG{37.955}{0.005}{10.9091}{mi}{73}\RG{42.762}{-1.631}{7.9701}{mi}{114}%
}\hskip.1em\allowbreak{} ，则

\[
\RB{224.194}{10.307}{3.224}{%
\RG{8.182}{4.507}{7.9701}{rm}{50}\RG{0.000}{0.003}{10.9091}{mi}{65}\RG{15.945}{0.003}{10.9091}{rm}{61}\RG{27.443}{0.003}{10.9091}{rm}{40}\RG{31.702}{0.003}{10.9091}{mi}{66}\RG{40.517}{0.003}{10.9091}{mi}{67}\RG{49.102}{0.003}{10.9091}{rm}{41}\RG{53.335}{0.003}{10.9091}{rm}{40}\RG{57.586}{0.003}{10.9091}{mi}{66}\RG{66.401}{0.003}{10.9091}{mi}{67}\RG{74.987}{0.003}{10.9091}{rm}{41}\RG{82.241}{0.003}{10.9091}{rm}{61}\RG{93.774}{0.003}{10.9091}{mi}{66}\RG{102.596}{0.003}{10.9091}{rm}{40}\RG{106.839}{0.003}{10.9091}{mi}{67}\RG{115.403}{0.003}{10.9091}{mi}{66}\RG{124.238}{0.003}{10.9091}{rm}{41}\RG{128.481}{0.003}{10.9091}{mi}{67}\RG{140.088}{0.003}{10.9091}{rm}{61}\RG{151.603}{0.003}{10.9091}{mi}{66}\RG{160.425}{0.003}{10.9091}{rm}{40}\RG{164.668}{0.003}{10.9091}{mi}{107}\RG{170.679}{0.003}{10.9091}{mi}{73}\RG{180.040}{0.003}{10.9091}{rm}{41}\RG{184.283}{0.003}{10.9091}{mi}{67}\RG{195.891}{0.003}{10.9091}{rm}{61}\RG{207.407}{0.003}{10.9091}{mi}{107}\RG{213.418}{0.003}{10.9091}{mi}{65}\RG{221.600}{0.003}{10.9091}{mi}{58}\RG{175.486}{-1.634}{7.9701}{mi}{114}%
}
\]
当\allowbreak{}\hskip.1em \RB{25.272}{7.765}{0.740}{%
\RG{0.000}{-0.000}{10.9091}{mi}{114}\RG{8.255}{-0.000}{10.9091}{rm}{61}\RG{19.753}{-0.000}{10.9091}{rm}{48}%
}\hskip.1em\allowbreak{} 时\allowbreak{}\hskip.1em \RB{31.300}{8.311}{0.740}{%
\RG{0.000}{-0.000}{10.9091}{mi}{65}\RG{11.212}{-0.000}{10.9091}{rm}{61}\RG{22.727}{-0.000}{10.9091}{mi}{79}%
}\hskip.1em\allowbreak{} ，两边的条件按零维单位矩阵的约定同时成立。\allowbreak{}\hskip.1em \RB{8.366}{7.995}{0.500}{%
\RG{0.000}{-0.000}{10.9091}{am}{3}%
}\hskip.1em\allowbreak{}

方法提炼：借满秩因子消去。\allowbreak{}\hskip.1em \RB{8.747}{7.950}{0.501}{%
\RG{0.000}{-0.001}{10.9091}{mi}{66}%
}\hskip.1em\allowbreak{} 列满秩、\allowbreak{}\hskip.1em \RB{8.791}{8.190}{0.741}{%
\RG{0.000}{-0.001}{10.9091}{mi}{67}%
}\hskip.1em\allowbreak{} 行满秩分别保证左逆、右逆存在；于是矩阵方程\allowbreak{}\hskip.1em \RB{50.411}{8.196}{0.735}{%
\RG{0.000}{0.005}{10.9091}{mi}{66}\RG{8.815}{0.005}{10.9091}{mi}{88}\RG{18.698}{0.005}{10.9091}{mi}{67}\RG{30.323}{0.005}{10.9091}{rm}{61}\RG{41.838}{0.005}{10.9091}{mi}{79}%
}\hskip.1em\allowbreak{} 可等价地化为\allowbreak{}\hskip.1em \RB{33.012}{8.196}{0.735}{%
\RG{0.000}{0.005}{10.9091}{mi}{88}\RG{12.924}{0.005}{10.9091}{rm}{61}\RG{24.439}{0.005}{10.9091}{mi}{79}%
}\hskip.1em\allowbreak{} 。

题目来源：谢启鸿《复旦大学高等代数习题集（2023版）》，第109页。

% Original page 5.
\section*{2021A13}
设\allowbreak{}\hskip.1em \RB{20.595}{7.955}{2.601}{%
\RG{0.000}{0.004}{10.9091}{mi}{86}\RG{6.960}{0.004}{10.9091}{mi}{59}\RG{11.804}{0.004}{10.9091}{mi}{85}%
}\hskip.1em\allowbreak{} 分别为数域\allowbreak{}\hskip.1em \RB{10.198}{7.955}{0.500}{%
\RG{0.000}{0.004}{10.9091}{mi}{75}%
}\hskip.1em\allowbreak{} 上的有限维线性空间，\allowbreak{}\hskip.1em \RB{50.775}{7.955}{2.874}{%
\RG{0.000}{0.004}{10.9091}{mi}{39}\RG{10.166}{0.004}{10.9091}{rm}{58}\RG{16.227}{0.004}{10.9091}{mi}{86}\RG{28.045}{0.004}{10.9091}{sy}{33}\RG{41.984}{0.004}{10.9091}{mi}{85}%
}\hskip.1em\allowbreak{} 是线性映射。

(1)证明存在\allowbreak{}\hskip.1em \RB{10.198}{7.951}{0.500}{%
\RG{0.000}{-0.000}{10.9091}{mi}{75}%
}\hskip.1em\allowbreak{} 上的线性空间\allowbreak{}\hskip.1em \RB{11.933}{7.951}{0.740}{%
\RG{0.000}{-0.000}{10.9091}{mi}{87}%
}\hskip.1em\allowbreak{} 、满线性映射\allowbreak{}\hskip.1em \RB{53.880}{8.104}{2.736}{%
\RG{0.000}{-0.000}{10.9091}{mi}{24}\RG{8.761}{-0.000}{10.9091}{rm}{58}\RG{15.278}{-0.000}{10.9091}{mi}{86}\RG{27.552}{-0.000}{10.9091}{sy}{33}\RG{41.947}{-0.000}{10.9091}{mi}{87}%
}\hskip.1em\allowbreak{} 以及单线性映射\allowbreak{}\hskip.1em \RB{52.252}{7.951}{0.740}{%
\RG{0.000}{-0.000}{10.9091}{mi}{105}\RG{7.244}{-0.000}{10.9091}{rm}{58}\RG{13.761}{-0.000}{10.9091}{mi}{87}\RG{29.065}{-0.000}{10.9091}{sy}{33}\RG{43.461}{-0.000}{10.9091}{mi}{85}%
}\hskip.1em\allowbreak{} ，使得\allowbreak{}\hskip.1em \RB{41.109}{8.100}{2.882}{%
\RG{0.000}{-0.004}{10.9091}{mi}{39}\RG{10.166}{-0.004}{10.9091}{rm}{61}\RG{21.681}{-0.004}{10.9091}{mi}{105}\RG{27.864}{-0.004}{10.9091}{sy}{14}\RG{35.743}{-0.004}{10.9091}{mi}{24}%
}\hskip.1em\allowbreak{} 。

(2)若另有\allowbreak{}\hskip.1em \RB{14.366}{7.954}{2.133}{%
\RG{0.000}{0.003}{10.9091}{mi}{87}\RG{10.303}{-1.633}{7.9701}{rm}{49}%
}\hskip.1em\allowbreak{} 、满线性映射\allowbreak{}\hskip.1em \RB{58.719}{8.107}{2.733}{%
\RG{0.000}{0.003}{10.9091}{mi}{24}\RG{4.773}{-1.633}{7.9701}{rm}{49}\RG{12.535}{0.003}{10.9091}{rm}{58}\RG{18.596}{0.003}{10.9091}{mi}{86}\RG{30.414}{0.003}{10.9091}{sy}{33}\RG{44.353}{0.003}{10.9091}{mi}{87}\RG{54.656}{-1.633}{7.9701}{rm}{49}%
}\hskip.1em\allowbreak{} 和单线性映射\allowbreak{}\hskip.1em \RB{58.378}{7.954}{2.133}{%
\RG{0.000}{0.003}{10.9091}{mi}{105}\RG{3.758}{-1.633}{7.9701}{rm}{49}\RG{11.521}{0.003}{10.9091}{rm}{58}\RG{17.582}{0.003}{10.9091}{mi}{87}\RG{27.885}{-1.633}{7.9701}{rm}{49}\RG{35.647}{0.003}{10.9091}{sy}{33}\RG{49.587}{0.003}{10.9091}{mi}{85}%
}\hskip.1em\allowbreak{} ，满足\allowbreak{}\hskip.1em \RB{48.563}{8.107}{2.875}{%
\RG{0.000}{0.003}{10.9091}{mi}{39}\RG{10.167}{0.003}{10.9091}{rm}{61}\RG{21.682}{0.003}{10.9091}{mi}{105}\RG{25.440}{-1.633}{7.9701}{rm}{49}\RG{32.222}{0.003}{10.9091}{sy}{14}\RG{39.727}{0.003}{10.9091}{mi}{24}\RG{44.500}{-1.633}{7.9701}{rm}{49}%
}\hskip.1em\allowbreak{} ，证明存在同构\allowbreak{}\hskip.1em \RB{58.053}{7.950}{2.857}{%
\RG{0.000}{-0.001}{10.9091}{mi}{17}\RG{8.838}{-0.001}{10.9091}{rm}{58}\RG{14.899}{-0.001}{10.9091}{mi}{87}\RG{29.747}{-0.001}{10.9091}{sy}{33}\RG{43.687}{-0.001}{10.9091}{mi}{87}\RG{53.990}{-1.637}{7.9701}{rm}{49}%
}\hskip.1em\allowbreak{} ，使得

\[
\RB{117.588}{8.100}{2.860}{%
\RG{0.000}{-0.004}{10.9091}{mi}{24}\RG{12.535}{-0.004}{10.9091}{rm}{61}\RG{24.050}{-0.004}{10.9091}{mi}{17}\RG{32.282}{-0.004}{10.9091}{sy}{14}\RG{40.161}{-0.004}{10.9091}{mi}{24}\RG{45.430}{-0.004}{10.9091}{mi}{59}\RG{72.082}{-0.004}{10.9091}{mi}{105}\RG{78.892}{-0.004}{10.9091}{rm}{61}\RG{90.407}{-0.004}{10.9091}{mi}{105}\RG{101.322}{-0.004}{10.9091}{sy}{14}\RG{109.201}{-0.004}{10.9091}{mi}{17}\RG{114.994}{-0.004}{10.9091}{mi}{58}\RG{4.772}{-1.641}{7.9701}{rm}{49}\RG{94.165}{-1.641}{7.9701}{rm}{49}%
}
\]
解：(1)取\allowbreak{}\hskip.1em \RB{23.214}{7.949}{0.742}{%
\RG{0.000}{-0.002}{10.9091}{mi}{87}\RG{14.848}{-0.002}{10.9091}{rm}{61}%
}\hskip.1em\allowbreak{} Im\allowbreak{}\hskip.1em \RB{30.473}{7.949}{2.880}{%
\RG{0.000}{-0.002}{10.9091}{mi}{39}\RG{10.167}{-0.002}{10.9091}{sy}{18}\RG{21.682}{-0.002}{10.9091}{mi}{85}%
}\hskip.1em\allowbreak{} ，定义

\[
\RB{263.736}{8.687}{3.222}{%
\RG{0.000}{0.005}{10.9091}{mi}{24}\RG{8.304}{0.005}{10.9091}{rm}{58}\RG{14.365}{0.005}{10.9091}{mi}{86}\RG{26.183}{0.005}{10.9091}{sy}{0}\RG{32.849}{0.005}{10.9091}{sy}{33}\RG{46.789}{0.005}{10.9091}{mi}{87}\RG{56.793}{0.005}{10.9091}{mi}{59}\RG{72.524}{0.005}{10.9091}{mi}{118}\RG{81.257}{0.005}{10.9091}{sy}{55}\RG{81.257}{0.005}{10.9091}{sy}{0}\RG{87.922}{0.005}{10.9091}{sy}{33}\RG{101.863}{0.005}{10.9091}{mi}{39}\RG{109.000}{0.005}{10.9091}{rm}{40}\RG{113.242}{0.005}{10.9091}{mi}{118}\RG{118.922}{0.005}{10.9091}{rm}{41}\RG{123.164}{0.005}{10.9091}{mi}{59}\RG{149.815}{0.005}{10.9091}{mi}{105}\RG{156.619}{0.005}{10.9091}{rm}{58}\RG{162.680}{0.005}{10.9091}{mi}{87}\RG{177.528}{0.005}{10.9091}{sy}{0}\RG{184.194}{0.005}{10.9091}{sy}{33}\RG{198.135}{0.005}{10.9091}{mi}{85}\RG{205.553}{0.005}{10.9091}{mi}{59}\RG{221.295}{0.005}{10.9091}{mi}{119}\RG{232.452}{0.005}{10.9091}{sy}{55}\RG{232.452}{0.005}{10.9091}{sy}{0}\RG{239.117}{0.005}{10.9091}{sy}{33}\RG{253.058}{0.005}{10.9091}{mi}{119}\RG{261.142}{0.005}{10.9091}{mi}{58}%
}
\]
由像空间的定义，\allowbreak{}\hskip.1em \RB{5.365}{8.107}{2.733}{%
\RG{0.000}{0.003}{10.9091}{mi}{24}%
}\hskip.1em\allowbreak{} 是满射；\allowbreak{}\hskip.1em \RB{3.696}{7.714}{0.617}{%
\RG{0.000}{0.003}{10.9091}{mi}{105}%
}\hskip.1em\allowbreak{} 为包含映射，因而是单射。对任意\allowbreak{}\hskip.1em \RB{27.902}{7.954}{0.933}{%
\RG{0.000}{0.003}{10.9091}{mi}{118}\RG{8.710}{0.003}{10.9091}{sy}{50}\RG{19.013}{0.003}{10.9091}{mi}{86}%
}\hskip.1em\allowbreak{} ，

\[
\RB{128.636}{8.682}{3.227}{%
\RG{0.000}{-0.000}{10.9091}{rm}{40}\RG{4.242}{-0.000}{10.9091}{mi}{105}\RG{10.425}{-0.000}{10.9091}{sy}{14}\RG{18.303}{-0.000}{10.9091}{mi}{24}\RG{23.578}{-0.000}{10.9091}{rm}{41}\RG{27.811}{-0.000}{10.9091}{rm}{40}\RG{32.063}{-0.000}{10.9091}{mi}{118}\RG{37.742}{-0.000}{10.9091}{rm}{41}\RG{44.997}{-0.000}{10.9091}{rm}{61}\RG{56.530}{-0.000}{10.9091}{mi}{105}\RG{60.288}{-0.000}{10.9091}{rm}{40}\RG{64.532}{-0.000}{10.9091}{mi}{39}\RG{71.668}{-0.000}{10.9091}{rm}{40}\RG{75.911}{-0.000}{10.9091}{mi}{118}\RG{81.590}{-0.000}{10.9091}{rm}{41}\RG{85.822}{-0.000}{10.9091}{rm}{41}\RG{93.077}{-0.000}{10.9091}{rm}{61}\RG{104.620}{-0.000}{10.9091}{mi}{39}\RG{111.757}{-0.000}{10.9091}{rm}{40}\RG{115.999}{-0.000}{10.9091}{mi}{118}\RG{121.678}{-0.000}{10.9091}{rm}{41}\RG{125.921}{-0.000}{10.9091}{mi}{59}%
}
\]
故\allowbreak{}\hskip.1em \RB{41.109}{8.102}{2.880}{%
\RG{0.000}{-0.002}{10.9091}{mi}{39}\RG{10.166}{-0.002}{10.9091}{rm}{61}\RG{21.681}{-0.002}{10.9091}{mi}{105}\RG{27.864}{-0.002}{10.9091}{sy}{14}\RG{35.743}{-0.002}{10.9091}{mi}{24}%
}\hskip.1em\allowbreak{} 。(2)使用(1)中的标准分解：由于\allowbreak{}\hskip.1em \RB{7.821}{7.715}{2.132}{%
\RG{0.000}{0.004}{10.9091}{mi}{105}\RG{3.758}{-1.632}{7.9701}{rm}{49}%
}\hskip.1em\allowbreak{} 单射且\allowbreak{}\hskip.1em \RB{39.009}{8.108}{2.874}{%
\RG{0.000}{0.004}{10.9091}{mi}{39}\RG{10.167}{0.004}{10.9091}{rm}{61}\RG{21.682}{0.004}{10.9091}{mi}{105}\RG{25.440}{-1.632}{7.9701}{rm}{49}\RG{30.173}{0.004}{10.9091}{mi}{24}\RG{34.946}{-1.632}{7.9701}{rm}{49}%
}\hskip.1em\allowbreak{} ，有

\[
\RB{103.681}{9.295}{2.876}{%
\RT{0.000}{0.002}{10.9091}{k}\RT{5.760}{0.002}{10.9091}{e}\RT{10.604}{0.002}{10.9091}{r}\RG{16.698}{0.002}{10.9091}{mi}{24}\RG{29.234}{0.002}{10.9091}{rm}{61}\RT{40.749}{0.002}{10.9091}{k}\RT{46.509}{0.002}{10.9091}{e}\RT{51.353}{0.002}{10.9091}{r}\RG{57.447}{0.002}{10.9091}{mi}{39}\RG{67.614}{0.002}{10.9091}{rm}{61}\RT{79.130}{0.002}{10.9091}{k}\RT{84.890}{0.002}{10.9091}{e}\RT{89.733}{0.002}{10.9091}{r}\RG{95.828}{0.002}{10.9091}{mi}{24}\RG{101.086}{0.002}{10.9091}{mi}{58}\RG{21.471}{-1.635}{7.9701}{rm}{49}%
}
\]
对\allowbreak{}\hskip.1em \RB{44.651}{7.950}{0.937}{%
\RG{0.000}{-0.001}{10.9091}{mi}{119}\RG{11.134}{-0.001}{10.9091}{sy}{50}\RG{21.437}{-0.001}{10.9091}{mi}{87}\RG{36.285}{-0.001}{10.9091}{rm}{61}%
}\hskip.1em\allowbreak{} Im\allowbreak{}\hskip.1em \RB{7.253}{5.321}{2.879}{%
\RG{0.000}{-0.001}{10.9091}{mi}{39}%
}\hskip.1em\allowbreak{} ，任选\allowbreak{}\hskip.1em \RB{27.902}{7.950}{0.937}{%
\RG{0.000}{-0.001}{10.9091}{mi}{118}\RG{8.710}{-0.001}{10.9091}{sy}{50}\RG{19.013}{-0.001}{10.9091}{mi}{86}%
}\hskip.1em\allowbreak{} 使得\allowbreak{}\hskip.1em \RB{43.885}{8.681}{3.228}{%
\RG{0.000}{-0.001}{10.9091}{mi}{39}\RG{7.136}{-0.001}{10.9091}{rm}{40}\RG{11.379}{-0.001}{10.9091}{mi}{118}\RG{17.058}{-0.001}{10.9091}{rm}{41}\RG{24.313}{-0.001}{10.9091}{rm}{61}\RG{35.846}{-0.001}{10.9091}{mi}{119}%
}\hskip.1em\allowbreak{} ，定义

\[
\RB{63.206}{8.679}{3.230}{%
\RG{0.000}{-0.003}{10.9091}{mi}{17}\RG{5.808}{-0.003}{10.9091}{rm}{40}\RG{10.051}{-0.003}{10.9091}{mi}{119}\RG{18.154}{-0.003}{10.9091}{rm}{41}\RG{25.409}{-0.003}{10.9091}{rm}{61}\RG{36.942}{-0.003}{10.9091}{mi}{24}\RG{46.447}{-0.003}{10.9091}{rm}{40}\RG{50.690}{-0.003}{10.9091}{mi}{118}\RG{56.369}{-0.003}{10.9091}{rm}{41}\RG{60.611}{-0.003}{10.9091}{mi}{58}\RG{41.715}{-1.640}{7.9701}{rm}{49}%
}
\]
若同一个\allowbreak{}\hskip.1em \RB{8.038}{5.326}{0.616}{%
\RG{0.000}{0.004}{10.9091}{mi}{119}%
}\hskip.1em\allowbreak{} 还可写为\allowbreak{}\hskip.1em \RB{23.507}{8.918}{3.223}{%
\RG{0.000}{0.004}{10.9091}{mi}{39}\RG{7.137}{0.004}{10.9091}{rm}{40}\RG{11.379}{0.004}{10.9091}{mi}{118}\RG{17.058}{3.963}{7.9701}{sy}{48}\RG{19.854}{0.004}{10.9091}{rm}{41}%
}\hskip.1em\allowbreak{} ，则

\[
\RB{259.331}{9.461}{3.225}{%
\RG{36.070}{4.506}{7.9701}{sy}{48}\RG{118.703}{4.506}{7.9701}{sy}{48}\RG{249.578}{4.506}{7.9701}{sy}{48}\RG{0.000}{0.002}{10.9091}{mi}{39}\RG{7.136}{0.002}{10.9091}{rm}{40}\RG{11.379}{0.002}{10.9091}{mi}{118}\RG{19.482}{0.002}{10.9091}{sy}{0}\RG{30.391}{0.002}{10.9091}{mi}{118}\RG{38.866}{0.002}{10.9091}{rm}{41}\RG{46.121}{0.002}{10.9091}{rm}{61}\RG{57.630}{0.002}{10.9091}{rm}{48}\RG{69.750}{0.002}{10.9091}{rm}{61}\RG{76.439}{0.002}{10.9091}{sy}{41}\RG{94.011}{0.002}{10.9091}{mi}{118}\RG{102.114}{0.002}{10.9091}{sy}{0}\RG{113.023}{0.002}{10.9091}{mi}{118}\RG{124.528}{0.002}{10.9091}{sy}{50}\RT{134.831}{0.002}{10.9091}{k}\RT{140.591}{0.002}{10.9091}{e}\RT{145.435}{0.002}{10.9091}{r}\RG{151.529}{0.002}{10.9091}{mi}{24}\RG{167.697}{0.002}{10.9091}{rm}{61}\RG{174.364}{0.002}{10.9091}{sy}{41}\RG{191.937}{0.002}{10.9091}{mi}{24}\RG{201.442}{0.002}{10.9091}{rm}{40}\RG{205.684}{0.002}{10.9091}{mi}{118}\RG{211.364}{0.002}{10.9091}{rm}{41}\RG{218.618}{0.002}{10.9091}{rm}{61}\RG{230.152}{0.002}{10.9091}{mi}{24}\RG{239.657}{0.002}{10.9091}{rm}{40}\RG{243.899}{0.002}{10.9091}{mi}{118}\RG{252.373}{0.002}{10.9091}{rm}{41}\RG{256.616}{0.002}{10.9091}{mi}{59}\RG{156.302}{-1.635}{7.9701}{rm}{49}\RG{196.710}{-1.635}{7.9701}{rm}{49}\RG{234.924}{-1.635}{7.9701}{rm}{49}%
}
\]
故\allowbreak{}\hskip.1em \RB{5.911}{5.321}{2.857}{%
\RG{0.000}{-0.001}{10.9091}{mi}{17}%
}\hskip.1em\allowbreak{} 定义良好。由\allowbreak{}\hskip.1em \RB{7.253}{5.321}{2.879}{%
\RG{0.000}{-0.001}{10.9091}{mi}{39}%
}\hskip.1em\allowbreak{} 和\allowbreak{}\hskip.1em \RB{8.836}{8.103}{2.737}{%
\RG{0.000}{-0.001}{10.9091}{mi}{24}\RG{4.773}{-1.637}{7.9701}{rm}{49}%
}\hskip.1em\allowbreak{} 的线性性，\allowbreak{}\hskip.1em \RB{5.911}{5.321}{2.857}{%
\RG{0.000}{-0.001}{10.9091}{mi}{17}%
}\hskip.1em\allowbreak{} 也保持加法和数乘。

任取\allowbreak{}\hskip.1em \RB{40.242}{7.947}{2.141}{%
\RG{0.000}{-0.004}{10.9091}{mi}{119}\RG{7.810}{-1.641}{7.9701}{rm}{49}\RG{15.573}{-0.004}{10.9091}{sy}{50}\RG{25.876}{-0.004}{10.9091}{mi}{87}\RG{36.179}{-1.641}{7.9701}{rm}{49}%
}\hskip.1em\allowbreak{} ，\allowbreak{}\hskip.1em \RB{8.836}{8.100}{2.740}{%
\RG{0.000}{-0.004}{10.9091}{mi}{24}\RG{4.773}{-1.641}{7.9701}{rm}{49}%
}\hskip.1em\allowbreak{} 的满射性给出\allowbreak{}\hskip.1em \RB{27.902}{7.947}{0.940}{%
\RG{0.000}{-0.004}{10.9091}{mi}{118}\RG{8.710}{-0.004}{10.9091}{sy}{50}\RG{19.013}{-0.004}{10.9091}{mi}{86}%
}\hskip.1em\allowbreak{} ，使得\allowbreak{}\hskip.1em \RB{50.087}{8.678}{3.231}{%
\RG{0.000}{-0.004}{10.9091}{mi}{24}\RG{4.772}{-1.641}{7.9701}{rm}{49}\RG{9.505}{-0.004}{10.9091}{rm}{40}\RG{13.747}{-0.004}{10.9091}{mi}{118}\RG{19.426}{-0.004}{10.9091}{rm}{41}\RG{26.681}{-0.004}{10.9091}{rm}{61}\RG{38.214}{-0.004}{10.9091}{mi}{119}\RG{46.024}{-1.641}{7.9701}{rm}{49}%
}\hskip.1em\allowbreak{} ，于是

\[
\RB{65.396}{8.685}{3.224}{%
\RG{0.000}{0.003}{10.9091}{mi}{17}\RG{5.808}{0.003}{10.9091}{rm}{40}\RG{10.051}{0.003}{10.9091}{mi}{39}\RG{17.187}{0.003}{10.9091}{rm}{40}\RG{21.430}{0.003}{10.9091}{mi}{118}\RG{27.109}{0.003}{10.9091}{rm}{41}\RG{31.342}{0.003}{10.9091}{rm}{41}\RG{38.597}{0.003}{10.9091}{rm}{61}\RG{50.139}{0.003}{10.9091}{mi}{119}\RG{62.682}{0.003}{10.9091}{mi}{59}\RG{57.949}{-1.633}{7.9701}{rm}{49}%
}
\]
即\allowbreak{}\hskip.1em \RB{5.911}{5.323}{2.855}{%
\RG{0.000}{0.001}{10.9091}{mi}{17}%
}\hskip.1em\allowbreak{} 为满射。若\allowbreak{}\hskip.1em \RB{42.436}{8.683}{3.226}{%
\RG{0.000}{0.001}{10.9091}{mi}{17}\RG{5.808}{0.001}{10.9091}{rm}{40}\RG{10.051}{0.001}{10.9091}{mi}{119}\RG{18.154}{0.001}{10.9091}{rm}{41}\RG{25.409}{0.001}{10.9091}{rm}{61}\RG{36.918}{0.001}{10.9091}{rm}{48}%
}\hskip.1em\allowbreak{} 且\allowbreak{}\hskip.1em \RB{43.360}{8.683}{3.226}{%
\RG{0.000}{0.001}{10.9091}{mi}{119}\RG{11.134}{0.001}{10.9091}{rm}{61}\RG{22.649}{0.001}{10.9091}{mi}{39}\RG{29.785}{0.001}{10.9091}{rm}{40}\RG{34.028}{0.001}{10.9091}{mi}{118}\RG{39.707}{0.001}{10.9091}{rm}{41}%
}\hskip.1em\allowbreak{} ，则\allowbreak{}\hskip.1em \RB{43.709}{8.683}{3.226}{%
\RG{0.000}{0.001}{10.9091}{mi}{24}\RG{4.773}{-1.636}{7.9701}{rm}{49}\RG{9.505}{0.001}{10.9091}{rm}{40}\RG{13.748}{0.001}{10.9091}{mi}{118}\RG{19.427}{0.001}{10.9091}{rm}{41}\RG{26.682}{0.001}{10.9091}{rm}{61}\RG{38.191}{0.001}{10.9091}{rm}{48}%
}\hskip.1em\allowbreak{} ，从而\allowbreak{}\hskip.1em \RB{74.849}{8.683}{3.226}{%
\RG{0.000}{0.001}{10.9091}{mi}{119}\RG{11.134}{0.001}{10.9091}{rm}{61}\RG{22.649}{0.001}{10.9091}{mi}{105}\RG{26.407}{-1.636}{7.9701}{rm}{49}\RG{31.140}{0.001}{10.9091}{mi}{24}\RG{35.912}{-1.636}{7.9701}{rm}{49}\RG{40.645}{0.001}{10.9091}{rm}{40}\RG{44.887}{0.001}{10.9091}{mi}{118}\RG{50.567}{0.001}{10.9091}{rm}{41}\RG{57.822}{0.001}{10.9091}{rm}{61}\RG{69.331}{0.001}{10.9091}{rm}{48}%
}\hskip.1em\allowbreak{} ，即\allowbreak{}\hskip.1em \RB{5.911}{5.323}{2.855}{%
\RG{0.000}{0.001}{10.9091}{mi}{17}%
}\hskip.1em\allowbreak{} 为单射。因此\allowbreak{}\hskip.1em \RB{5.911}{5.319}{2.859}{%
\RG{0.000}{-0.003}{10.9091}{mi}{17}%
}\hskip.1em\allowbreak{} 是线性同构。

最后，对任意\allowbreak{}\hskip.1em \RB{27.902}{7.955}{0.932}{%
\RG{0.000}{0.004}{10.9091}{mi}{118}\RG{8.710}{0.004}{10.9091}{sy}{50}\RG{19.013}{0.004}{10.9091}{mi}{86}%
}\hskip.1em\allowbreak{} 及\allowbreak{}\hskip.1em \RB{69.216}{8.686}{3.223}{%
\RG{0.000}{0.004}{10.9091}{mi}{119}\RG{11.134}{0.004}{10.9091}{rm}{61}\RG{22.649}{0.004}{10.9091}{mi}{39}\RG{29.786}{0.004}{10.9091}{rm}{40}\RG{34.028}{0.004}{10.9091}{mi}{118}\RG{39.707}{0.004}{10.9091}{rm}{41}\RG{46.980}{0.004}{10.9091}{sy}{50}\RG{57.283}{0.004}{10.9091}{mi}{87}%
}\hskip.1em\allowbreak{} ，

\[
\RB{337.255}{8.683}{3.226}{%
\RG{0.000}{0.001}{10.9091}{rm}{40}\RG{4.242}{0.001}{10.9091}{mi}{17}\RG{12.475}{0.001}{10.9091}{sy}{14}\RG{20.353}{0.001}{10.9091}{mi}{24}\RG{25.628}{0.001}{10.9091}{rm}{41}\RG{29.861}{0.001}{10.9091}{rm}{40}\RG{34.113}{0.001}{10.9091}{mi}{118}\RG{39.792}{0.001}{10.9091}{rm}{41}\RG{47.047}{0.001}{10.9091}{rm}{61}\RG{58.580}{0.001}{10.9091}{mi}{17}\RG{64.389}{0.001}{10.9091}{rm}{40}\RG{68.632}{0.001}{10.9091}{mi}{39}\RG{75.768}{0.001}{10.9091}{rm}{40}\RG{80.010}{0.001}{10.9091}{mi}{118}\RG{85.690}{0.001}{10.9091}{rm}{41}\RG{89.922}{0.001}{10.9091}{rm}{41}\RG{97.177}{0.001}{10.9091}{rm}{61}\RG{108.720}{0.001}{10.9091}{mi}{24}\RG{118.225}{0.001}{10.9091}{rm}{40}\RG{122.467}{0.001}{10.9091}{mi}{118}\RG{128.147}{0.001}{10.9091}{rm}{41}\RG{132.389}{0.001}{10.9091}{mi}{59}\RG{159.056}{0.001}{10.9091}{rm}{40}\RG{163.298}{0.001}{10.9091}{mi}{105}\RG{174.213}{0.001}{10.9091}{sy}{14}\RG{182.093}{0.001}{10.9091}{mi}{17}\RG{187.901}{0.001}{10.9091}{rm}{41}\RG{192.133}{0.001}{10.9091}{rm}{40}\RG{196.386}{0.001}{10.9091}{mi}{119}\RG{204.489}{0.001}{10.9091}{rm}{41}\RG{211.743}{0.001}{10.9091}{rm}{61}\RG{223.277}{0.001}{10.9091}{mi}{105}\RG{231.768}{0.001}{10.9091}{rm}{40}\RG{236.010}{0.001}{10.9091}{mi}{24}\RG{245.515}{0.001}{10.9091}{rm}{40}\RG{249.758}{0.001}{10.9091}{mi}{118}\RG{255.437}{0.001}{10.9091}{rm}{41}\RG{259.670}{0.001}{10.9091}{rm}{41}\RG{266.924}{0.001}{10.9091}{rm}{61}\RG{278.467}{0.001}{10.9091}{mi}{39}\RG{285.604}{0.001}{10.9091}{rm}{40}\RG{289.846}{0.001}{10.9091}{mi}{118}\RG{295.525}{0.001}{10.9091}{rm}{41}\RG{302.780}{0.001}{10.9091}{rm}{61}\RG{314.314}{0.001}{10.9091}{mi}{105}\RG{318.072}{0.001}{10.9091}{rm}{40}\RG{322.315}{0.001}{10.9091}{mi}{119}\RG{330.418}{0.001}{10.9091}{rm}{41}\RG{334.661}{0.001}{10.9091}{mi}{58}\RG{113.493}{-1.635}{7.9701}{rm}{49}\RG{167.057}{-1.635}{7.9701}{rm}{49}\RG{227.035}{-1.635}{7.9701}{rm}{49}\RG{240.783}{-1.635}{7.9701}{rm}{49}%
}
\]
所要求的两式均成立。由于\allowbreak{}\hskip.1em \RB{5.365}{8.103}{2.737}{%
\RG{0.000}{-0.001}{10.9091}{mi}{24}%
}\hskip.1em\allowbreak{} 为满射，这样的\allowbreak{}\hskip.1em \RB{5.911}{5.321}{2.857}{%
\RG{0.000}{-0.001}{10.9091}{mi}{17}%
}\hskip.1em\allowbreak{} 还是唯一的。\allowbreak{}\hskip.1em \RB{8.366}{7.994}{0.501}{%
\RG{0.000}{-0.001}{10.9091}{am}{3}%
}\hskip.1em\allowbreak{}

方法提炼：像分解的唯一性。每个线性映射均能唯一到自然同构地分解为满射后接单射；关键是单射因子保证两个满射具有相同的核。

题目来源：谢启鸿《复旦大学高等代数习题集（2023版）》，第109页。

% Original page 6.
\section*{2021A14}
设\allowbreak{}\hskip.1em \RB{8.889}{7.955}{0.736}{%
\RG{0.000}{0.004}{10.9091}{mi}{86}%
}\hskip.1em\allowbreak{} 是\allowbreak{}\hskip.1em \RB{6.729}{5.326}{0.616}{%
\RG{0.000}{0.004}{10.9091}{mi}{110}%
}\hskip.1em\allowbreak{} 维复线性空间，\allowbreak{}\hskip.1em \RB{50.874}{7.955}{2.874}{%
\RG{0.000}{0.004}{10.9091}{mi}{39}\RG{10.167}{0.004}{10.9091}{rm}{58}\RG{16.227}{0.004}{10.9091}{mi}{86}\RG{28.045}{0.004}{10.9091}{sy}{33}\RG{41.985}{0.004}{10.9091}{mi}{86}%
}\hskip.1em\allowbreak{} 为复线性变换。将\allowbreak{}\hskip.1em \RB{8.889}{7.955}{0.736}{%
\RG{0.000}{0.004}{10.9091}{mi}{86}%
}\hskip.1em\allowbreak{} 看成\allowbreak{}\hskip.1em \RB{12.183}{7.770}{0.616}{%
\RG{0.000}{0.004}{10.9091}{rm}{50}\RG{5.454}{0.004}{10.9091}{mi}{110}%
}\hskip.1em\allowbreak{} 维实线性空间\allowbreak{}\hskip.1em \RB{10.761}{7.955}{2.300}{%
\RG{0.000}{0.004}{10.9091}{mi}{86}\RG{6.364}{-1.633}{7.9701}{rm}{48}%
}\hskip.1em\allowbreak{} ，并将\allowbreak{}\hskip.1em \RB{7.253}{5.326}{2.874}{%
\RG{0.000}{0.004}{10.9091}{mi}{39}%
}\hskip.1em\allowbreak{} 视作实线性变换\allowbreak{}\hskip.1em \RB{11.534}{5.322}{2.878}{%
\RG{0.000}{-0.000}{10.9091}{mi}{39}\RG{7.137}{-1.636}{7.9701}{rm}{48}%
}\hskip.1em\allowbreak{} 。证明：

\[
\RB{96.032}{10.301}{3.231}{%
\RG{88.706}{4.501}{7.9701}{rm}{50}\RT{0.000}{-0.004}{10.9091}{d}\RT{6.066}{-0.004}{10.9091}{e}\RT{10.909}{-0.004}{10.9091}{t}\RG{15.153}{-0.004}{10.9091}{rm}{40}\RG{19.396}{-0.004}{10.9091}{mi}{39}\RG{31.264}{-0.004}{10.9091}{rm}{41}\RG{38.519}{-0.004}{10.9091}{rm}{61}\RG{50.052}{-0.004}{10.9091}{sy}{106}\RT{54.900}{-0.004}{10.9091}{d}\RT{60.966}{-0.004}{10.9091}{e}\RT{65.810}{-0.004}{10.9091}{t}\RG{70.054}{-0.004}{10.9091}{rm}{40}\RG{74.297}{-0.004}{10.9091}{mi}{39}\RG{81.433}{-0.004}{10.9091}{rm}{41}\RG{85.675}{-0.004}{10.9091}{sy}{106}\RG{93.438}{-0.004}{10.9091}{mi}{58}\RG{26.532}{-1.640}{7.9701}{rm}{48}%
}
\]
解：任取\allowbreak{}\hskip.1em \RB{8.889}{7.949}{0.742}{%
\RG{0.000}{-0.002}{10.9091}{mi}{86}%
}\hskip.1em\allowbreak{} 的复基\allowbreak{}\hskip.1em \RB{44.511}{5.320}{2.607}{%
\RG{0.000}{-0.002}{10.9091}{mi}{101}\RG{5.079}{-1.638}{7.9701}{rm}{49}\RG{9.811}{-0.002}{10.9091}{mi}{59}\RG{14.644}{-0.002}{10.9091}{mi}{58}\RG{19.488}{-0.002}{10.9091}{mi}{58}\RG{24.331}{-0.002}{10.9091}{mi}{58}\RG{29.164}{-0.002}{10.9091}{mi}{59}\RG{34.008}{-0.002}{10.9091}{mi}{101}\RG{39.133}{-1.638}{7.9701}{mi}{110}%
}\hskip.1em\allowbreak{} ，设\allowbreak{}\hskip.1em \RB{7.253}{5.320}{2.880}{%
\RG{0.000}{-0.002}{10.9091}{mi}{39}%
}\hskip.1em\allowbreak{} 在该基下的表示矩阵为

\[
\RB{153.313}{9.289}{3.231}{%
\RG{0.000}{-0.004}{10.9091}{mi}{65}\RG{11.212}{-0.004}{10.9091}{rm}{61}\RG{22.727}{-0.004}{10.9091}{mi}{66}\RG{33.973}{-0.004}{10.9091}{rm}{43}\RT{44.883}{-0.004}{10.9091}{i}\RG{47.915}{-0.004}{10.9091}{mi}{67}\RG{55.868}{-0.004}{10.9091}{mi}{59}\RG{82.530}{-0.004}{10.9091}{mi}{66}\RG{91.345}{-0.004}{10.9091}{mi}{59}\RG{96.189}{-0.004}{10.9091}{mi}{67}\RG{107.832}{-0.004}{10.9091}{sy}{50}\RG{118.135}{-0.004}{10.9091}{mi}{77}\RG{134.355}{-0.004}{10.9091}{rm}{40}\RG{138.597}{-0.004}{10.9091}{bb}{82}\RG{146.476}{-0.004}{10.9091}{rm}{41}\RG{150.718}{-0.004}{10.9091}{mi}{58}\RG{128.718}{-1.640}{7.9701}{mi}{110}%
}
\]
考虑\allowbreak{}\hskip.1em \RB{8.889}{7.954}{0.737}{%
\RG{0.000}{0.003}{10.9091}{mi}{86}%
}\hskip.1em\allowbreak{} 的实基

\[
\RB{106.682}{9.293}{2.616}{%
\RG{0.000}{-0.000}{10.9091}{mi}{101}\RG{9.812}{-0.000}{10.9091}{mi}{59}\RG{14.645}{-0.000}{10.9091}{mi}{58}\RG{19.489}{-0.000}{10.9091}{mi}{58}\RG{24.332}{-0.000}{10.9091}{mi}{58}\RG{29.165}{-0.000}{10.9091}{mi}{59}\RG{34.009}{-0.000}{10.9091}{mi}{101}\RG{44.770}{-0.000}{10.9091}{mi}{59}\RT{53.252}{-0.000}{10.9091}{i}\RG{56.284}{-0.000}{10.9091}{mi}{101}\RG{66.097}{-0.000}{10.9091}{mi}{59}\RG{70.929}{-0.000}{10.9091}{mi}{58}\RG{75.773}{-0.000}{10.9091}{mi}{58}\RG{80.617}{-0.000}{10.9091}{mi}{58}\RG{85.449}{-0.000}{10.9091}{mi}{59}\RT{90.340}{-0.000}{10.9091}{i}\RG{93.372}{-0.000}{10.9091}{mi}{101}\RG{104.088}{-0.000}{10.9091}{mi}{58}\RG{5.080}{-1.636}{7.9701}{rm}{49}\RG{39.134}{-1.636}{7.9701}{mi}{110}\RG{61.365}{-1.636}{7.9701}{rm}{49}\RG{98.452}{-1.636}{7.9701}{mi}{110}%
}
\]
它确为实基：任意\allowbreak{}\hskip.1em \RB{41.079}{8.683}{5.351}{%
\RG{0.000}{8.182}{10.9092}{ex}{80}\RG{11.516}{-3.225}{7.9701}{mi}{106}\RG{15.898}{0.001}{10.9091}{rm}{40}\RG{20.140}{0.001}{10.9091}{mi}{97}\RG{25.907}{-1.636}{7.9701}{mi}{106}\RG{32.713}{0.001}{10.9091}{rm}{43}%
}\hskip.1em\allowbreak{} i\allowbreak{}\hskip.1em \RB{22.178}{8.683}{3.762}{%
\RG{0.000}{0.001}{10.9091}{mi}{98}\RG{4.682}{-1.636}{7.9701}{mi}{106}\RG{9.064}{0.001}{10.9091}{rm}{41}\RG{13.306}{0.001}{10.9091}{mi}{101}\RG{18.386}{-1.636}{7.9701}{mi}{106}%
}\hskip.1em\allowbreak{} 都可唯一地写成\allowbreak{}\hskip.1em \RB{81.550}{8.683}{5.351}{%
\RG{0.000}{8.182}{10.9092}{ex}{80}\RG{11.515}{-3.225}{7.9701}{mi}{106}\RG{17.715}{0.001}{10.9091}{mi}{97}\RG{23.482}{-1.636}{7.9701}{mi}{106}\RG{27.864}{0.001}{10.9091}{mi}{101}\RG{32.944}{-1.636}{7.9701}{mi}{106}\RG{39.750}{0.001}{10.9091}{rm}{43}\RG{50.659}{8.182}{10.9092}{ex}{80}\RG{62.175}{-3.225}{7.9701}{mi}{106}\RG{68.375}{0.001}{10.9091}{mi}{98}\RG{73.057}{-1.636}{7.9701}{mi}{106}\RG{77.439}{0.001}{10.9091}{rm}{40}%
}\hskip.1em\allowbreak{} i\allowbreak{}\hskip.1em \RB{13.114}{8.683}{3.762}{%
\RG{0.000}{0.001}{10.9091}{mi}{101}\RG{5.079}{-1.636}{7.9701}{mi}{106}\RG{9.461}{0.001}{10.9091}{rm}{41}%
}\hskip.1em\allowbreak{} ，其中\allowbreak{}\hskip.1em \RB{45.563}{8.072}{3.762}{%
\RG{0.000}{0.001}{10.9091}{mi}{97}\RG{5.766}{-1.636}{7.9701}{mi}{106}\RG{10.148}{0.001}{10.9091}{mi}{59}\RG{14.981}{0.001}{10.9091}{mi}{98}\RG{19.679}{-1.636}{7.9701}{mi}{106}\RG{27.091}{0.001}{10.9091}{sy}{50}\RG{37.394}{0.001}{10.9091}{bb}{82}%
}\hskip.1em\allowbreak{} 。

由\allowbreak{}\hskip.1em \RB{81.504}{8.679}{3.808}{%
\RG{0.000}{-0.003}{10.9091}{mi}{39}\RG{7.136}{-0.003}{10.9091}{rm}{40}\RG{11.379}{-0.003}{10.9091}{mi}{101}\RG{16.458}{-1.639}{7.9701}{mi}{106}\RG{20.840}{-0.003}{10.9091}{rm}{41}\RG{28.095}{-0.003}{10.9091}{rm}{61}\RG{39.628}{8.179}{10.9092}{ex}{80}\RG{51.143}{-3.228}{7.9701}{mi}{105}\RG{54.525}{-0.003}{10.9091}{rm}{40}\RG{58.767}{-0.003}{10.9091}{mi}{98}\RG{63.450}{-1.639}{7.9701}{mi}{105}\RG{66.327}{-1.639}{7.9701}{mi}{106}\RG{73.139}{-0.003}{10.9091}{rm}{43}%
}\hskip.1em\allowbreak{} i\allowbreak{}\hskip.1em \RB{24.336}{8.679}{3.765}{%
\RG{0.000}{-0.003}{10.9091}{mi}{99}\RG{4.721}{-1.639}{7.9701}{mi}{105}\RG{7.598}{-1.639}{7.9701}{mi}{106}\RG{11.987}{-0.003}{10.9091}{rm}{41}\RG{16.229}{-0.003}{10.9091}{mi}{101}\RG{21.309}{-1.639}{7.9701}{mi}{105}%
}\hskip.1em\allowbreak{} 和复线性性\allowbreak{}\hskip.1em \RB{11.247}{8.679}{3.230}{%
\RG{0.000}{-0.003}{10.9091}{mi}{39}\RG{7.136}{-0.003}{10.9091}{rm}{40}%
}\hskip.1em\allowbreak{} i\allowbreak{}\hskip.1em \RB{25.082}{8.679}{3.765}{%
\RG{0.000}{-0.003}{10.9091}{mi}{101}\RG{5.080}{-1.639}{7.9701}{mi}{106}\RG{9.462}{-0.003}{10.9091}{rm}{41}\RG{16.717}{-0.003}{10.9091}{rm}{61}%
}\hskip.1em\allowbreak{} i\allowbreak{}\hskip.1em \RB{24.494}{8.679}{3.765}{%
\RG{0.000}{-0.003}{10.9091}{mi}{39}\RG{7.137}{-0.003}{10.9091}{rm}{40}\RG{11.379}{-0.003}{10.9091}{mi}{101}\RG{16.459}{-1.639}{7.9701}{mi}{106}\RG{20.841}{-0.003}{10.9091}{rm}{41}%
}\hskip.1em\allowbreak{} 可知，\allowbreak{}\hskip.1em \RB{11.534}{5.319}{2.881}{%
\RG{0.000}{-0.003}{10.9091}{mi}{39}\RG{7.137}{-1.639}{7.9701}{rm}{48}%
}\hskip.1em\allowbreak{} 在上述实基下的矩阵为

\[
\RB{80.445}{12.026}{21.690}{%
\RG{22.912}{11.090}{10.9092}{ex}{32}\RG{67.396}{11.090}{10.9092}{ex}{33}\RG{31.549}{0.000}{10.9091}{mi}{66}\RG{50.333}{0.000}{10.9091}{sy}{0}\RG{58.818}{0.000}{10.9091}{mi}{67}\RG{0.000}{-7.565}{10.9091}{mi}{82}\RG{11.397}{-7.565}{10.9091}{rm}{61}\RG{77.851}{-7.565}{10.9091}{mi}{58}\RG{31.671}{-16.124}{10.9091}{mi}{67}\RG{54.439}{-16.124}{10.9091}{mi}{66}%
}
\]
在复数域上取非异分块矩阵

\[
\RB{88.926}{28.154}{5.563}{%
\RG{21.864}{27.217}{10.9092}{ex}{32}\RG{75.878}{27.217}{10.9092}{ex}{33}\RG{36.259}{16.127}{10.9091}{mi}{73}\RG{63.930}{16.127}{10.9091}{mi}{73}\RG{41.054}{14.491}{7.9701}{mi}{110}\RG{68.725}{14.491}{7.9701}{mi}{110}\RG{0.000}{8.563}{10.9091}{mi}{83}\RG{10.349}{8.563}{10.9091}{rm}{61}\RG{86.332}{8.563}{10.9091}{mi}{58}\RG{30.500}{0.004}{10.9091}{sy}{0}\RT{38.985}{0.003}{10.9091}{i}\RG{42.018}{0.004}{10.9091}{mi}{73}\RT{62.413}{0.004}{10.9091}{i}\RG{65.446}{0.003}{10.9091}{mi}{73}\RG{46.813}{-1.633}{7.9701}{mi}{110}\RG{70.241}{-1.633}{7.9701}{mi}{110}%
}
\]
直接分块相乘可得

\[
\RB{138.667}{12.029}{21.687}{%
\RG{39.367}{11.093}{10.9092}{ex}{32}\RG{125.498}{11.093}{10.9092}{ex}{33}\RG{48.004}{0.003}{10.9091}{mi}{66}\RG{59.251}{0.003}{10.9091}{rm}{43}\RT{70.160}{0.003}{10.9091}{i}\RG{73.193}{0.003}{10.9091}{mi}{67}\RG{104.295}{0.003}{10.9091}{mi}{79}\RG{0.000}{-7.561}{10.9091}{mi}{82}\RG{8.356}{-7.561}{10.9091}{mi}{83}\RG{18.716}{-7.561}{10.9091}{rm}{61}\RG{30.231}{-7.561}{10.9091}{mi}{83}\RG{135.952}{-7.561}{10.9091}{mi}{59}\RG{60.574}{-16.120}{10.9091}{mi}{79}\RG{91.720}{-16.120}{10.9091}{mi}{66}\RG{102.979}{-16.120}{10.9091}{sy}{0}\RT{113.888}{-16.120}{10.9091}{i}\RG{116.921}{-16.120}{10.9091}{mi}{67}%
}
\]
所以\allowbreak{}\hskip.1em \RB{8.736}{7.955}{0.736}{%
\RG{0.000}{0.004}{10.9091}{mi}{82}%
}\hskip.1em\allowbreak{} 与diag\allowbreak{}\hskip.1em \RB{29.108}{8.686}{3.223}{%
\RG{0.000}{0.004}{10.9091}{rm}{40}\RG{4.243}{0.004}{10.9091}{mi}{65}\RG{12.425}{0.004}{10.9091}{mi}{59}\RG{17.273}{0.004}{10.9091}{mi}{65}\RG{25.455}{0.004}{10.9091}{rm}{41}\RR{17.273}{8.986}{8.182}{0.436}%
}\hskip.1em\allowbreak{} 在\allowbreak{}\hskip.1em \RB{7.973}{8.184}{0.703}{%
\RG{0.000}{0.004}{10.9091}{bb}{67}%
}\hskip.1em\allowbreak{} 上相似。于是

\[
\RB{215.011}{28.404}{3.228}{%
\RT{0.000}{19.111}{10.9091}{d}\RT{6.066}{19.111}{10.9091}{e}\RT{10.909}{19.111}{10.9091}{t}\RG{15.153}{19.111}{10.9091}{rm}{40}\RG{19.395}{19.111}{10.9091}{mi}{39}\RG{31.264}{19.111}{10.9091}{rm}{41}\RG{38.519}{19.111}{10.9091}{rm}{61}\RT{50.052}{19.111}{10.9091}{d}\RT{56.118}{19.111}{10.9091}{e}\RT{60.962}{19.111}{10.9091}{t}\RG{67.024}{19.111}{10.9091}{mi}{82}\RG{78.421}{19.111}{10.9091}{rm}{61}\RT{89.937}{19.111}{10.9091}{d}\RT{96.002}{19.111}{10.9091}{e}\RT{100.846}{19.111}{10.9091}{t}\RG{106.907}{19.111}{10.9091}{mi}{65}\RT{118.726}{19.111}{10.9091}{d}\RT{124.791}{19.111}{10.9091}{e}\RT{129.635}{19.111}{10.9091}{t}\RG{135.696}{19.111}{10.9091}{mi}{65}\RG{26.532}{17.475}{7.9701}{rm}{48}\RG{149.753}{4.503}{7.9701}{rm}{50}\RG{207.685}{4.503}{7.9701}{rm}{50}\RG{38.537}{-0.001}{10.9091}{rm}{61}\RT{50.052}{-0.001}{10.9091}{d}\RT{56.118}{-0.001}{10.9091}{e}\RT{60.962}{-0.001}{10.9091}{t}\RG{67.024}{-0.001}{10.9091}{mi}{65}\RT{77.024}{-0.001}{10.9091}{d}\RT{83.089}{-0.001}{10.9091}{e}\RT{87.933}{-0.001}{10.9091}{t}\RG{93.995}{-0.001}{10.9091}{mi}{65}\RG{105.207}{-0.001}{10.9091}{rm}{61}\RG{116.722}{-0.001}{10.9091}{sy}{106}\RT{121.570}{-0.001}{10.9091}{d}\RT{127.636}{-0.001}{10.9091}{e}\RT{132.479}{-0.001}{10.9091}{t}\RG{138.541}{-0.001}{10.9091}{mi}{65}\RG{146.723}{-0.001}{10.9091}{sy}{106}\RG{157.516}{-0.001}{10.9091}{rm}{61}\RG{169.031}{-0.001}{10.9091}{sy}{106}\RT{173.879}{-0.001}{10.9091}{d}\RT{179.945}{-0.001}{10.9091}{e}\RT{184.789}{-0.001}{10.9091}{t}\RG{189.033}{-0.001}{10.9091}{rm}{40}\RG{193.276}{-0.001}{10.9091}{mi}{39}\RG{200.412}{-0.001}{10.9091}{rm}{41}\RG{204.654}{-0.001}{10.9091}{sy}{106}\RG{212.417}{-0.001}{10.9091}{mi}{58}\RR{135.696}{28.093}{8.182}{0.436}\RR{77.024}{9.097}{25.152}{0.436}%
}
\]
这里并未假定\allowbreak{}\hskip.1em \RB{22.438}{8.188}{2.608}{%
\RG{0.000}{-0.003}{10.9091}{mi}{66}\RG{8.815}{-0.003}{10.9091}{mi}{59}\RG{13.647}{-0.003}{10.9091}{mi}{67}%
}\hskip.1em\allowbreak{} 乘法可交换。\allowbreak{}\hskip.1em \RB{8.366}{7.992}{0.503}{%
\RG{0.000}{-0.003}{10.9091}{am}{3}%
}\hskip.1em\allowbreak{}

\noindent \RB{465.386}{10.105}{18.800}{%
\RG{327.555}{8.842}{10.9092}{ex}{0}\RG{359.186}{8.842}{10.9092}{ex}{1}\RG{334.373}{3.939}{7.9701}{mi}{66}\RG{344.183}{3.939}{7.9701}{sy}{0}\RG{350.769}{3.939}{7.9701}{mi}{67}\RT{0.000}{0.005}{10.9091}{方}\RT{10.909}{0.005}{10.9091}{法}\RT{21.818}{0.005}{10.9091}{提}\RT{32.728}{0.005}{10.9091}{炼}\RT{43.637}{0.005}{10.9091}{：}\RT{56.466}{0.005}{10.9091}{复}\RT{67.375}{0.005}{10.9091}{线}\RT{78.284}{0.005}{10.9091}{性}\RT{89.194}{0.005}{10.9091}{映}\RT{100.103}{0.005}{10.9091}{射}\RT{111.012}{0.005}{10.9091}{的}\RT{121.921}{0.005}{10.9091}{实}\RT{132.830}{0.005}{10.9091}{化}\RT{143.740}{0.005}{10.9091}{。}\RT{153.248}{0.005}{10.9091}{复}\RT{164.158}{0.005}{10.9091}{矩}\RT{175.067}{0.005}{10.9091}{阵}\RG{189.298}{0.005}{10.9091}{mi}{65}\RG{200.510}{0.005}{10.9091}{rm}{61}\RG{212.025}{0.005}{10.9091}{mi}{66}\RG{222.649}{0.005}{10.9091}{rm}{43}\RT{232.936}{0.005}{10.9091}{i}\RG{235.969}{0.005}{10.9091}{mi}{67}\RT{247.868}{0.005}{10.9091}{对}\RT{258.777}{0.005}{10.9091}{应}\RT{269.686}{0.005}{10.9091}{的}\RT{280.596}{0.005}{10.9091}{实}\RT{291.505}{0.005}{10.9091}{矩}\RT{302.414}{0.005}{10.9091}{阵}\RT{313.323}{0.005}{10.9091}{为}\RT{364.186}{0.005}{10.9091}{；}\RT{373.699}{0.005}{10.9091}{复}\RT{384.608}{0.005}{10.9091}{化}\RT{395.517}{0.005}{10.9091}{后}\RT{406.426}{0.005}{10.9091}{分}\RT{417.336}{0.005}{10.9091}{解}\RT{428.245}{0.005}{10.9091}{成}\RG{442.473}{0.005}{10.9091}{mi}{65}\RT{453.977}{0.005}{10.9091}{与}\RG{334.464}{-3.920}{7.9701}{mi}{67}\RG{347.375}{-3.920}{7.9701}{mi}{66}\RG{0.000}{-16.118}{10.9091}{mi}{65}\RT{8.182}{-16.118}{10.9091}{，}\RT{19.091}{-16.118}{10.9091}{从}\RT{30.000}{-16.118}{10.9091}{而}\RT{40.910}{-16.118}{10.9091}{行}\RT{51.819}{-16.118}{10.9091}{列}\RT{62.728}{-16.118}{10.9091}{式}\RT{73.637}{-16.118}{10.9091}{等}\RT{84.546}{-16.118}{10.9091}{于}\RT{95.456}{-16.118}{10.9091}{模}\RT{106.365}{-16.118}{10.9091}{平}\RT{117.274}{-16.118}{10.9091}{方}\RT{128.183}{-16.118}{10.9091}{。}\RR{0.000}{-7.136}{8.182}{0.436}%
}

题目来源：谢启鸿《复旦大学高等代数习题集（2023版）》，第109页。

% Original page 7.
\section*{2021A15}
(1)设\allowbreak{}\hskip.1em \RB{229.370}{9.757}{3.229}{%
\RG{0.000}{-0.002}{10.9091}{mi}{102}\RG{6.515}{-0.002}{10.9091}{rm}{40}\RG{10.758}{-0.002}{10.9091}{mi}{120}\RG{16.993}{-0.002}{10.9091}{rm}{41}\RG{24.248}{-0.002}{10.9091}{rm}{61}\RG{35.781}{-0.002}{10.9091}{mi}{97}\RG{41.547}{-1.638}{7.9701}{mi}{110}\RG{47.183}{-0.002}{10.9091}{mi}{120}\RG{53.418}{3.957}{7.9701}{mi}{110}\RG{61.480}{-0.002}{10.9091}{rm}{43}\RG{72.389}{-0.002}{10.9091}{mi}{97}\RG{78.155}{-1.638}{7.9701}{mi}{110}\RG{83.293}{-1.638}{7.9701}{sy}{0}\RG{89.880}{-1.638}{7.9701}{rm}{49}\RG{94.612}{-0.002}{10.9091}{mi}{120}\RG{100.847}{3.957}{7.9701}{mi}{110}\RG{105.985}{3.957}{7.9701}{sy}{0}\RG{112.572}{3.957}{7.9701}{rm}{49}\RG{119.728}{-0.002}{10.9091}{rm}{43}\RG{130.637}{-0.002}{10.9091}{sy}{1}\RG{135.470}{-0.002}{10.9091}{sy}{1}\RG{140.314}{-0.002}{10.9091}{sy}{1}\RG{145.789}{-0.002}{10.9091}{rm}{43}\RG{156.698}{-0.002}{10.9091}{mi}{97}\RG{162.464}{-1.638}{7.9701}{rm}{49}\RG{167.197}{-0.002}{10.9091}{mi}{120}\RG{175.856}{-0.002}{10.9091}{rm}{43}\RG{186.766}{-0.002}{10.9091}{mi}{97}\RG{192.532}{-1.638}{7.9701}{rm}{48}\RG{200.295}{-0.002}{10.9091}{sy}{50}\RG{210.598}{-0.002}{10.9091}{bb}{90}\RG{217.870}{-0.002}{10.9091}{rm}{91}\RG{220.901}{-0.002}{10.9091}{mi}{120}\RG{227.136}{-0.002}{10.9091}{rm}{93}%
}\hskip.1em\allowbreak{} ，\allowbreak{}\hskip.1em \RB{47.298}{8.069}{1.997}{%
\RG{0.000}{-0.002}{10.9091}{rm}{49}\RG{8.485}{-0.002}{10.9091}{sy}{20}\RG{20.000}{-0.002}{10.9091}{mi}{107}\RG{29.053}{-0.002}{10.9091}{sy}{20}\RG{40.568}{-0.002}{10.9091}{mi}{110}%
}\hskip.1em\allowbreak{} 。若存在素数\allowbreak{}\hskip.1em \RB{5.846}{5.320}{2.618}{%
\RG{0.000}{-0.002}{10.9091}{mi}{112}%
}\hskip.1em\allowbreak{} ，满足

\[
\RB{198.217}{10.302}{3.229}{%
\RG{171.180}{4.502}{7.9701}{rm}{50}\RG{0.000}{-0.002}{10.9091}{mi}{112}\RG{8.519}{-0.002}{10.9091}{bb}{45}\RG{14.580}{-0.002}{10.9091}{mi}{97}\RG{25.983}{-0.002}{10.9091}{mi}{59}\RG{52.634}{-0.002}{10.9091}{mi}{112}\RG{61.169}{-0.002}{10.9091}{sy}{106}\RG{67.230}{-0.002}{10.9091}{mi}{97}\RG{77.728}{-0.002}{10.9091}{mi}{59}\RG{82.560}{-0.002}{10.9091}{mi}{97}\RG{93.076}{-0.002}{10.9091}{mi}{59}\RG{97.909}{-0.002}{10.9091}{mi}{58}\RG{102.752}{-0.002}{10.9091}{mi}{58}\RG{107.585}{-0.002}{10.9091}{mi}{58}\RG{112.429}{-0.002}{10.9091}{mi}{59}\RG{117.272}{-0.002}{10.9091}{mi}{97}\RG{139.025}{-0.002}{10.9091}{mi}{59}\RG{165.676}{-0.002}{10.9091}{mi}{112}\RG{178.943}{-0.002}{10.9091}{bb}{45}\RG{185.004}{-0.002}{10.9091}{mi}{97}\RG{195.503}{-0.002}{10.9091}{mi}{59}\RG{20.346}{-1.639}{7.9701}{mi}{110}\RG{72.996}{-1.639}{7.9701}{rm}{48}\RG{88.343}{-1.639}{7.9701}{rm}{49}\RG{190.771}{-1.639}{7.9701}{rm}{48}\RG{123.084}{-1.780}{7.9701}{mi}{107}\RG{127.706}{-1.780}{7.9701}{sy}{0}\RG{134.293}{-1.780}{7.9701}{rm}{49}%
}
\]
证明：\allowbreak{}\hskip.1em \RB{6.522}{8.188}{2.739}{%
\RG{0.000}{-0.003}{10.9091}{mi}{102}%
}\hskip.1em\allowbreak{} 在\allowbreak{}\hskip.1em \RB{19.985}{8.679}{3.230}{%
\RG{0.000}{-0.003}{10.9091}{bb}{81}\RG{8.484}{-0.003}{10.9091}{rm}{91}\RG{11.515}{-0.003}{10.9091}{mi}{120}\RG{17.750}{-0.003}{10.9091}{rm}{93}%
}\hskip.1em\allowbreak{} 中有一个次数不小于\allowbreak{}\hskip.1em \RB{6.042}{8.068}{0.623}{%
\RG{0.000}{-0.003}{10.9091}{mi}{107}%
}\hskip.1em\allowbreak{} 的不可约因子。

(2)对任意整数\allowbreak{}\hskip.1em \RB{26.491}{7.770}{1.991}{%
\RG{0.000}{0.004}{10.9091}{mi}{110}\RG{9.578}{0.004}{10.9091}{sy}{21}\RG{21.093}{0.004}{10.9091}{rm}{50}%
}\hskip.1em\allowbreak{} 和素数\allowbreak{}\hskip.1em \RB{5.846}{5.326}{2.612}{%
\RG{0.000}{0.004}{10.9091}{mi}{112}%
}\hskip.1em\allowbreak{} ，证明\allowbreak{}\hskip.1em \RB{99.837}{9.762}{3.223}{%
\RG{0.000}{0.004}{10.9091}{mi}{120}\RG{6.235}{3.962}{7.9701}{mi}{110}\RG{14.296}{0.004}{10.9091}{rm}{43}\RG{25.194}{0.004}{10.9091}{rm}{40}\RG{29.447}{0.004}{10.9091}{mi}{112}\RG{37.360}{0.004}{10.9091}{rm}{43}\RG{48.259}{0.004}{10.9091}{rm}{50}\RG{53.713}{0.004}{10.9091}{rm}{41}\RG{57.966}{0.004}{10.9091}{mi}{120}\RG{64.202}{3.962}{7.9701}{mi}{110}\RG{69.340}{3.962}{7.9701}{sy}{0}\RG{75.927}{3.962}{7.9701}{rm}{49}\RG{83.083}{0.004}{10.9091}{rm}{43}\RG{93.992}{0.004}{10.9091}{mi}{112}%
}\hskip.1em\allowbreak{} 在\allowbreak{}\hskip.1em \RB{8.627}{8.184}{2.449}{%
\RG{0.000}{0.004}{10.9091}{bb}{81}%
}\hskip.1em\allowbreak{} 上不可约。

解：(1)由Gauss引理，取\allowbreak{}\hskip.1em \RB{6.522}{8.193}{2.734}{%
\RG{0.000}{0.002}{10.9091}{mi}{102}%
}\hskip.1em\allowbreak{} 在整数多项式环中的不可约因子分解

\[
\RB{110.391}{8.683}{3.226}{%
\RG{0.000}{0.001}{10.9091}{mi}{102}\RG{6.515}{0.001}{10.9091}{rm}{40}\RG{10.758}{0.001}{10.9091}{mi}{120}\RG{16.993}{0.001}{10.9091}{rm}{41}\RG{24.248}{0.001}{10.9091}{rm}{61}\RG{35.780}{0.001}{10.9091}{mi}{99}\RG{42.304}{0.001}{10.9091}{mi}{103}\RG{52.255}{0.001}{10.9091}{rm}{40}\RG{56.498}{0.001}{10.9091}{mi}{120}\RG{62.732}{0.001}{10.9091}{rm}{41}\RG{68.794}{0.001}{10.9091}{sy}{1}\RG{73.626}{0.001}{10.9091}{sy}{1}\RG{78.470}{0.001}{10.9091}{sy}{1}\RG{83.339}{0.001}{10.9091}{mi}{103}\RG{92.957}{0.001}{10.9091}{rm}{40}\RG{97.199}{0.001}{10.9091}{mi}{120}\RG{103.434}{0.001}{10.9091}{rm}{41}\RG{107.677}{0.001}{10.9091}{mi}{59}\RG{47.523}{-1.635}{7.9701}{rm}{49}\RG{88.543}{-1.635}{7.9701}{mi}{115}%
}
\]
其中\allowbreak{}\hskip.1em \RB{51.698}{8.683}{3.226}{%
\RG{0.000}{0.001}{10.9091}{mi}{99}\RG{7.751}{0.001}{10.9091}{sy}{50}\RG{18.054}{0.001}{10.9091}{bb}{90}\RG{27.751}{0.001}{10.9091}{sy}{110}\RG{35.628}{0.001}{10.9091}{sy}{102}\RG{41.084}{0.001}{10.9091}{rm}{48}\RG{46.539}{0.001}{10.9091}{sy}{103}%
}\hskip.1em\allowbreak{} ，每个\allowbreak{}\hskip.1em \RB{41.691}{8.683}{3.761}{%
\RG{0.000}{0.001}{10.9091}{mi}{103}\RG{5.203}{-1.635}{7.9701}{mi}{106}\RG{12.615}{0.001}{10.9091}{sy}{50}\RG{22.918}{0.001}{10.9091}{bb}{90}\RG{30.191}{0.001}{10.9091}{rm}{91}\RG{33.221}{0.001}{10.9091}{mi}{120}\RG{39.456}{0.001}{10.9091}{rm}{93}%
}\hskip.1em\allowbreak{} 均为本原、正次数、在\allowbreak{}\hskip.1em \RB{19.985}{8.683}{3.226}{%
\RG{0.000}{0.001}{10.9091}{bb}{81}\RG{8.485}{0.001}{10.9091}{rm}{91}\RG{11.515}{0.001}{10.9091}{mi}{120}\RG{17.750}{0.001}{10.9091}{rm}{93}%
}\hskip.1em\allowbreak{} 中不可约的多项式。

由\allowbreak{}\hskip.1em \RB{25.724}{8.679}{3.230}{%
\RG{0.000}{-0.003}{10.9091}{mi}{112}\RG{8.519}{-0.003}{10.9091}{bb}{45}\RG{14.579}{-0.003}{10.9091}{mi}{97}\RG{20.346}{-1.639}{7.9701}{mi}{110}%
}\hskip.1em\allowbreak{} 知\allowbreak{}\hskip.1em \RB{19.770}{8.679}{3.230}{%
\RG{0.000}{-0.003}{10.9091}{mi}{112}\RG{8.519}{-0.003}{10.9091}{bb}{45}\RG{14.579}{-0.003}{10.9091}{mi}{99}%
}\hskip.1em\allowbreak{} ，而且每个\allowbreak{}\hskip.1em \RB{8.996}{5.319}{3.765}{%
\RG{0.000}{-0.003}{10.9091}{mi}{103}\RG{5.204}{-1.639}{7.9701}{mi}{106}%
}\hskip.1em\allowbreak{} 的首项系数均不被\allowbreak{}\hskip.1em \RB{5.846}{5.319}{2.619}{%
\RG{0.000}{-0.003}{10.9091}{mi}{112}%
}\hskip.1em\allowbreak{} 整除。由

\[
\RB{142.355}{17.651}{15.214}{%
\RG{51.557}{13.636}{7.9701}{mi}{115}\RG{46.546}{10.363}{10.9092}{ex}{89}\RG{118.367}{4.504}{7.9701}{rm}{50}\RG{0.000}{-0.001}{10.9091}{mi}{112}\RG{8.519}{-0.001}{10.9091}{sy}{106}\RG{14.580}{-0.001}{10.9091}{mi}{97}\RG{28.109}{-0.001}{10.9091}{rm}{61}\RG{39.624}{-0.001}{10.9091}{mi}{99}\RG{62.687}{-0.001}{10.9091}{mi}{103}\RG{72.272}{-0.001}{10.9091}{rm}{40}\RG{76.504}{-0.001}{10.9091}{rm}{48}\RG{81.959}{-0.001}{10.9091}{rm}{41}\RG{86.212}{-0.001}{10.9091}{mi}{59}\RG{112.863}{-0.001}{10.9091}{mi}{112}\RG{126.130}{-0.001}{10.9091}{bb}{45}\RG{132.190}{-0.001}{10.9091}{mi}{97}\RG{20.346}{-1.637}{7.9701}{rm}{48}\RG{67.890}{-1.637}{7.9701}{mi}{106}\RG{137.957}{-1.637}{7.9701}{rm}{48}\RG{46.163}{-13.088}{7.9701}{mi}{106}\RG{50.047}{-13.088}{7.9701}{rm}{61}\RG{56.631}{-13.088}{7.9701}{rm}{49}%
}
\]
可知，恰有一个因子（记为\allowbreak{}\hskip.1em \RB{12.416}{5.323}{3.761}{%
\RG{0.000}{0.001}{10.9091}{mi}{103}\RG{5.203}{-1.635}{7.9701}{mi}{106}\RG{8.628}{-2.742}{5.9776}{rm}{48}%
}\hskip.1em\allowbreak{} ）的常数项被\allowbreak{}\hskip.1em \RB{5.846}{5.323}{2.615}{%
\RG{0.000}{0.001}{10.9091}{mi}{112}%
}\hskip.1em\allowbreak{} 整除，其余因子的常数项均不被\allowbreak{}\hskip.1em \RB{5.846}{5.323}{2.615}{%
\RG{0.000}{0.001}{10.9091}{mi}{112}%
}\hskip.1em\allowbreak{} 整除。

在\allowbreak{}\hskip.1em \RB{22.928}{8.680}{3.685}{%
\RG{0.000}{-0.002}{10.9091}{bb}{70}\RG{6.667}{-1.639}{7.9701}{mi}{112}\RG{11.428}{-0.002}{10.9091}{rm}{91}\RG{14.458}{-0.002}{10.9091}{mi}{120}\RG{20.693}{-0.002}{10.9091}{rm}{93}%
}\hskip.1em\allowbreak{} 中用横线表示模\allowbreak{}\hskip.1em \RB{5.846}{5.320}{2.618}{%
\RG{0.000}{-0.002}{10.9091}{mi}{112}%
}\hskip.1em\allowbreak{} 约化。题设给出

\[
\RB{131.499}{10.533}{3.230}{%
\RG{6.235}{4.502}{7.9701}{mi}{107}\RG{0.000}{-0.003}{10.9091}{mi}{120}\RG{14.385}{-0.003}{10.9091}{sy}{106}\RG{20.445}{-0.003}{10.9091}{mi}{102}\RG{26.961}{-0.003}{10.9091}{rm}{40}\RG{31.203}{-0.003}{10.9091}{mi}{120}\RG{37.438}{-0.003}{10.9091}{rm}{41}\RG{44.693}{-0.003}{10.9091}{rm}{61}\RG{56.227}{-0.003}{10.9091}{mi}{99}\RG{62.766}{-0.003}{10.9091}{mi}{103}\RG{73.093}{-0.003}{10.9091}{rm}{40}\RG{77.335}{-0.003}{10.9091}{mi}{120}\RG{83.570}{-0.003}{10.9091}{rm}{41}\RG{89.631}{-0.003}{10.9091}{sy}{1}\RG{94.464}{-0.003}{10.9091}{sy}{1}\RG{99.307}{-0.003}{10.9091}{sy}{1}\RG{104.176}{-0.003}{10.9091}{mi}{103}\RG{114.185}{-0.003}{10.9091}{rm}{40}\RG{118.427}{-0.003}{10.9091}{mi}{120}\RG{124.662}{-0.003}{10.9091}{rm}{41}\RG{128.904}{-0.003}{10.9091}{mi}{58}\RG{68.360}{-2.622}{7.9701}{rm}{49}\RG{109.771}{-2.622}{7.9701}{mi}{115}\RR{20.445}{9.100}{6.516}{0.436}\RR{56.227}{6.221}{4.721}{0.436}\RR{62.766}{6.221}{5.594}{0.436}\RR{104.176}{6.221}{5.595}{0.436}%
}
\]
对\allowbreak{}\hskip.1em \RB{28.552}{8.308}{2.859}{%
\RG{0.000}{-0.003}{10.9091}{mi}{106}\RG{8.147}{-0.003}{10.9091}{sy}{54}\RG{8.147}{-0.003}{10.9091}{rm}{61}\RG{19.662}{-0.003}{10.9091}{mi}{106}\RG{24.154}{-1.639}{7.9701}{rm}{48}%
}\hskip.1em\allowbreak{} ，因\allowbreak{}\hskip.1em \RB{43.963}{8.679}{4.748}{%
\RG{0.000}{-0.003}{10.9091}{mi}{103}\RG{5.594}{-2.622}{7.9701}{mi}{106}\RG{9.976}{-0.003}{10.9091}{rm}{40}\RG{14.209}{-0.003}{10.9091}{rm}{48}\RG{19.663}{-0.003}{10.9091}{rm}{41}\RG{26.946}{-0.003}{10.9091}{sy}{54}\RG{26.946}{-0.003}{10.9091}{rm}{61}\RG{38.445}{-0.003}{10.9091}{rm}{48}\RR{0.000}{6.221}{5.594}{0.436}%
}\hskip.1em\allowbreak{} ，故\allowbreak{}\hskip.1em \RB{24.712}{8.679}{4.748}{%
\RG{0.000}{-0.003}{10.9091}{mi}{120}\RG{9.265}{-0.003}{10.9091}{bb}{45}\RG{15.325}{-0.003}{10.9091}{mi}{103}\RG{20.920}{-2.622}{7.9701}{mi}{106}\RR{15.325}{6.221}{5.595}{0.436}%
}\hskip.1em\allowbreak{} ；又\allowbreak{}\hskip.1em \RB{24.768}{8.308}{2.859}{%
\RG{0.000}{-0.003}{10.9091}{mi}{99}\RG{7.751}{-0.003}{10.9091}{sy}{54}\RG{7.751}{-0.003}{10.9091}{rm}{61}\RG{19.249}{-0.003}{10.9091}{rm}{48}\RR{0.000}{6.221}{4.721}{0.436}%
}\hskip.1em\allowbreak{} 。由\allowbreak{}\hskip.1em \RB{22.928}{8.679}{3.685}{%
\RG{0.000}{-0.003}{10.9091}{bb}{70}\RG{6.667}{-1.639}{7.9701}{mi}{112}\RG{11.428}{-0.003}{10.9091}{rm}{91}\RG{14.458}{-0.003}{10.9091}{mi}{120}\RG{20.693}{-0.003}{10.9091}{rm}{93}%
}\hskip.1em\allowbreak{} 中因子的整除性可得

\[
\RB{51.429}{10.532}{4.749}{%
\RG{6.235}{4.501}{7.9701}{mi}{107}\RG{0.000}{-0.003}{10.9091}{mi}{120}\RG{14.385}{-0.003}{10.9091}{sy}{106}\RG{20.445}{-0.003}{10.9091}{mi}{103}\RG{34.114}{-0.003}{10.9091}{rm}{40}\RG{38.356}{-0.003}{10.9091}{mi}{120}\RG{44.591}{-0.003}{10.9091}{rm}{41}\RG{48.834}{-0.003}{10.9091}{mi}{58}\RG{26.040}{-2.623}{7.9701}{mi}{106}\RG{29.465}{-3.730}{5.9776}{rm}{48}\RR{20.445}{6.221}{5.595}{0.436}%
}
\]
由于\allowbreak{}\hskip.1em \RB{12.416}{5.318}{3.766}{%
\RG{0.000}{-0.004}{10.9091}{mi}{103}\RG{5.203}{-1.640}{7.9701}{mi}{106}\RG{8.628}{-2.747}{5.9776}{rm}{48}%
}\hskip.1em\allowbreak{} 的首项系数不被\allowbreak{}\hskip.1em \RB{5.846}{5.318}{2.620}{%
\RG{0.000}{-0.004}{10.9091}{mi}{112}%
}\hskip.1em\allowbreak{} 整除，约化不降低其次数，因此

\[
\RB{101.006}{9.289}{4.749}{%
\RT{0.000}{-0.004}{10.9091}{d}\RT{6.066}{-0.004}{10.9091}{e}\RT{10.909}{-0.004}{10.9091}{g}\RG{18.182}{-0.004}{10.9091}{mi}{103}\RG{34.489}{-0.004}{10.9091}{rm}{61}\RT{46.005}{-0.004}{10.9091}{d}\RT{52.071}{-0.004}{10.9091}{e}\RT{56.915}{-0.004}{10.9091}{g}\RG{64.187}{-0.004}{10.9091}{mi}{103}\RG{80.886}{-0.004}{10.9091}{sy}{21}\RG{92.401}{-0.004}{10.9091}{mi}{107}\RG{98.412}{-0.004}{10.9091}{mi}{58}\RG{23.385}{-1.640}{7.9701}{mi}{106}\RG{69.782}{-2.623}{7.9701}{mi}{106}\RG{26.810}{-2.747}{5.9776}{rm}{48}\RG{73.207}{-3.730}{5.9776}{rm}{48}\RR{64.187}{6.220}{5.595}{0.436}%
}
\]
所求不可约因子即\allowbreak{}\hskip.1em \RB{12.416}{5.318}{3.767}{%
\RG{0.000}{-0.004}{10.9091}{mi}{103}\RG{5.204}{-1.641}{7.9701}{mi}{106}\RG{8.628}{-2.748}{5.9776}{rm}{48}%
}\hskip.1em\allowbreak{} 。(2)令

\[
\RB{139.870}{10.303}{3.229}{%
\RG{44.030}{4.503}{7.9701}{mi}{110}\RG{101.997}{4.503}{7.9701}{mi}{110}\RG{107.136}{4.503}{7.9701}{sy}{0}\RG{113.722}{4.503}{7.9701}{rm}{49}\RG{0.000}{-0.002}{10.9091}{mi}{70}\RG{8.531}{-0.002}{10.9091}{rm}{40}\RG{12.773}{-0.002}{10.9091}{mi}{120}\RG{19.008}{-0.002}{10.9091}{rm}{41}\RG{26.263}{-0.002}{10.9091}{rm}{61}\RG{37.796}{-0.002}{10.9091}{mi}{120}\RG{52.091}{-0.002}{10.9091}{rm}{43}\RG{62.990}{-0.002}{10.9091}{rm}{40}\RG{67.244}{-0.002}{10.9091}{mi}{112}\RG{75.156}{-0.002}{10.9091}{rm}{43}\RG{86.054}{-0.002}{10.9091}{rm}{50}\RG{91.508}{-0.002}{10.9091}{rm}{41}\RG{95.762}{-0.002}{10.9091}{mi}{120}\RG{120.879}{-0.002}{10.9091}{rm}{43}\RG{131.788}{-0.002}{10.9091}{mi}{112}\RG{137.275}{-0.002}{10.9091}{mi}{58}%
}
\]
取(1)中的\allowbreak{}\hskip.1em \RB{46.584}{8.071}{0.620}{%
\RG{0.000}{0.000}{10.9091}{mi}{107}\RG{9.433}{0.000}{10.9091}{rm}{61}\RG{21.328}{0.000}{10.9091}{mi}{110}\RG{30.452}{0.000}{10.9091}{sy}{0}\RG{41.513}{0.000}{10.9091}{rm}{49}%
}\hskip.1em\allowbreak{} 。\allowbreak{}\hskip.1em \RB{8.669}{7.918}{0.500}{%
\RG{0.000}{0.000}{10.9091}{mi}{70}%
}\hskip.1em\allowbreak{} 的首项系数为\allowbreak{}\hskip.1em \RB{5.071}{7.766}{0.500}{%
\RG{0.000}{0.000}{10.9091}{rm}{49}%
}\hskip.1em\allowbreak{} ，其常数项为\allowbreak{}\hskip.1em \RB{5.846}{5.322}{2.616}{%
\RG{0.000}{0.000}{10.9091}{mi}{112}%
}\hskip.1em\allowbreak{} ，而\allowbreak{}\hskip.1em \RB{68.554}{9.759}{2.605}{%
\RG{0.000}{0.000}{10.9091}{mi}{120}\RG{6.229}{0.000}{10.9091}{mi}{59}\RG{11.062}{0.000}{10.9091}{mi}{120}\RG{17.318}{3.959}{7.9701}{rm}{50}\RG{22.050}{0.000}{10.9091}{mi}{59}\RG{26.883}{0.000}{10.9091}{mi}{58}\RG{31.727}{0.000}{10.9091}{mi}{58}\RG{36.570}{0.000}{10.9091}{mi}{58}\RG{41.403}{0.000}{10.9091}{mi}{59}\RG{46.247}{0.000}{10.9091}{mi}{120}\RG{52.527}{3.959}{7.9701}{mi}{110}\RG{57.666}{3.959}{7.9701}{sy}{0}\RG{64.253}{3.959}{7.9701}{rm}{50}%
}\hskip.1em\allowbreak{} 的系数均为零，故(1)保证\allowbreak{}\hskip.1em \RB{8.669}{7.915}{0.503}{%
\RG{0.000}{-0.003}{10.9091}{mi}{70}%
}\hskip.1em\allowbreak{} 有一个次数至少为\allowbreak{}\hskip.1em \RB{25.214}{7.763}{0.623}{%
\RG{0.000}{-0.003}{10.9091}{mi}{110}\RG{9.103}{-0.003}{10.9091}{sy}{0}\RG{20.143}{-0.003}{10.9091}{rm}{49}%
}\hskip.1em\allowbreak{} 的不可约因子。若\allowbreak{}\hskip.1em \RB{8.669}{7.915}{0.503}{%
\RG{0.000}{-0.003}{10.9091}{mi}{70}%
}\hskip.1em\allowbreak{} 可约，其余因子的总次数至多为\allowbreak{}\hskip.1em \RB{5.071}{7.763}{0.503}{%
\RG{0.000}{-0.003}{10.9091}{rm}{49}%
}\hskip.1em\allowbreak{} ，因此\allowbreak{}\hskip.1em \RB{8.669}{7.921}{0.500}{%
\RG{0.000}{0.003}{10.9091}{mi}{70}%
}\hskip.1em\allowbreak{} 必有一个一次因子，即有理根。

\allowbreak{}\hskip.1em \RB{8.669}{7.917}{0.501}{%
\RG{0.000}{-0.001}{10.9091}{mi}{70}%
}\hskip.1em\allowbreak{} 为首一整系数多项式且常数项为\allowbreak{}\hskip.1em \RB{5.846}{5.321}{2.617}{%
\RG{0.000}{-0.001}{10.9091}{mi}{112}%
}\hskip.1em\allowbreak{} ，故有理根定理表明其有理根只能取\allowbreak{}\hskip.1em \RB{33.119}{7.765}{2.617}{%
\RG{0.000}{-0.001}{10.9091}{sy}{6}\RG{8.485}{-0.001}{10.9091}{rm}{49}\RG{13.939}{-0.001}{10.9091}{mi}{59}\RG{18.788}{-0.001}{10.9091}{sy}{6}\RG{27.273}{-0.001}{10.9091}{mi}{112}%
}\hskip.1em\allowbreak{} 。但

\[
\RB{89.461}{8.681}{3.228}{%
\RG{0.000}{-0.001}{10.9091}{mi}{70}\RG{8.530}{-0.001}{10.9091}{rm}{40}\RG{12.763}{-0.001}{10.9091}{rm}{49}\RG{18.217}{-0.001}{10.9091}{rm}{41}\RG{25.472}{-0.001}{10.9091}{rm}{61}\RG{36.981}{-0.001}{10.9091}{rm}{50}\RG{42.469}{-0.001}{10.9091}{mi}{112}\RG{50.382}{-0.001}{10.9091}{rm}{43}\RG{61.281}{-0.001}{10.9091}{rm}{51}\RG{69.777}{-0.001}{10.9091}{mi}{62}\RG{81.292}{-0.001}{10.9091}{rm}{48}\RG{86.747}{-0.001}{10.9091}{mi}{59}%
}
\]
\noindent \RB{144.492}{33.668}{6.583}{%
\RG{45.500}{33.179}{10.9092}{ex}{56}\RG{45.500}{23.361}{10.9092}{ex}{60}\RG{55.197}{19.347}{10.9091}{sy}{0}\RG{63.683}{19.347}{10.9091}{rm}{49}\RG{69.138}{19.347}{10.9091}{mi}{59}\RG{98.854}{19.347}{10.9091}{mi}{110}\RT{109.049}{19.347}{10.9091}{为}\RT{119.958}{19.347}{10.9091}{偶}\RT{130.867}{19.347}{10.9091}{数}\RG{141.777}{19.347}{10.9091}{mi}{59}\RG{0.000}{10.815}{10.9091}{mi}{70}\RG{8.531}{10.815}{10.9091}{rm}{40}\RG{12.773}{10.815}{10.9091}{sy}{0}\RG{21.258}{10.815}{10.9091}{rm}{49}\RG{26.713}{10.815}{10.9091}{rm}{41}\RG{33.967}{10.815}{10.9091}{rm}{61}\RG{45.500}{3.724}{10.9092}{ex}{58}\RG{55.197}{-0.001}{10.9091}{rm}{50}\RG{60.653}{-0.001}{10.9091}{mi}{112}\RG{68.566}{-0.001}{10.9091}{rm}{43}\RG{79.464}{-0.001}{10.9091}{rm}{49}\RG{84.929}{-0.001}{10.9091}{mi}{59}\RG{98.860}{-0.001}{10.9091}{mi}{110}\RT{109.049}{-0.001}{10.9091}{为}\RT{119.958}{-0.001}{10.9091}{奇}\RT{130.867}{-0.001}{10.9091}{数}\RG{141.777}{-0.001}{10.9091}{mi}{59}%
}

以及

\[
\RB{138.749}{10.299}{3.232}{%
\RG{80.755}{4.499}{7.9701}{mi}{110}\RG{85.893}{4.499}{7.9701}{sy}{0}\RG{92.480}{4.499}{7.9701}{rm}{49}\RG{0.000}{-0.005}{10.9091}{mi}{70}\RG{8.531}{-0.005}{10.9091}{rm}{40}\RG{12.773}{-0.005}{10.9091}{mi}{112}\RG{18.262}{-0.005}{10.9091}{rm}{41}\RG{25.517}{-0.005}{10.9091}{rm}{61}\RG{37.026}{-0.005}{10.9091}{rm}{50}\RG{42.481}{-0.005}{10.9091}{rm}{40}\RG{46.747}{-0.005}{10.9091}{mi}{112}\RG{54.659}{-0.005}{10.9091}{rm}{43}\RG{65.558}{-0.005}{10.9091}{rm}{49}\RG{71.012}{-0.005}{10.9091}{rm}{41}\RG{75.267}{-0.005}{10.9091}{mi}{112}\RG{99.636}{-0.005}{10.9091}{rm}{43}\RG{110.545}{-0.005}{10.9091}{mi}{112}\RG{119.054}{-0.005}{10.9091}{mi}{62}\RG{130.580}{-0.005}{10.9091}{rm}{48}\RG{136.034}{-0.005}{10.9091}{mi}{59}%
}
\]
\noindent \RB{181.016}{14.316}{25.935}{%
\RG{45.534}{13.827}{10.9092}{ex}{56}\RG{45.534}{4.008}{10.9092}{ex}{60}\RG{84.997}{3.953}{7.9701}{mi}{110}\RG{90.135}{3.953}{7.9701}{sy}{0}\RG{96.722}{3.953}{7.9701}{rm}{49}\RG{55.231}{-0.005}{10.9091}{mi}{112}\RG{63.145}{-0.005}{10.9091}{sy}{0}\RG{74.054}{-0.005}{10.9091}{rm}{50}\RG{79.508}{-0.005}{10.9091}{mi}{112}\RG{104.484}{-0.005}{10.9091}{mi}{60}\RG{115.999}{-0.005}{10.9091}{rm}{48}\RG{121.454}{-0.005}{10.9091}{mi}{59}\RG{135.374}{-0.005}{10.9091}{mi}{110}\RT{145.574}{-0.005}{10.9091}{为}\RT{156.483}{-0.005}{10.9091}{偶}\RT{167.393}{-0.005}{10.9091}{数}\RG{178.301}{-0.005}{10.9091}{mi}{59}\RG{0.000}{-8.537}{10.9091}{mi}{70}\RG{8.530}{-8.537}{10.9091}{rm}{40}\RG{12.772}{-8.537}{10.9091}{sy}{0}\RG{21.257}{-8.537}{10.9091}{mi}{112}\RG{26.746}{-8.537}{10.9091}{rm}{41}\RG{34.001}{-8.537}{10.9091}{rm}{61}\RG{66.175}{-15.395}{7.9701}{mi}{110}\RG{71.313}{-15.395}{7.9701}{sy}{0}\RG{77.900}{-15.395}{7.9701}{rm}{49}\RG{45.534}{-15.628}{10.9092}{ex}{58}\RG{55.231}{-19.354}{10.9091}{rm}{50}\RG{60.685}{-19.354}{10.9091}{mi}{112}\RG{85.056}{-19.354}{10.9091}{rm}{43}\RG{95.965}{-19.354}{10.9091}{mi}{112}\RG{104.474}{-19.354}{10.9091}{mi}{62}\RG{115.999}{-19.354}{10.9091}{rm}{48}\RG{121.454}{-19.354}{10.9091}{mi}{59}\RG{135.374}{-19.354}{10.9091}{mi}{110}\RT{145.574}{-19.354}{10.9091}{为}\RT{156.483}{-19.354}{10.9091}{奇}\RT{167.393}{-19.354}{10.9091}{数}\RG{178.301}{-19.354}{10.9091}{mi}{58}%
}

因此\allowbreak{}\hskip.1em \RB{8.669}{7.917}{0.501}{%
\RG{0.000}{-0.001}{10.9091}{mi}{70}%
}\hskip.1em\allowbreak{} 不存在有理根，与可约的假设矛盾，故\allowbreak{}\hskip.1em \RB{8.669}{7.917}{0.501}{%
\RG{0.000}{-0.001}{10.9091}{mi}{70}%
}\hskip.1em\allowbreak{} 在\allowbreak{}\hskip.1em \RB{19.985}{8.681}{3.228}{%
\RG{0.000}{-0.001}{10.9091}{bb}{81}\RG{8.485}{-0.001}{10.9091}{rm}{91}\RG{11.515}{-0.001}{10.9091}{mi}{120}\RG{17.750}{-0.001}{10.9091}{rm}{93}%
}\hskip.1em\allowbreak{} 中不可约。\allowbreak{}\hskip.1em \RB{8.366}{7.994}{0.501}{%
\RG{0.000}{-0.001}{10.9091}{am}{3}%
}\hskip.1em\allowbreak{}方法提炼：模素数后的根的阶数。一般化Eisenstein判别法的核心是：\allowbreak{}\hskip.1em \RB{29.832}{9.757}{3.228}{%
\RG{0.000}{-0.001}{10.9091}{mi}{112}\RG{5.489}{3.957}{7.9701}{rm}{50}\RG{13.429}{-0.001}{10.9091}{bb}{45}\RG{19.667}{-0.001}{10.9091}{mi}{97}\RG{25.434}{-1.638}{7.9701}{rm}{48}%
}\hskip.1em\allowbreak{} 保证恰有一个不可约因子具有被\allowbreak{}\hskip.1em \RB{5.846}{5.317}{2.621}{%
\RG{0.000}{-0.005}{10.9091}{mi}{112}%
}\hskip.1em\allowbreak{} 整除的常数项；而\allowbreak{}\hskip.1em \RB{27.923}{9.985}{3.232}{%
\RG{0.000}{-0.005}{10.9091}{mi}{120}\RG{6.235}{3.954}{7.9701}{mi}{107}\RG{14.863}{-0.005}{10.9091}{sy}{106}\RG{21.401}{-0.005}{10.9091}{mi}{102}\RR{21.401}{9.098}{6.516}{0.436}%
}\hskip.1em\allowbreak{} 强迫此因子的次数至少为\allowbreak{}\hskip.1em \RB{6.042}{8.066}{0.625}{%
\RG{0.000}{-0.005}{10.9091}{mi}{107}%
}\hskip.1em\allowbreak{} 。再结合有理根定理便可处理第(2)问。

题目来源：谢启鸿《复旦大学高等代数习题集（2023版）》，第109–110页。

% Original page 8.
\section*{2022S01}
设\allowbreak{}\hskip.1em \RB{99.330}{5.326}{2.732}{%
\RG{0.000}{0.004}{10.9091}{mi}{120}\RG{6.235}{-1.633}{7.9701}{rm}{49}\RG{10.967}{0.004}{10.9091}{mi}{59}\RG{15.800}{0.004}{10.9091}{mi}{58}\RG{20.644}{0.004}{10.9091}{mi}{58}\RG{25.476}{0.004}{10.9091}{mi}{58}\RG{30.320}{0.004}{10.9091}{mi}{59}\RG{35.164}{0.004}{10.9091}{mi}{120}\RG{41.444}{-1.633}{7.9701}{mi}{110}\RG{47.080}{0.004}{10.9091}{mi}{59}\RG{51.913}{0.004}{10.9091}{mi}{121}\RG{57.277}{-1.633}{7.9701}{rm}{49}\RG{62.011}{0.004}{10.9091}{mi}{59}\RG{66.843}{0.004}{10.9091}{mi}{58}\RG{71.687}{0.004}{10.9091}{mi}{58}\RG{76.520}{0.004}{10.9091}{mi}{58}\RG{81.363}{0.004}{10.9091}{mi}{59}\RG{86.207}{0.004}{10.9091}{mi}{121}\RG{91.601}{-1.633}{7.9701}{mi}{109}%
}\hskip.1em\allowbreak{} 均为未定元，

\[
\RB{240.895}{17.650}{51.537}{%
\RG{37.886}{13.635}{7.9701}{mi}{110}\RG{33.485}{10.362}{10.9092}{ex}{89}\RG{101.284}{4.503}{7.9701}{mi}{110}\RG{135.159}{4.503}{7.9701}{mi}{110}\RG{140.297}{4.503}{7.9701}{sy}{0}\RG{146.884}{4.503}{7.9701}{rm}{49}\RG{213.435}{4.503}{7.9701}{mi}{110}\RG{0.000}{-0.002}{10.9091}{mi}{102}\RG{6.515}{-0.002}{10.9091}{rm}{40}\RG{10.758}{-0.002}{10.9091}{mi}{116}\RG{14.697}{-0.002}{10.9091}{rm}{41}\RG{21.952}{-0.002}{10.9091}{rm}{61}\RG{47.425}{-0.002}{10.9091}{rm}{40}\RG{51.667}{-0.002}{10.9091}{mi}{116}\RG{58.031}{-0.002}{10.9091}{sy}{0}\RG{68.941}{-0.002}{10.9091}{mi}{120}\RG{78.557}{-0.002}{10.9091}{rm}{41}\RG{85.811}{-0.002}{10.9091}{rm}{61}\RG{97.345}{-0.002}{10.9091}{mi}{116}\RG{109.345}{-0.002}{10.9091}{sy}{0}\RG{120.254}{-0.002}{10.9091}{mi}{27}\RG{131.220}{-0.002}{10.9091}{mi}{116}\RG{154.040}{-0.002}{10.9091}{rm}{43}\RG{164.949}{-0.002}{10.9091}{sy}{1}\RG{169.782}{-0.002}{10.9091}{sy}{1}\RG{174.626}{-0.002}{10.9091}{sy}{1}\RG{180.101}{-0.002}{10.9091}{rm}{43}\RG{191.000}{-0.002}{10.9091}{rm}{40}\RG{195.253}{-0.002}{10.9091}{sy}{0}\RG{203.738}{-0.002}{10.9091}{rm}{49}\RG{209.192}{-0.002}{10.9091}{rm}{41}\RG{219.071}{-0.002}{10.9091}{mi}{27}\RG{230.941}{-0.002}{10.9091}{mi}{59}\RG{75.175}{-1.638}{7.9701}{mi}{105}\RG{126.487}{-1.638}{7.9701}{rm}{49}\RG{225.305}{-1.638}{7.9701}{mi}{110}\RG{33.603}{-13.089}{7.9701}{mi}{105}\RG{36.486}{-13.089}{7.9701}{rm}{61}\RG{43.070}{-13.089}{7.9701}{rm}{49}\RG{37.092}{-22.687}{7.9701}{mi}{109}\RG{33.868}{-25.960}{10.9092}{ex}{89}\RG{102.164}{-31.819}{7.9701}{mi}{109}\RG{136.926}{-31.819}{7.9701}{mi}{109}\RG{144.417}{-31.819}{7.9701}{sy}{0}\RG{151.004}{-31.819}{7.9701}{rm}{49}\RG{217.555}{-31.819}{7.9701}{mi}{109}\RG{0.921}{-36.324}{10.9091}{mi}{103}\RG{6.515}{-36.324}{10.9091}{rm}{40}\RG{10.758}{-36.324}{10.9091}{mi}{116}\RG{14.697}{-36.324}{10.9091}{rm}{41}\RG{21.952}{-36.324}{10.9091}{rm}{61}\RG{48.190}{-36.324}{10.9091}{rm}{40}\RG{52.432}{-36.324}{10.9091}{mi}{116}\RG{58.796}{-36.324}{10.9091}{sy}{0}\RG{69.706}{-36.324}{10.9091}{mi}{121}\RG{79.437}{-36.324}{10.9091}{rm}{41}\RG{86.691}{-36.324}{10.9091}{rm}{61}\RG{98.224}{-36.324}{10.9091}{mi}{116}\RG{112.577}{-36.324}{10.9091}{sy}{0}\RG{123.486}{-36.324}{10.9091}{mi}{28}\RG{132.987}{-36.324}{10.9091}{mi}{116}\RG{158.160}{-36.324}{10.9091}{rm}{43}\RG{169.069}{-36.324}{10.9091}{sy}{1}\RG{173.902}{-36.324}{10.9091}{sy}{1}\RG{178.746}{-36.324}{10.9091}{sy}{1}\RG{184.221}{-36.324}{10.9091}{rm}{43}\RG{195.120}{-36.324}{10.9091}{rm}{40}\RG{199.373}{-36.324}{10.9091}{sy}{0}\RG{207.858}{-36.324}{10.9091}{rm}{49}\RG{213.312}{-36.324}{10.9091}{rm}{41}\RG{225.544}{-36.324}{10.9091}{mi}{28}\RG{238.301}{-36.324}{10.9091}{mi}{58}\RG{75.055}{-37.960}{7.9701}{mi}{106}\RG{128.255}{-37.960}{7.9701}{rm}{49}\RG{230.313}{-37.960}{7.9701}{mi}{109}\RG{33.485}{-49.411}{7.9701}{mi}{106}\RG{37.369}{-49.411}{7.9701}{rm}{61}\RG{43.953}{-49.411}{7.9701}{rm}{49}%
}
\]
其中\allowbreak{}\hskip.1em \RB{23.024}{5.199}{3.765}{%
\RG{0.000}{-0.003}{10.9091}{mi}{27}\RG{6.233}{-1.639}{7.9701}{mi}{105}\RG{9.614}{-0.003}{10.9091}{mi}{59}\RG{14.447}{-0.003}{10.9091}{mi}{28}\RG{19.232}{-1.639}{7.9701}{mi}{106}%
}\hskip.1em\allowbreak{} 分别是两组未定元的初等对称多项式。设\allowbreak{}\hskip.1em \RB{32.615}{8.679}{3.230}{%
\RG{0.000}{-0.003}{10.9091}{mi}{82}\RG{8.367}{-0.003}{10.9091}{rm}{40}\RG{12.610}{-0.003}{10.9091}{mi}{102}\RG{18.512}{-0.003}{10.9091}{mi}{59}\RG{23.345}{-0.003}{10.9091}{mi}{103}\RG{28.962}{-0.003}{10.9091}{rm}{41}%
}\hskip.1em\allowbreak{} 为多项式\allowbreak{}\hskip.1em \RB{16.406}{8.188}{2.739}{%
\RG{0.000}{-0.003}{10.9091}{mi}{102}\RG{5.902}{-0.003}{10.9091}{mi}{59}\RG{10.735}{-0.003}{10.9091}{mi}{103}%
}\hskip.1em\allowbreak{} 的结式，利用行列式求值的“求根法”证明

\[
\RB{121.973}{17.648}{15.217}{%
\RG{52.151}{13.633}{7.9701}{mi}{110}\RG{67.116}{13.633}{7.9701}{mi}{109}\RG{47.750}{10.360}{10.9092}{ex}{89}\RG{63.891}{10.360}{10.9092}{ex}{89}\RG{0.000}{-0.004}{10.9091}{mi}{82}\RG{8.368}{-0.004}{10.9091}{rm}{40}\RG{12.610}{-0.004}{10.9091}{mi}{102}\RG{18.512}{-0.004}{10.9091}{mi}{59}\RG{23.345}{-0.004}{10.9091}{mi}{103}\RG{28.962}{-0.004}{10.9091}{rm}{41}\RG{36.217}{-0.004}{10.9091}{rm}{61}\RG{78.213}{-0.004}{10.9091}{rm}{40}\RG{82.456}{-0.004}{10.9091}{mi}{120}\RG{94.496}{-0.004}{10.9091}{sy}{0}\RG{105.405}{-0.004}{10.9091}{mi}{121}\RG{115.136}{-0.004}{10.9091}{rm}{41}\RG{119.378}{-0.004}{10.9091}{mi}{58}\RG{88.691}{-1.640}{7.9701}{mi}{105}\RG{110.754}{-1.640}{7.9701}{mi}{106}\RG{47.868}{-13.091}{7.9701}{mi}{105}\RG{50.751}{-13.091}{7.9701}{rm}{61}\RG{57.335}{-13.091}{7.9701}{rm}{49}\RG{63.508}{-13.091}{7.9701}{mi}{106}\RG{67.393}{-13.091}{7.9701}{rm}{61}\RG{73.976}{-13.091}{7.9701}{rm}{49}%
}
\]
解：记\allowbreak{}\hskip.1em \RB{83.293}{8.684}{3.225}{%
\RG{0.000}{0.002}{10.9091}{mi}{70}\RG{8.530}{0.002}{10.9091}{rm}{40}\RG{12.773}{0.002}{10.9091}{mi}{120}\RG{19.962}{0.002}{10.9091}{mi}{59}\RG{24.811}{0.002}{10.9091}{mi}{121}\RG{31.654}{0.002}{10.9091}{rm}{41}\RG{39.029}{0.002}{10.9091}{rm}{61}\RG{50.678}{0.002}{10.9091}{mi}{82}\RG{59.046}{0.002}{10.9091}{rm}{40}\RG{63.289}{0.002}{10.9091}{mi}{102}\RG{69.190}{0.002}{10.9091}{mi}{59}\RG{74.023}{0.002}{10.9091}{mi}{103}\RG{79.641}{0.002}{10.9091}{rm}{41}%
}\hskip.1em\allowbreak{} 。由Sylvester行列式的定义，\allowbreak{}\hskip.1em \RB{8.669}{7.920}{0.500}{%
\RG{0.000}{0.002}{10.9091}{mi}{70}%
}\hskip.1em\allowbreak{} 是多元多项式，其前\allowbreak{}\hskip.1em \RB{9.751}{5.324}{0.618}{%
\RG{0.000}{0.002}{10.9091}{mi}{109}%
}\hskip.1em\allowbreak{} 行来自\allowbreak{}\hskip.1em \RB{6.522}{8.193}{2.734}{%
\RG{0.000}{0.002}{10.9091}{mi}{102}%
}\hskip.1em\allowbreak{} 的系数，后\allowbreak{}\hskip.1em \RB{6.729}{5.321}{0.621}{%
\RG{0.000}{-0.001}{10.9091}{mi}{110}%
}\hskip.1em\allowbreak{} 行来自\allowbreak{}\hskip.1em \RB{5.671}{5.321}{2.737}{%
\RG{0.000}{-0.001}{10.9091}{mi}{103}%
}\hskip.1em\allowbreak{} 的系数。固定\allowbreak{}\hskip.1em \RB{3.696}{7.710}{0.621}{%
\RG{0.000}{-0.001}{10.9091}{mi}{105}%
}\hskip.1em\allowbreak{} ，\allowbreak{}\hskip.1em \RB{6.522}{8.190}{2.737}{%
\RG{0.000}{-0.001}{10.9091}{mi}{102}%
}\hskip.1em\allowbreak{} 的每个系数关于\allowbreak{}\hskip.1em \RB{9.262}{5.321}{2.218}{%
\RG{0.000}{-0.001}{10.9091}{mi}{120}\RG{6.235}{-1.638}{7.9701}{mi}{105}%
}\hskip.1em\allowbreak{} 至多一次，而\allowbreak{}\hskip.1em \RB{5.671}{5.321}{2.737}{%
\RG{0.000}{-0.001}{10.9091}{mi}{103}%
}\hskip.1em\allowbreak{} 与\allowbreak{}\hskip.1em \RB{9.262}{5.321}{2.218}{%
\RG{0.000}{-0.001}{10.9091}{mi}{120}\RG{6.235}{-1.638}{7.9701}{mi}{105}%
}\hskip.1em\allowbreak{} 无关。故利用行列式对每行的线性性，得到

\[
\RB{137.013}{9.295}{4.519}{%
\RT{0.000}{0.002}{10.9091}{d}\RT{6.066}{0.002}{10.9091}{e}\RT{10.909}{0.002}{10.9091}{g}\RG{26.609}{0.002}{10.9091}{mi}{70}\RG{38.169}{0.002}{10.9091}{sy}{20}\RG{49.684}{0.002}{10.9091}{mi}{109}\RG{81.082}{0.002}{10.9091}{rm}{40}\RG{85.314}{0.002}{10.9091}{rm}{49}\RG{93.809}{0.002}{10.9091}{sy}{20}\RG{105.324}{0.002}{10.9091}{mi}{105}\RG{112.113}{0.002}{10.9091}{sy}{20}\RG{123.628}{0.002}{10.9091}{mi}{110}\RG{130.176}{0.002}{10.9091}{rm}{41}\RG{134.418}{0.002}{10.9091}{mi}{58}\RG{16.364}{-2.744}{7.9701}{mi}{120}\RG{21.131}{-3.959}{5.9776}{mi}{105}%
}
\]
任取\allowbreak{}\hskip.1em \RB{13.416}{7.713}{2.734}{%
\RG{0.000}{0.002}{10.9091}{mi}{105}\RG{3.753}{0.002}{10.9091}{mi}{59}\RG{8.586}{0.002}{10.9091}{mi}{106}%
}\hskip.1em\allowbreak{} 。当\allowbreak{}\hskip.1em \RB{33.302}{5.324}{3.760}{%
\RG{0.000}{0.002}{10.9091}{mi}{120}\RG{6.234}{-1.634}{7.9701}{mi}{105}\RG{12.646}{0.002}{10.9091}{rm}{61}\RG{24.161}{0.002}{10.9091}{mi}{121}\RG{29.510}{-1.634}{7.9701}{mi}{106}%
}\hskip.1em\allowbreak{} 时，\allowbreak{}\hskip.1em \RB{6.522}{8.193}{2.734}{%
\RG{0.000}{0.002}{10.9091}{mi}{102}%
}\hskip.1em\allowbreak{} 与\allowbreak{}\hskip.1em \RB{5.671}{5.324}{2.734}{%
\RG{0.000}{0.002}{10.9091}{mi}{103}%
}\hskip.1em\allowbreak{} 有公共根\allowbreak{}\hskip.1em \RB{9.262}{5.324}{2.214}{%
\RG{0.000}{0.002}{10.9091}{mi}{120}\RG{6.235}{-1.634}{7.9701}{mi}{105}%
}\hskip.1em\allowbreak{} 。由于Sylvester行列式对应的线性映射\allowbreak{}\hskip.1em \RB{39.076}{8.684}{3.225}{%
\RG{0.000}{0.002}{10.9091}{rm}{40}\RG{4.242}{0.002}{10.9091}{mi}{117}\RG{10.482}{0.002}{10.9091}{mi}{59}\RG{15.315}{0.002}{10.9091}{mi}{118}\RG{21.015}{0.002}{10.9091}{rm}{41}\RG{28.288}{0.002}{10.9091}{sy}{55}\RG{28.288}{0.002}{10.9091}{sy}{33}%
}\hskip.1em\allowbreak{}\allowbreak{}\hskip.1em \RB{36.605}{8.190}{2.737}{%
\RG{0.000}{-0.001}{10.9091}{mi}{117}\RG{6.240}{-0.001}{10.9091}{mi}{102}\RG{14.773}{-0.001}{10.9091}{rm}{43}\RG{25.272}{-0.001}{10.9091}{mi}{118}\RG{30.934}{-0.001}{10.9091}{mi}{103}%
}\hskip.1em\allowbreak{} 的像中每个多项式均在\allowbreak{}\hskip.1em \RB{27.747}{7.328}{2.218}{%
\RG{0.000}{-0.001}{10.9091}{mi}{116}\RG{6.970}{-0.001}{10.9091}{rm}{61}\RG{18.485}{-0.001}{10.9091}{mi}{120}\RG{24.720}{-1.638}{7.9701}{mi}{105}%
}\hskip.1em\allowbreak{} 处为零，故该映射不可能满射，结式为零。因此作为多元多项式有\allowbreak{}\hskip.1em \RB{57.650}{8.677}{3.767}{%
\RG{0.000}{-0.005}{10.9091}{rm}{40}\RG{4.242}{-0.005}{10.9091}{mi}{120}\RG{10.477}{-1.641}{7.9701}{mi}{105}\RG{15.645}{-0.005}{10.9091}{sy}{0}\RG{25.917}{-0.005}{10.9091}{mi}{121}\RG{31.266}{-1.641}{7.9701}{mi}{106}\RG{35.648}{-0.005}{10.9091}{rm}{41}\RG{42.920}{-0.005}{10.9091}{sy}{106}\RG{48.981}{-0.005}{10.9091}{mi}{70}%
}\hskip.1em\allowbreak{} 。各个\allowbreak{}\hskip.1em \RB{30.816}{5.317}{3.767}{%
\RG{0.000}{-0.005}{10.9091}{mi}{120}\RG{6.235}{-1.641}{7.9701}{mi}{105}\RG{11.403}{-0.005}{10.9091}{sy}{0}\RG{21.675}{-0.005}{10.9091}{mi}{121}\RG{27.024}{-1.641}{7.9701}{mi}{106}%
}\hskip.1em\allowbreak{} 在多项式环\allowbreak{}\hskip.1em \RB{34.192}{8.677}{3.232}{%
\RG{0.000}{-0.005}{10.9091}{mi}{75}\RG{10.046}{-0.005}{10.9091}{rm}{91}\RG{13.076}{-0.005}{10.9091}{mi}{120}\RG{20.265}{-0.005}{10.9091}{mi}{59}\RG{25.114}{-0.005}{10.9091}{mi}{121}\RG{31.957}{-0.005}{10.9091}{rm}{93}%
}\hskip.1em\allowbreak{} 中都是不可约且两两不相伴的多项式，于是

\[
\RB{90.622}{17.646}{15.218}{%
\RG{4.400}{13.631}{7.9701}{mi}{110}\RG{19.364}{13.631}{7.9701}{mi}{109}\RG{0.000}{10.359}{10.9092}{ex}{89}\RG{16.140}{10.359}{10.9092}{ex}{89}\RG{30.462}{-0.005}{10.9091}{rm}{40}\RG{34.704}{-0.005}{10.9091}{mi}{120}\RG{46.745}{-0.005}{10.9091}{sy}{0}\RG{57.654}{-0.005}{10.9091}{mi}{121}\RG{67.386}{-0.005}{10.9091}{rm}{41}\RG{74.658}{-0.005}{10.9091}{sy}{106}\RG{80.719}{-0.005}{10.9091}{mi}{70}\RG{88.028}{-0.005}{10.9091}{mi}{58}\RG{40.939}{-1.641}{7.9701}{mi}{105}\RG{63.003}{-1.641}{7.9701}{mi}{106}\RG{0.117}{-13.092}{7.9701}{mi}{105}\RG{3.001}{-13.092}{7.9701}{rm}{61}\RG{9.584}{-13.092}{7.9701}{rm}{49}\RG{15.757}{-13.092}{7.9701}{mi}{106}\RG{19.641}{-13.092}{7.9701}{rm}{61}\RG{26.225}{-13.092}{7.9701}{rm}{49}%
}
\]
另一方面，此乘积关于每个\allowbreak{}\hskip.1em \RB{9.262}{5.319}{2.220}{%
\RG{0.000}{-0.003}{10.9091}{mi}{120}\RG{6.235}{-1.640}{7.9701}{mi}{105}%
}\hskip.1em\allowbreak{} 的次数都等于\allowbreak{}\hskip.1em \RB{9.751}{5.319}{0.623}{%
\RG{0.000}{-0.003}{10.9091}{mi}{109}%
}\hskip.1em\allowbreak{} ，结合上述次数上界，必存在仅依赖于\allowbreak{}\hskip.1em \RB{6.849}{5.428}{2.707}{%
\RG{0.000}{-0.003}{10.9091}{mi}{121}%
}\hskip.1em\allowbreak{} 的多项式\allowbreak{}\hskip.1em \RB{10.111}{7.954}{0.500}{%
\RG{0.000}{0.003}{10.9091}{mi}{72}%
}\hskip.1em\allowbreak{} ，使得

\[
\RB{151.766}{17.654}{15.211}{%
\RG{81.944}{13.639}{7.9701}{mi}{110}\RG{96.908}{13.639}{7.9701}{mi}{109}\RG{77.544}{10.366}{10.9092}{ex}{89}\RG{93.684}{10.366}{10.9092}{ex}{89}\RG{0.000}{0.002}{10.9091}{mi}{70}\RG{8.530}{0.002}{10.9091}{rm}{40}\RG{12.773}{0.002}{10.9091}{mi}{120}\RG{19.962}{0.002}{10.9091}{mi}{59}\RG{24.810}{0.002}{10.9091}{mi}{121}\RG{31.654}{0.002}{10.9091}{rm}{41}\RG{38.909}{0.002}{10.9091}{rm}{61}\RG{50.442}{0.002}{10.9091}{mi}{72}\RG{60.396}{0.002}{10.9091}{rm}{40}\RG{64.640}{0.002}{10.9091}{mi}{121}\RG{71.483}{0.002}{10.9091}{rm}{41}\RG{108.006}{0.002}{10.9091}{rm}{40}\RG{112.248}{0.002}{10.9091}{mi}{120}\RG{124.289}{0.002}{10.9091}{sy}{0}\RG{135.198}{0.002}{10.9091}{mi}{121}\RG{144.928}{0.002}{10.9091}{rm}{41}\RG{149.171}{0.002}{10.9091}{mi}{58}\RG{118.483}{-1.634}{7.9701}{mi}{105}\RG{140.546}{-1.634}{7.9701}{mi}{106}\RG{77.661}{-13.085}{7.9701}{mi}{105}\RG{80.545}{-13.085}{7.9701}{rm}{61}\RG{87.128}{-13.085}{7.9701}{rm}{49}\RG{93.301}{-13.085}{7.9701}{mi}{106}\RG{97.185}{-13.085}{7.9701}{rm}{61}\RG{103.768}{-13.085}{7.9701}{rm}{49}%
}
\]
最后令\allowbreak{}\hskip.1em \RB{87.437}{7.761}{2.221}{%
\RG{0.000}{-0.005}{10.9091}{mi}{120}\RG{6.235}{-1.641}{7.9701}{rm}{49}\RG{14.453}{-0.005}{10.9091}{rm}{61}\RG{26.423}{-0.005}{10.9091}{sy}{1}\RG{31.256}{-0.005}{10.9091}{sy}{1}\RG{36.100}{-0.005}{10.9091}{sy}{1}\RG{42.636}{-0.005}{10.9091}{rm}{61}\RG{54.606}{-0.005}{10.9091}{mi}{120}\RG{60.842}{-1.641}{7.9701}{mi}{110}\RG{69.963}{-0.005}{10.9091}{rm}{61}\RG{81.919}{-0.005}{10.9091}{rm}{48}%
}\hskip.1em\allowbreak{} 。此时\allowbreak{}\hskip.1em \RB{43.713}{8.677}{3.232}{%
\RG{0.000}{-0.005}{10.9091}{mi}{102}\RG{6.515}{-0.005}{10.9091}{rm}{40}\RG{10.758}{-0.005}{10.9091}{mi}{116}\RG{14.697}{-0.005}{10.9091}{rm}{41}\RG{22.410}{-0.005}{10.9091}{rm}{61}\RG{34.395}{-0.005}{10.9091}{mi}{116}\RG{38.335}{3.954}{7.9701}{mi}{110}%
}\hskip.1em\allowbreak{} ，Sylvester行列式来自\allowbreak{}\hskip.1em \RB{6.522}{8.186}{2.741}{%
\RG{0.000}{-0.005}{10.9091}{mi}{102}%
}\hskip.1em\allowbreak{} 的前\allowbreak{}\hskip.1em \RB{9.751}{5.317}{0.625}{%
\RG{0.000}{-0.005}{10.9091}{mi}{109}%
}\hskip.1em\allowbreak{} 行是前\allowbreak{}\hskip.1em \RB{9.751}{5.317}{0.625}{%
\RG{0.000}{-0.005}{10.9091}{mi}{109}%
}\hskip.1em\allowbreak{} 个标准单位行向量。沿这些行展开，剩下的\allowbreak{}\hskip.1em \RB{6.729}{5.324}{0.618}{%
\RG{0.000}{0.002}{10.9091}{mi}{110}%
}\hskip.1em\allowbreak{} 阶子阵是以\allowbreak{}\hskip.1em \RB{18.935}{8.684}{3.225}{%
\RG{0.000}{0.002}{10.9091}{mi}{103}\RG{5.595}{0.002}{10.9091}{rm}{40}\RG{9.828}{0.002}{10.9091}{rm}{48}\RG{15.282}{0.002}{10.9091}{rm}{41}%
}\hskip.1em\allowbreak{} 为主对角元的三角阵，故

\[
\RB{245.132}{8.681}{3.228}{%
\RG{66.075}{4.504}{7.9701}{mi}{110}\RG{132.901}{4.504}{7.9701}{mi}{110}\RG{175.507}{4.504}{7.9701}{mi}{109}\RG{182.991}{4.504}{7.9701}{mi}{110}\RG{236.901}{4.504}{7.9701}{mi}{110}\RG{0.000}{-0.001}{10.9091}{mi}{70}\RG{8.530}{-0.001}{10.9091}{rm}{40}\RG{12.773}{-0.001}{10.9091}{rm}{48}\RG{19.045}{-0.001}{10.9091}{mi}{59}\RG{23.894}{-0.001}{10.9091}{mi}{121}\RG{30.737}{-0.001}{10.9091}{rm}{41}\RG{37.992}{-0.001}{10.9091}{rm}{61}\RG{49.525}{-0.001}{10.9091}{mi}{82}\RG{57.893}{-0.001}{10.9091}{rm}{40}\RG{62.136}{-0.001}{10.9091}{mi}{116}\RG{71.712}{-0.001}{10.9091}{mi}{59}\RG{76.544}{-0.001}{10.9091}{mi}{103}\RG{82.155}{-0.001}{10.9091}{rm}{40}\RG{86.397}{-0.001}{10.9091}{mi}{116}\RG{90.337}{-0.001}{10.9091}{rm}{41}\RG{94.569}{-0.001}{10.9091}{rm}{41}\RG{101.824}{-0.001}{10.9091}{rm}{61}\RG{113.367}{-0.001}{10.9091}{mi}{103}\RG{118.962}{-0.001}{10.9091}{rm}{40}\RG{123.195}{-0.001}{10.9091}{rm}{48}\RG{128.649}{-0.001}{10.9091}{rm}{41}\RG{141.568}{-0.001}{10.9091}{rm}{61}\RG{153.066}{-0.001}{10.9091}{rm}{40}\RG{157.325}{-0.001}{10.9091}{sy}{0}\RG{165.810}{-0.001}{10.9091}{rm}{49}\RG{171.265}{-0.001}{10.9091}{rm}{41}\RG{188.635}{-0.001}{10.9091}{rm}{40}\RG{192.877}{-0.001}{10.9091}{mi}{121}\RG{204.776}{-0.001}{10.9091}{sy}{1}\RG{209.608}{-0.001}{10.9091}{sy}{1}\RG{214.452}{-0.001}{10.9091}{sy}{1}\RG{219.322}{-0.001}{10.9091}{mi}{121}\RG{232.659}{-0.001}{10.9091}{rm}{41}\RG{242.538}{-0.001}{10.9091}{mi}{58}\RG{198.226}{-1.637}{7.9701}{rm}{49}\RG{224.670}{-1.637}{7.9701}{mi}{109}%
}
\]
将\allowbreak{}\hskip.1em \RB{29.202}{7.643}{0.623}{%
\RG{0.000}{-0.003}{10.9091}{mi}{120}\RG{10.818}{-0.003}{10.9091}{rm}{61}\RG{22.931}{-0.003}{10.9091}{rm}{48}%
}\hskip.1em\allowbreak{} 代入上述因式分解式，又有\allowbreak{}\hskip.1em \RB{165.201}{8.679}{3.230}{%
\RG{0.000}{-0.003}{10.9091}{mi}{70}\RG{8.530}{-0.003}{10.9091}{rm}{40}\RG{12.773}{-0.003}{10.9091}{rm}{48}\RG{19.045}{-0.003}{10.9091}{mi}{59}\RG{23.894}{-0.003}{10.9091}{mi}{121}\RG{30.737}{-0.003}{10.9091}{rm}{41}\RG{38.592}{-0.003}{10.9091}{rm}{61}\RG{50.721}{-0.003}{10.9091}{mi}{72}\RG{60.677}{-0.003}{10.9091}{rm}{40}\RG{64.919}{-0.003}{10.9091}{mi}{121}\RG{71.763}{-0.003}{10.9091}{rm}{41}\RG{75.995}{-0.003}{10.9091}{rm}{40}\RG{80.247}{-0.003}{10.9091}{sy}{0}\RG{88.732}{-0.003}{10.9091}{rm}{49}\RG{94.186}{-0.003}{10.9091}{rm}{41}\RG{98.429}{3.956}{7.9701}{mi}{109}\RG{105.913}{3.956}{7.9701}{mi}{110}\RG{111.556}{-0.003}{10.9091}{rm}{40}\RG{115.799}{-0.003}{10.9091}{mi}{121}\RG{121.147}{-1.640}{7.9701}{rm}{49}\RG{127.698}{-0.003}{10.9091}{sy}{1}\RG{132.531}{-0.003}{10.9091}{sy}{1}\RG{137.374}{-0.003}{10.9091}{sy}{1}\RG{142.243}{-0.003}{10.9091}{mi}{121}\RG{147.592}{-1.640}{7.9701}{mi}{109}\RG{155.580}{-0.003}{10.9091}{rm}{41}\RG{159.823}{3.956}{7.9701}{mi}{110}%
}\hskip.1em\allowbreak{} 。由于\allowbreak{}\hskip.1em \RB{62.971}{8.679}{3.230}{%
\RG{0.000}{-0.003}{10.9091}{mi}{75}\RG{10.046}{-0.003}{10.9091}{rm}{91}\RG{13.076}{-0.003}{10.9091}{mi}{121}\RG{18.425}{-1.640}{7.9701}{rm}{49}\RG{23.157}{-0.003}{10.9091}{mi}{59}\RG{27.990}{-0.003}{10.9091}{mi}{58}\RG{32.834}{-0.003}{10.9091}{mi}{58}\RG{37.666}{-0.003}{10.9091}{mi}{58}\RG{42.510}{-0.003}{10.9091}{mi}{59}\RG{47.354}{-0.003}{10.9091}{mi}{121}\RG{52.748}{-1.640}{7.9701}{mi}{109}\RG{60.736}{-0.003}{10.9091}{rm}{93}%
}\hskip.1em\allowbreak{} 是整环，消去非零多项式即得\allowbreak{}\hskip.1em \RB{29.554}{7.954}{0.500}{%
\RG{0.000}{0.003}{10.9091}{mi}{72}\RG{12.985}{0.003}{10.9091}{rm}{61}\RG{24.483}{0.003}{10.9091}{rm}{49}%
}\hskip.1em\allowbreak{} ，从而

\[
\RB{121.973}{17.648}{15.217}{%
\RG{52.151}{13.633}{7.9701}{mi}{110}\RG{67.116}{13.633}{7.9701}{mi}{109}\RG{47.750}{10.360}{10.9092}{ex}{89}\RG{63.890}{10.360}{10.9092}{ex}{89}\RG{0.000}{-0.004}{10.9091}{mi}{82}\RG{8.368}{-0.004}{10.9091}{rm}{40}\RG{12.610}{-0.004}{10.9091}{mi}{102}\RG{18.512}{-0.004}{10.9091}{mi}{59}\RG{23.345}{-0.004}{10.9091}{mi}{103}\RG{28.962}{-0.004}{10.9091}{rm}{41}\RG{36.217}{-0.004}{10.9091}{rm}{61}\RG{78.213}{-0.004}{10.9091}{rm}{40}\RG{82.456}{-0.004}{10.9091}{mi}{120}\RG{94.496}{-0.004}{10.9091}{sy}{0}\RG{105.405}{-0.004}{10.9091}{mi}{121}\RG{115.136}{-0.004}{10.9091}{rm}{41}\RG{119.378}{-0.004}{10.9091}{mi}{58}\RG{88.691}{-1.640}{7.9701}{mi}{105}\RG{110.754}{-1.640}{7.9701}{mi}{106}\RG{47.868}{-13.091}{7.9701}{mi}{105}\RG{50.751}{-13.091}{7.9701}{rm}{61}\RG{57.335}{-13.091}{7.9701}{rm}{49}\RG{63.509}{-13.091}{7.9701}{mi}{106}\RG{67.393}{-13.091}{7.9701}{rm}{61}\RG{73.976}{-13.091}{7.9701}{rm}{49}%
}
\]
\[
\RB{8.366}{7.998}{0.500}{%
\RG{0.000}{0.003}{10.9091}{am}{3}%
}
\]
方法提炼：求根法的完整闭环。利用公共根识别所有线性因子，再比较Sylvester行列式的次数排除额外因子，最后以\allowbreak{}\hskip.1em \RB{42.803}{8.680}{3.229}{%
\RG{0.000}{-0.002}{10.9091}{mi}{102}\RG{6.516}{-0.002}{10.9091}{rm}{40}\RG{10.758}{-0.002}{10.9091}{mi}{116}\RG{14.697}{-0.002}{10.9091}{rm}{41}\RG{21.952}{-0.002}{10.9091}{rm}{61}\RG{33.485}{-0.002}{10.9091}{mi}{116}\RG{37.425}{3.957}{7.9701}{mi}{110}%
}\hskip.1em\allowbreak{} 的特例确定常数。此处没有预先使用待证的结式根式公式。

题目来源：谢启鸿《复旦大学高等代数习题集（2023版）》，第115页。


\end{document}
